% L’Afrique : décolonisation et néocolonialisme — Histoire Tle Tech — Cours signé OnBuch+
% Programme officiel MINESEC. Compiler avec : tectonic N04-afrique-decolonisation-neocolonialisme.tex
\newcommand{\DOCMATIERE}{Histoire}
\newcommand{\DOCNIVEAU}{Tle Tech}
\newcommand{\DOCMODULE}{Histoire – classe de Terminale technique}
\newcommand{\DOCLECON}{N04}
\newcommand{\DOCTITRE}{L’Afrique : décolonisation et néocolonialisme}
% OnBuch+ — préambule « cours signés » ENRICHI (compilé avec tectonic / XeLaTeX)
% Système visuel « Pop » d'OnBuch+ : fond crème (jamais blanc), bordures encre
% épaisses, ombres « dures » décalées, titres Archivo Black en capitales.
% Version MATHÉMATIQUES : graphes (pgfplots/tikz), boîtes pédagogiques
% (définition, propriété, méthode, attention, le savais-tu, exercice résolu,
% exercice, TP, bilan), page de garde et signature OnBuch+ ancrée en bas.
%
% Macros attendues dans main.tex AVANT \input{preamble} :
%   \DOCMATIERE  \DOCNIVEAU  \DOCMODULE  \DOCLECON (numéro)  \DOCTITRE
\documentclass[11pt]{article}

\usepackage{fontspec}
\usepackage[french]{babel}
\usepackage[a4paper,margin=2cm,top=3.4cm,bottom=2.8cm,headheight=1.4cm,footskip=1.3cm]{geometry}
\usepackage{amsmath,amssymb,amsfonts}
\usepackage{mathtools} % \xmapsto, \coloneqq... souvent utilisés par les modèles
\usepackage{cancel} % simplifications barrées (bilans de forces)
% \abs{z} (module, valeur absolue) et \norm{u} (norme), utilisés par les modèles
\providecommand{\abs}{}\let\abs\relax\DeclarePairedDelimiter{\abs}{\lvert}{\rvert}
\providecommand{\norm}{}\let\norm\relax\DeclarePairedDelimiter{\norm}{\lVert}{\rVert}
\usepackage[table]{xcolor} % [table] : fournit aussi \rowcolors (alternance de lignes)
\usepackage{pagecolor}
\usepackage{tikz}
\usetikzlibrary{arrows.meta,positioning,calc,shapes.geometric,shapes.misc,shapes.arrows,decorations.pathmorphing,decorations.pathreplacing,decorations.markings,patterns,shadows,mindmap,trees,fit,backgrounds,matrix,intersections,angles,quotes}
\usepackage{pgfplots}
\usepackage[version=4]{mhchem} % équations nucléaires \ce{^{235}_{92}U}
\usepackage[european,siunitx]{circuitikz} % schémas électriques
\pgfplotsset{compat=1.18}
\usepgfplotslibrary{fillbetween,groupplots} % aires entre courbes (calcul intégral)
\usepackage{enumitem}
\usepackage{fancyhdr}
% [most] charge toutes les bibliothèques tcolorbox (listings, minted, raster,
% documentation...) et gonfle inutilement la mémoire TeX sur un document long
% (~300 boîtes) : on ne charge que ce qui est réellement utilisé.
\usepackage[skins,breakable]{tcolorbox}
\usepackage{graphicx}
\usepackage{colortbl}
\usepackage{booktabs}
\usepackage{tabularx}
\usepackage{multirow}
\usepackage{diagbox} % case d'en-tête diagonale des tableaux à double entrée
% Colonnes extensibles centrée / gauche / droite (C, L, R), souvent utilisées par les modèles
\newcolumntype{C}{>{\centering\arraybackslash}X}
\newcolumntype{L}{>{\raggedright\arraybackslash}X}
\newcolumntype{R}{>{\raggedleft\arraybackslash}X}
\usepackage{array}
\usepackage{multicol}
\usepackage{siunitx}
% parse-numbers=false : accepte aussi \SI{2,5 \times 10^{-3}}{...}
\sisetup{locale=FR,output-decimal-marker={,},group-minimum-digits=4,parse-numbers=false}
\DeclareSIUnit\ohm{\ensuremath{\symup{\Omega}}} % U+2126 absent de la police
\DeclareSIUnit\division{div} % réglages d'oscilloscope (ms/div, V/div)
\DeclareSIUnit\tr{tr} % tours (vitesse de rotation, ex. \SI{1500}{\tr\per\minute})
\DeclareSIUnit\molar{M} % concentration molaire (ex. \SI{3}{\molar}, \SI{1}{\micro\molar})
\DeclareSIUnit\an{an} % année (le modèle confond parfois avec \year, un compteur TeX) — ex. \SI{5}{\kilo\meter\per\an}
\DeclareSIUnit\FCFA{FCFA} % franc CFA (ex. \SI{500}{\FCFA\per\kilo\gram})
\DeclareSIUnit\degreeBrix{\textdegree Brix} % teneur en sucre d'un sirop/jus (réfractométrie)
\DeclareSIUnit\month{mois} % le modèle utilise parfois \month, un compteur TeX, comme unité de durée
\DeclareSIUnit\microliter{\micro\liter} % le modèle écrit parfois \microliter en un seul token
\DeclareSIUnit\million{millions} % dénombrement (ex. \SI{15}{\million\per\mL} spermatozoïdes)
\usepackage{qrcode}
% \degree : commande fréquemment utilisée par les modèles pour un angle ou une
% température en dehors de \SI{}{} (non fournie par les paquets chargés).
\providecommand{\degree}{^{\circ}}
% \qed (fin de démonstration), utilisé par les modèles sans amsthm
\providecommand{\qed}{\hfill$\square$}
% \barre : double barre des valeurs interdites dans un tableau de variations (tkz-tab)
\providecommand{\barre}{\ensuremath{\|}}
% environnement proof (amsthm non chargé) utilisé par les modèles
\ifdefined\proof\else\newenvironment{proof}[1][Démonstration]{\par\noindent\textit{#1.}\ }{\qed\par}\fi
\usepackage{lastpage}
\usepackage{hyperref}
\hypersetup{hidelinks}

% ============================================================
% Palette « Pop » OnBuch+ — mêmes hex que la classe P (plus_ui.dart)
% ============================================================
\definecolor{poporange}{HTML}{F59321}
\definecolor{popdark}{HTML}{E07A0C}
\definecolor{popcream}{HTML}{FFF6E8}
\definecolor{popink}{HTML}{1C1714}
\definecolor{poppurple}{HTML}{7A5AE0}
\definecolor{popgreen}{HTML}{2E9E6A}
\definecolor{poppink}{HTML}{E0457B}
\definecolor{popblue}{HTML}{2F7DE1}
\definecolor{popgold}{HTML}{C98A12}
\definecolor{popmuted}{HTML}{8A7F76}
\definecolor{popline}{HTML}{EADFD2}
\definecolor{popgray}{HTML}{8A7F76}
\colorlet{popgrey}{popgray}
% Alias tolérants (le modèle nomme parfois les couleurs différemment)
\colorlet{popwhite}{white}
\colorlet{popblack}{popink}
\colorlet{popred}{poppink}
\colorlet{popyellow}{popgold}
% Teintes claires (fonds de tableaux/encadrés) — le modèle nomme parfois
% les couleurs avec un suffixe L (ex. popblueL) sans les avoir définies.
\colorlet{poporangeL}{poporange!20}
\colorlet{popdarkL}{popdark!20}
\colorlet{poppurpleL}{poppurple!20}
\colorlet{popgreenL}{popgreen!20}
\colorlet{poppinkL}{poppink!20}
\colorlet{popblueL}{popblue!20}
\colorlet{popgoldL}{popgold!20}
\colorlet{popredL}{popred!20}
\colorlet{popgrayL}{popgray!20}
% Teintes claires pour fonds de tableaux / zones
\colorlet{popblueL}{popblue!10!white}
\colorlet{poporangeL}{poporange!14!white}
\colorlet{popgreenL}{popgreen!12!white}
\colorlet{poppurpleL}{poppurple!12!white}
\colorlet{poppinkL}{poppink!12!white}
\colorlet{popgoldL}{popgold!14!white}

\pagecolor{popcream}

% ============================================================
% Typographie
% ============================================================
\defaultfontfeatures{Ligatures=TeX}
\setmainfont{PlusJakartaSans-Medium.ttf}[Path=../../../fonts/,BoldFont=PlusJakartaSans-Bold.ttf]
% Maths en sans-serif (Fira Math) avec chiffres et lettres droites de Plus
% Jakarta, pour que les formules se fondent dans le texte.
\usepackage[math-style=ISO,bold-style=ISO]{unicode-math}
\setmathfont{FiraMath-Regular.otf}[Path=../../../fonts/]
\setmathfont{PlusJakartaSans-Medium.ttf}[Path=../../../fonts/,range={up/num,up/{Latin,latin},\mathup}]
\usepackage{icomma} % virgule décimale sans espace en mode math
% Caractères mathématiques Unicode absents des polices de texte (Plus Jakarta,
% Archivo Black) : rendus par leur équivalent mathématique, en texte comme en
% formule — sinon ils s'affichent en carré vide (ex. « Arithmétique dans ℤ »).
\usepackage{newunicodechar}
\newunicodechar{ℝ}{\ensuremath{\mathbb{R}}}
\newunicodechar{ℤ}{\ensuremath{\mathbb{Z}}}
\newunicodechar{ℂ}{\ensuremath{\mathbb{C}}}
\newunicodechar{ℕ}{\ensuremath{\mathbb{N}}}
\newunicodechar{ℚ}{\ensuremath{\mathbb{Q}}}
\newunicodechar{𝔻}{\ensuremath{\mathbb{D}}}
\newunicodechar{α}{\ensuremath{\alpha}}
\newunicodechar{β}{\ensuremath{\beta}}
\newunicodechar{γ}{\ensuremath{\gamma}}
\newunicodechar{δ}{\ensuremath{\delta}}
\newunicodechar{ε}{\ensuremath{\varepsilon}}
\newunicodechar{θ}{\ensuremath{\theta}}
\newunicodechar{λ}{\ensuremath{\lambda}}
\newunicodechar{μ}{\ensuremath{\mu}}
\newunicodechar{σ}{\ensuremath{\sigma}}
\newunicodechar{τ}{\ensuremath{\tau}}
\newunicodechar{φ}{\ensuremath{\varphi}}
\newunicodechar{ψ}{\ensuremath{\psi}}
\newunicodechar{ω}{\ensuremath{\omega}}
\newunicodechar{Σ}{\ensuremath{\Sigma}}
% majuscules produites par \MakeUppercase dans les titres (α-aminés -> Α)
\newunicodechar{Α}{\ensuremath{\alpha}}
\newunicodechar{Β}{\ensuremath{\beta}}
\newunicodechar{Γ}{\ensuremath{\Gamma}}
\newunicodechar{Δ}{\ensuremath{\Delta}}
\newunicodechar{Ε}{\ensuremath{\varepsilon}}
\newunicodechar{Θ}{\ensuremath{\Theta}}
\newunicodechar{Λ}{\ensuremath{\Lambda}}
\newunicodechar{Μ}{\ensuremath{\mu}}
\newunicodechar{Τ}{\ensuremath{\tau}}
\newunicodechar{Φ}{\ensuremath{\Phi}}
\newunicodechar{Ψ}{\ensuremath{\Psi}}
\newunicodechar{Ω}{\ensuremath{\Omega}}
\newunicodechar{↦}{\ensuremath{\mapsto}}
\newunicodechar{∈}{\ensuremath{\in}}
\newunicodechar{∉}{\ensuremath{\notin}}
\newunicodechar{∧}{\ensuremath{\wedge}}
\newunicodechar{⇔}{\ensuremath{\Leftrightarrow}}
\newunicodechar{⟺}{\ensuremath{\Longleftrightarrow}}
\newunicodechar{⇒}{\ensuremath{\Rightarrow}}
\newunicodechar{⟹}{\ensuremath{\Longrightarrow}}
\newunicodechar{∘}{\ensuremath{\circ}}
\newunicodechar{∪}{\ensuremath{\cup}}
\newunicodechar{∩}{\ensuremath{\cap}}
\newunicodechar{≡}{\ensuremath{\equiv}}
\newunicodechar{⊂}{\ensuremath{\subset}}
\newunicodechar{⊥}{\ensuremath{\perp}}
\newunicodechar{∃}{\ensuremath{\exists}}
\newunicodechar{∀}{\ensuremath{\forall}}
\newunicodechar{∛}{\ensuremath{\sqrt[3]{\,}}}
\newunicodechar{∜}{\ensuremath{\sqrt[4]{\,}}}
\newunicodechar{⃗}{\ensuremath{{}^{\to}}}
\newunicodechar{ⁿ}{\ifmmode^{n}\else\textsuperscript{\ensuremath{n}}\fi}
\newunicodechar{ˣ}{\ifmmode^{x}\else\textsuperscript{\ensuremath{x}}\fi}
\newunicodechar{ᵇ}{\ifmmode^{b}\else\textsuperscript{\ensuremath{b}}\fi}
\newunicodechar{ʳ}{\ifmmode^{r}\else\textsuperscript{\ensuremath{r}}\fi}
\newunicodechar{ᵘ}{\ifmmode^{u}\else\textsuperscript{\ensuremath{u}}\fi}
\newunicodechar{ʸ}{\ifmmode^{y}\else\textsuperscript{\ensuremath{y}}\fi}
\newunicodechar{ᵃ}{\ifmmode^{a}\else\textsuperscript{\ensuremath{a}}\fi}
\newunicodechar{ᵉ}{\ifmmode^{e}\else\textsuperscript{\ensuremath{e}}\fi}
\newunicodechar{ᵏ}{\ifmmode^{k}\else\textsuperscript{\ensuremath{k}}\fi}
\newunicodechar{ᵖ}{\ifmmode^{p}\else\textsuperscript{\ensuremath{p}}\fi}
\newunicodechar{ᴺ}{\ifmmode^{N}\else\textsuperscript{\ensuremath{N}}\fi}
\newunicodechar{ᶜ}{\ifmmode^{c}\else\textsuperscript{\ensuremath{c}}\fi}
\newunicodechar{ʰ}{\ifmmode^{h}\else\textsuperscript{\ensuremath{h}}\fi}
\newunicodechar{ᵠ}{\ifmmode^{q}\else\textsuperscript{\ensuremath{q}}\fi}
\newunicodechar{ₙ}{\ifmmode_{n}\else\textsubscript{\ensuremath{n}}\fi}
\newunicodechar{ₐ}{\ifmmode_{a}\else\textsubscript{\ensuremath{a}}\fi}
\newunicodechar{ᵢ}{\ifmmode_{i}\else\textsubscript{\ensuremath{i}}\fi}
\newunicodechar{ᵣ}{\ifmmode_{r}\else\textsubscript{\ensuremath{r}}\fi}
\newunicodechar{ₖ}{\ifmmode_{k}\else\textsubscript{\ensuremath{k}}\fi}
\newunicodechar{ᵦ}{\ifmmode_{\beta}\else\textsubscript{\ensuremath{\beta}}\fi}
\newunicodechar{ₓ}{\ifmmode_{x}\else\textsubscript{\ensuremath{x}}\fi}
\newfontfamily\popdisplay{ArchivoBlack-Regular.ttf}[Path=../../../fonts/]
\newfontfamily\poplogo{SpaceGrotesk-Bold.ttf}[Path=../../../fonts/]
\newfontfamily\popsemi{PlusJakartaSans-SemiBold.ttf}[Path=../../../fonts/]
\newfontfamily\popxbold{PlusJakartaSans-ExtraBold.ttf}[Path=../../../fonts/]
\newfontfamily\popxboldls{PlusJakartaSans-ExtraBold.ttf}[Path=../../../fonts/,LetterSpace=3.0]

\setlist[enumerate]{leftmargin=1.7em,itemsep=5pt,topsep=4pt}
\setlist[itemize]{leftmargin=1.7em,itemsep=4pt,topsep=4pt,label={\color{poporange}\textbullet}}
\setlength{\parindent}{0pt}
\setlength{\parskip}{5pt}
\renewcommand{\arraystretch}{1.25}

% pgfplots : style Pop par défaut
\pgfplotsset{
  every axis/.append style={axis line style={popink,line width=1pt},tick style={popink},
    grid style={popline,line width=0.5pt},label style={font=\small},tick label style={font=\footnotesize}},
  popcurve/.style={poporange,line width=1.8pt,no markers,smooth},
}
% Même style disponible dans un \draw TikZ simple (hors pgfplots)
\tikzset{popcurve/.style={poporange,line width=1.8pt}}
\tikzset{popgrid/.style={popline,very thin}}
\colorlet{popcurve}{poporange} % le nom du style est parfois utilisé comme couleur

% ============================================================
% Pill / squiggle
% ============================================================
\newcommand{\poppill}[3]{%
  \tikz[baseline=-0.55ex]\node[draw=popink,line width=1.2pt,rounded corners=2.6pt,%
    fill=#2,inner xsep=6.5pt,inner ysep=3pt]{%
    \popxboldls\fontsize{7}{7}\selectfont\color{#3}\MakeUppercase{#1}};%
}
\newcommand{\popsquiggle}[1]{%
  \tikz[baseline=-0.4ex]{
    \draw[#1,line width=2.6pt,line cap=round]
      plot [smooth] coordinates {(0,0.09) (0.34,0.01) (0.68,0.21) (1.02,0.01) (1.36,0.10)};
  }%
}
% Mot-clé surligné (marqueur orange sous le texte)
\newcommand{\cle}[1]{{\popxbold\color{popdark}#1}}

% ============================================================
% En-tête / pied
% ============================================================
\pagestyle{fancy}
\fancyhf{}
\renewcommand{\headrulewidth}{0pt}
\renewcommand{\footrulewidth}{0pt}
\lhead{\raisebox{-0.15ex}{\poplogo\fontsize{13}{13}\selectfont\color{popink}OnBuch\color{poporange}+}}
\rhead{\poppill{\DOCMATIERE\ \textbullet\ \DOCNIVEAU\ \textbullet\ Leçon \DOCLECON}{white}{popink}}
\lfoot{\popxbold\fontsize{7.6}{7.6}\selectfont\color{popmuted}\MakeUppercase{Cours signé OnBuch+}\ \textbullet\ {\popsemi\DOCTITRE}}
\rfoot{\tikz[baseline=-0.6ex]\node[rounded corners=8.5pt,fill=popink,minimum height=17pt,minimum width=17pt,inner xsep=5pt,inner sep=0]{\popxbold\fontsize{8}{8}\selectfont\color{popcream}\thepage\,/\,\pageref*{LastPage}};}

% ============================================================
% Titres
% ============================================================
\newcommand{\chaptitre}[2]{%
  \begin{tcolorbox}[enhanced,colback=poporange,colframe=popink,boxrule=2.6pt,arc=11pt,
      drop shadow=popink,left=15pt,right=15pt,top=12pt,bottom=11pt,before skip=3pt,after skip=18pt]
    \poppill{#2}{popink}{popcream}\par\vspace{9pt}
    {\raggedright\hyphenpenalty=10000\exhyphenpenalty=10000\popdisplay\fontsize{22}{24}\selectfont\color{popcream}\MakeUppercase{#1}\par}
  \end{tcolorbox}
}
% Section numérotée : pastille orange avec le numéro + titre Archivo
\newcounter{popsec}
\newcommand{\coursec}[1]{%
  \stepcounter{popsec}\setcounter{popsub}{0}%
  \par\vspace{16pt}\noindent
  \begin{minipage}{\linewidth}
  \tikz[baseline=-0.7ex]\node[draw=popink,line width=1.6pt,fill=poporange,rounded corners=4pt,minimum size=22pt,inner sep=0]{\popdisplay\fontsize{12}{12}\selectfont\color{popcream}\thepopsec};%
  \hspace{8pt}{\popdisplay\fontsize{15}{17}\selectfont\color{popink}\MakeUppercase{#1}}\par
  \vspace{4pt}\noindent{\color{poporange}\rule{36pt}{3.6pt}}
  \end{minipage}\par\nopagebreak\vspace{8pt}
}
\newcounter{popsub}
\newcommand{\courssub}[1]{%
  \stepcounter{popsub}%
  \par\vspace{9pt}\noindent{\popxbold\fontsize{11.5}{13}\selectfont\color{popink}{\color{poporange}\thepopsec.\thepopsub}\hspace{6pt}#1}\par\nopagebreak\vspace{4pt}
}

% ============================================================
% Boîtes pédagogiques — même recette : fond blanc, bordure épaisse colorée,
% ombre dure encre, badge plein. Titre optionnel [#1].
% ============================================================
\tcbset{popbase/.style={breakable,enhanced,colback=white,boxrule=2.3pt,arc=9pt,
  shadow={3pt}{-3pt}{0pt}{popink},left=14pt,right=14pt,top=11pt,bottom=11pt,before skip=11pt,after skip=15pt,
  fontupper=\popsemi\fontsize{10.6}{14.5}\selectfont\color{popink},
  fontlower=\popsemi\fontsize{10.6}{14.5}\selectfont\color{popink}}}
% Coche dessinée (le glyphe ✓ n'existe pas dans les polices du document)
\newcommand{\popcheck}[1]{\tikz[baseline=-0.5ex]\draw[#1,line width=1.6pt,line cap=round,line join=round](-0.13,0.0)--(-0.04,-0.09)--(0.13,0.1);}
\providecommand{\coche}{\popcheck{popgreen}}
\providecommand{\resultat}[1]{\par\smallskip\noindent\fcolorbox{popgreen}{popgreenL}{\parbox{\dimexpr\linewidth-2\fboxsep-2\fboxrule\relax}{\textbf{Résultat.} #1}}\par\smallskip}
\newcommand{\popboxhead}[3]{\poppill{#1}{#2}{white}\ifx\relax#3\relax\else\hspace{6pt}{\popxbold\fontsize{10.6}{12}\selectfont\color{popink}#3}\fi\par\vspace{7pt}}

\newtcolorbox{aretenir}[1][]{popbase,colframe=popgreen,before upper={\popboxhead{À retenir}{popgreen}{#1}}}
\newtcolorbox{exemplebox}[1][]{popbase,colframe=poporange,before upper={\popboxhead{Exemple}{poporange}{#1}}}
\newtcolorbox{definition}[1][]{popbase,colframe=popblue,before upper={\popboxhead{Définition}{popblue}{#1}}}
\newtcolorbox{propriete}[1][]{popbase,colframe=poppurple,before upper={\popboxhead{Propriété}{poppurple}{#1}}}
\newtcolorbox{methode}[1][]{popbase,colframe=popgold,before upper={\popboxhead{Méthode}{popgold}{#1}}}
\newtcolorbox{attention}[1][]{popbase,colframe=poppink,before upper={\popboxhead{Attention}{poppink}{#1}}}
\newtcolorbox{experience}[1][]{popbase,colframe=popblue,colback=popblueL,before upper={\popboxhead{Expérience}{popblue}{#1}}}
% « Le savais-tu ? » : carte violette pleine (contexte, culture, Cameroun)
\newtcolorbox{savaistu}[1][]{popbase,colback=poppurple,colframe=popink,
  fontupper=\popsemi\fontsize{10.4}{14.2}\selectfont\color{white},
  before upper={\poppill{Le savais-tu\,?}{white}{poppurple}\ifx\relax#1\relax\else\hspace{6pt}{\popxbold\color{white}#1}\fi\par\vspace{7pt}}}
% Exercice résolu : énoncé en haut, \tcblower puis la solution sur fond crème
\newcounter{popexo}\newcounter{popexr}
\newtcolorbox{exoresolu}[1][]{popbase,colframe=poporange,colbacklower=popcream,
  segmentation style={popink,line width=1.2pt,dashed},
  before upper={\stepcounter{popexr}\popboxhead{Exercice résolu \thepopexr}{poporange}{#1}},
  before lower={\poppill{Solution}{popink}{popcream}\par\vspace{6pt}}}
% Exercice (non résolu en place) — la correction est dans la partie Corrigés
\newtcolorbox{exercice}[1][]{popbase,colframe=popink,
  before upper={\stepcounter{popexo}\popboxhead{Exercice \thepopexo}{popink}{#1}}}
% Corrigé
\newtcolorbox{corrige}[1][]{popbase,colframe=popgreen,colback=popgreenL,
  before upper={\popboxhead{Corrigé}{popgreen}{#1}}}
% Objectifs / prérequis / bilan
\newtcolorbox{objectifs}{popbase,colframe=popink,colback=poporangeL,
  before upper={\popboxhead{Objectifs de la leçon}{popink}{}}}
\newtcolorbox{prerequis}{popbase,colframe=popink,colback=popblueL,
  before upper={\popboxhead{Prérequis}{popblue}{}}}
\newtcolorbox{bilan}{popbase,colframe=popink,colback=white,boxrule=2.8pt,
  before upper={\popboxhead{Fiche bilan}{poporange}{L'essentiel à connaître pour l'examen}}}
% Figure encadrée avec légende
\newtcolorbox{popfigure}[1][]{enhanced,breakable=false,colback=white,colframe=popline,boxrule=1.4pt,arc=8pt,
  left=10pt,right=10pt,top=10pt,bottom=8pt,before skip=10pt,after skip=12pt,halign=center,
  lower separated=false,halign lower=center,
  fontlower=\popsemi\fontsize{9}{11.5}\selectfont\color{popmuted},
  #1}
\newcommand{\legende}[1]{\par\vspace{6pt}{\popsemi\fontsize{9}{11.5}\selectfont\color{popmuted}#1}}

% ============================================================
% Signature OnBuch+
% ============================================================
\newcommand{\poplogoinline}[1]{{\poplogo\fontsize{#1}{#1}\selectfont\color{popink}OnBuch\color{poporange}+}}
\newcommand{\signatureonbuch}{%
  \begin{tcolorbox}[enhanced,colback=popink,colframe=popink,boxrule=2.2pt,arc=10pt,
      drop shadow=poporange,left=14pt,right=14pt,top=10pt,bottom=10pt,before skip=0pt,after skip=0pt]
    \tikz[baseline=-0.7ex]\node[circle,fill=popgreen,draw=popcream,line width=1.3pt,minimum size=22pt,inner sep=0]{\popcheck{white}};%
    \hspace{9pt}\begin{minipage}[c]{0.62\linewidth}
      {\popxbold\fontsize{12}{13}\selectfont\color{popcream}Signé\ \textbullet\ {\poplogo OnBuch\color{poporange}+}}\par\vspace{2pt}
      {\raggedright\popsemi\fontsize{8}{10.5}\selectfont\color{popline}\MakeUppercase{Équipe pédagogique OnBuch+}\\\MakeUppercase{Conforme au programme officiel MINESEC}\par}
    \end{minipage}\hfill
    {\poplogo\fontsize{22}{22}\selectfont\color{popcream}OnBuch\color{poporange}+}
  \end{tcolorbox}%
}

% ============================================================
% Page de garde
% #1 titre · #2 matière · #3 niveau · #4 module · #5 n° leçon · #6 durée ·
% #7 description / objectifs (texte ou itemize)
% ============================================================
\newcommand{\pagegarde}[7]{%
  \thispagestyle{empty}
  \newgeometry{a4paper,margin=1.7cm,top=1.6cm,bottom=1.4cm}%
  \begin{tikzpicture}[remember picture,overlay]
    \fill[poporange!18] ([xshift=-3.2cm,yshift=-3.6cm]current page.north east) circle (4.6cm);
    \fill[poppurple!14] ([xshift=2.4cm,yshift=7.6cm]current page.south west) circle (3.8cm);
  \end{tikzpicture}%
  {\poplogo\fontsize{30}{30}\selectfont\color{popink}OnBuch\color{poporange}+}\hfill
  \raisebox{0.6ex}{\poppill{Cours signé \textbullet\ Premium}{popink}{popcream}}\par
  \vspace{1pt}\popsquiggle{poppurple}\par
  \vspace{8pt}
  \begin{tcolorbox}[enhanced,colback=poporange,colframe=popink,boxrule=2.8pt,arc=13pt,
      drop shadow=popink,left=18pt,right=18pt,top=12pt,bottom=13pt,after skip=10pt]
    \poppill{#2 \textbullet\ #3}{popink}{popcream}\hspace{5pt}\poppill{#4}{white}{popink}\par\vspace{8pt}
    {\popdisplay\fontsize{14}{14}\selectfont\color{popink}LEÇON #5}\par\vspace{4pt}
    {\raggedright\hyphenpenalty=10000\exhyphenpenalty=10000\popdisplay\fontsize{25}{27}\selectfont\color{popcream}\MakeUppercase{#1}\par}
    \vspace{8pt}
    {\popxbold\fontsize{9.5}{9.5}\selectfont\color{popink}\MakeUppercase{Durée conseillée : #6}}
  \end{tcolorbox}
  \begin{tcolorbox}[enhanced,colback=white,colframe=popink,boxrule=2.2pt,arc=10pt,
      drop shadow=popink,left=16pt,right=16pt,top=10pt,bottom=10pt,after skip=8pt]
    \poppill{Dans cette leçon}{poppurple}{white}\par\vspace{5pt}
    \popsemi\fontsize{9.3}{12.6}\selectfont\color{popink}#7
  \end{tcolorbox}
  \noindent\begin{minipage}[t]{0.487\linewidth}
    \begin{tcolorbox}[enhanced,colback=popgreen,colframe=popink,boxrule=2.2pt,arc=10pt,
        drop shadow=popink,left=11pt,right=11pt,top=10pt,bottom=10pt,height=3.1cm,valign=center]
      \tcbox[colback=white,colframe=popink,boxrule=1.3pt,arc=5pt,boxsep=0pt,nobeforeafter,
        left=3pt,right=3pt,top=3pt,bottom=3pt,valign=center]{\qrcode[height=1.9cm]{https://play.google.com/store/apps/details?id=com.luvvix.onbuch&hl=fr}}%
      \hspace{8pt}\begin{minipage}[c]{0.5\linewidth}
        {\popdisplay\fontsize{11}{12.5}\selectfont\color{white}\MakeUppercase{Télécharge OnBuch}}\par\vspace{4pt}
        {\raggedright\popsemi\fontsize{8.4}{11}\selectfont\color{white}Gratuit \textbullet\ Google Play\\Tuteur IA, annales, concours, résultats\par}
      \end{minipage}
    \end{tcolorbox}
  \end{minipage}\hfill
  \begin{minipage}[t]{0.487\linewidth}
    \begin{tcolorbox}[enhanced,colback=poppurple,colframe=popink,boxrule=2.2pt,arc=10pt,
        drop shadow=popink,left=11pt,right=11pt,top=10pt,bottom=10pt,height=3.1cm,valign=center]
      \tikz[baseline=-0.7ex]\node[circle,fill=white,draw=popink,line width=1.3pt,minimum size=1.6cm,inner sep=0]{\popdisplay\fontsize{24}{24}\selectfont\color{poppurple}+};%
      \hspace{8pt}\begin{minipage}[c]{0.6\linewidth}
        {\popdisplay\fontsize{11}{12.5}\selectfont\color{white}\MakeUppercase{Passe à OnBuch+}}\par\vspace{4pt}
        {\raggedright\popsemi\fontsize{8.4}{11}\selectfont\color{white}Tous les cours signés, Tuteur IA illimité, fiches PDF — onglet OnBuch+ de l'app\par}
      \end{minipage}
    \end{tcolorbox}
  \end{minipage}
  \vspace{8pt}
  \popcontenu
  \vfill
  \signatureonbuch
  \restoregeometry
  \newpage
}

% Rangée « Ce cours contient » (page de garde)
\newcommand{\poptile}[3]{% #1 couleur #2 gros texte #3 légende
  \begin{tcolorbox}[enhanced,colback=#1,colframe=popink,boxrule=2pt,arc=9pt,shadow={2.5pt}{-2.5pt}{0pt}{popink},
      left=6pt,right=6pt,top=9pt,bottom=9pt,halign=center,valign=center,height=1.75cm,width=\linewidth,nobeforeafter]
    {\popdisplay\fontsize{12.5}{13}\selectfont\color{white}\MakeUppercase{#2}}\par\vspace{3pt}
    {\centering\popsemi\fontsize{7.8}{9.5}\selectfont\color{white}#3\par}
  \end{tcolorbox}}
\newcommand{\popcontenu}{%
  \par\noindent{\popxboldls\fontsize{7.5}{7.5}\selectfont\color{popmuted}CE COURS SIGNÉ CONTIENT}\par\vspace{7pt}
  \noindent\begin{minipage}[t]{0.235\linewidth}\poptile{popblue}{Cours}{complet, illustré, pas à pas}\end{minipage}\hfill
  \begin{minipage}[t]{0.235\linewidth}\poptile{poporange}{Méthodes}{et exercices résolus}\end{minipage}\hfill
  \begin{minipage}[t]{0.235\linewidth}\poptile{poppink}{Exercices}{type examen + corrigés}\end{minipage}\hfill
  \begin{minipage}[t]{0.235\linewidth}\poptile{popgreen}{Fiche bilan}{l'essentiel à retenir}\end{minipage}\par}

% Sommaire
\newtcolorbox{sommaire}{enhanced,colback=white,colframe=popink,boxrule=2.2pt,arc=10pt,
  drop shadow=popink,left=18pt,right=18pt,top=13pt,bottom=13pt,before skip=6pt,after skip=20pt,
  before upper={\poppill{Sommaire}{poppurple}{white}\par\vspace{9pt}\popsemi\fontsize{10.6}{14.5}\selectfont\color{popink}}}

% Signature de fin de cours — poussée tout en bas de la dernière page
\newcommand{\findecours}{%
  \par\vfill\nopagebreak
  \begin{center}\popsquiggle{poporange}\end{center}\vspace{4pt}
  \signatureonbuch
}
\AtBeginDocument{\def\llbracket{\mathopen{[\mkern-3.5mu[}}\def\rrbracket{\mathclose{]\mkern-3.5mu]}}} % intervalles d'entiers (glyphe ⟦ absent de la police)
% Glyphes absents de Fira Math : redéfinis à partir de glyphes disponibles
\AtBeginDocument{%
  \def\setminus{\mathbin{\backslash}}\let\smallsetminus\setminus
  \def\vdots{\mathord{\vbox{\baselineskip3.2pt\lineskiplimit0pt\kern4pt\hbox{.}\hbox{.}\hbox{.}}}}%
  \def\ddots{\mathinner{\mkern1mu\raise6.5pt\hbox{.}\mkern2mu\raise3.6pt\hbox{.}\mkern2mu\raise0.7pt\hbox{.}\mkern1mu}}%
  \def\triangle{\mathord{\scalebox{0.9}{$\symup{\Delta}$}}}%
  \def\nabla{\mathord{\raisebox{\depth}{\scalebox{1}[-1]{$\symup{\Delta}$}}}}%
  \def\checkmark{\popcheck{popdark}}%
  \def\star{\mathbin{\ast}}\def\bigstar{\mathord{\ast}}%
  \def\diamond{\mathbin{\scalebox{0.6}{$\Diamond$}}}%
  \def\bigcup{\mathop{\mathchoice{\raisebox{-0.3em}{\scalebox{1.7}{$\cup$}}}{\scalebox{1.25}{$\cup$}}{\cup}{\cup}}\displaylimits}%
  \def\bigcap{\mathop{\mathchoice{\raisebox{-0.3em}{\scalebox{1.7}{$\cap$}}}{\scalebox{1.25}{$\cap$}}{\cap}{\cap}}\displaylimits}%
  \def\lesseqqgtr{\mathrel{\lessgtr}}\def\bigoplus{\mathop{\oplus}\displaylimits}%
}
\providecommand{\ssi}{si et seulement si} % abréviation fréquente
\AtBeginDocument{\providecommand{\wideparen}{\overparen}} % arc de cercle

%% FIGFIX-BEGIN v4 — correctifs globaux des figures (docs/onbuch-generation/tools/figfix)
\usepackage{adjustbox}\usepackage{etoolbox}
% toute figure TikZ/pgfplots plus large que la ligne est réduite à la largeur disponible
\BeforeBeginEnvironment{tikzpicture}{\begin{adjustbox}{max width=\linewidth}}
\AfterEndEnvironment{tikzpicture}{\end{adjustbox}}
% légendes pgfplots lisibles par-dessus les courbes
\pgfplotsset{every axis/.append style={legend style={fill=white,fill opacity=0.92,text opacity=1,draw=black!15,inner sep=3pt}}}
% étiquettes d'arêtes (syntaxe edge["..."]) avec fond blanc
\tikzset{every edge quotes/.append style={fill=white,inner sep=1.2pt}}
%% FIGFIX-END
\begin{document}

\pagegarde{L’Afrique : décolonisation et néocolonialisme}{Histoire}{Tle Tech}{Histoire – classe de Terminale technique}{N04}{5 séances (fiche de progression harmonisée)}{Cette leçon étudie les voies diverses de la décolonisation africaine (négociée au Ghana, armée en Algérie et en Angola, référendaire en Guinée) puis analyse les mécanismes du néocolonialisme qui perpétuent la domination économique et politique des anciennes puissances coloniales sur le continent.
\vspace{4pt}
\begin{multicols}{2}\small
\begin{itemize}
\item À la fin de cette leçon, je sais définir les notions de décolonisation, de néocolonialisme, de Françafrique et de coopération
\item À la fin de cette leçon, je sais retracer les étapes de l'émancipation du Ghana et le rôle de Kwame Nkrumah
\item À la fin de cette leçon, je sais expliquer le déroulement et les conséquences de la guerre d'Algérie (1954-1962)
\item À la fin de cette leçon, je sais analyser le « non » de la Guinée Conakry au référendum du 28 septembre 1958 et ses conséquences
\item À la fin de cette leçon, je sais décrire la lutte armée de libération de l'Angola contre le Portugal
\item À la fin de cette leçon, je sais identifier et illustrer, à partir d'exemples camerounais, les manifestations du néocolonialisme en Afrique
\end{itemize}\end{multicols}}

\begin{sommaire}
\begin{multicols}{2}
\begin{enumerate}
\item Introduction : les facteurs de la décolonisation africaine
\item L'émancipation des territoires britanniques : le cas du Ghana
\item L'émancipation des territoires français : le cas de l'Algérie
\item Dossier 2 : L'indépendance de la Guinée Conakry (activité guidée)
\item L'émancipation des autres territoires : le cas de l'Angola
\item L'Afrique sous l'emprise du néocolonialisme
\item Activité d'intégration
\item Méthodes et exercices résolus
\item Exercices
\item Corrigés des exercices
\item Fiche bilan
\end{enumerate}
\end{multicols}
\end{sommaire}

\begin{objectifs}
\begin{itemize}
\item À la fin de cette leçon, je sais définir les notions de décolonisation, de néocolonialisme, de Françafrique et de coopération
\item À la fin de cette leçon, je sais retracer les étapes de l'émancipation du Ghana et le rôle de Kwame Nkrumah
\item À la fin de cette leçon, je sais expliquer le déroulement et les conséquences de la guerre d'Algérie (1954-1962)
\item À la fin de cette leçon, je sais analyser le « non » de la Guinée Conakry au référendum du 28 septembre 1958 et ses conséquences
\item À la fin de cette leçon, je sais décrire la lutte armée de libération de l'Angola contre le Portugal
\item À la fin de cette leçon, je sais identifier et illustrer, à partir d'exemples camerounais, les manifestations du néocolonialisme en Afrique
\item À la fin de cette leçon, je sais construire une carte, un croquis ou un diagramme à partir de données historiques
\item À la fin de cette leçon, je sais rédiger et justifier une démarche, une réponse ou une conclusion argumentée
\item À la fin de cette leçon, je sais appliquer les notions de l'unité à des situations concrètes de la vie courante camerounaise
\end{itemize}
\end{objectifs}

\begin{prerequis}
\begin{itemize}
\item La colonisation de l'Afrique et la mise en valeur coloniale (classe de Première)
\item La conférence de Berlin (1884-1885) et le partage de l'Afrique
\item Les deux Guerres mondiales et l'affaiblissement des puissances coloniales
\item La création de l'ONU (1945) et le principe du droit des peuples à disposer d'eux-mêmes
\item Les notions d'indépendance, de souveraineté et de nationalisme
\end{itemize}
\end{prerequis}

\newpage

\coursec{Introduction : les facteurs de la décolonisation africaine}

\textbf{Situation d'accroche.} En 1960, le Cameroun accède à l'indépendance : drapeau hissé, hymne national, président camerounais au palais de Yaoundé. Mais soixante ans plus tard, sa monnaie, le franc CFA, est toujours arrimée au Trésor français ; plusieurs de ses grandes entreprises et concessions (Bolloré, TotalEnergies, Orange) restent aux mains de groupes étrangers ; et ses grandes décisions économiques se négocient souvent à Paris ou à Washington. Un élève de Terminale peut légitimement se demander : le Cameroun est-il \emph{vraiment} indépendant ? Comment les peuples africains ont-ils conquis leur souveraineté, et pourquoi cette indépendance semble-t-elle inachevée ?

\textbf{Problématique.} Comment et pourquoi l'Afrique s'est-elle libérée de la domination coloniale entre 1945 et 1975, et quelles limites cette émancipation a-t-elle rencontrées ? Pour répondre à cette question, la leçon suit les parcours contrastés du \cle{Ghana}, de l'\cle{Algérie}, de la \cle{Guinée Conakry} et de l'\cle{Angola}, avant d'interroger la réalité du \cle{néocolonialisme}. Cette première section pose le cadre général : la définition de la décolonisation, sa chronologie et ses facteurs.

\courssub{Qu'est-ce que la décolonisation ?}

\textbf{Intuition.} Imagine un adolescent qui vit sous l'autorité totale d'un tuteur étranger : celui-ci choisit son école, gère son argent, décide à sa place. Le jour de sa majorité, il obtient le droit de décider lui-même. La décolonisation, c'est la « majorité » d'un peuple : le passage d'un territoire administré par une puissance étrangère à un État qui gouverne lui-même ses affaires.

\begin{definition}[La décolonisation]
La \cle{décolonisation} est le processus historique par lequel les territoires colonisés accèdent à la \cle{souveraineté nationale}, c'est-à-dire deviennent des États indépendants dotés d'un gouvernement propre, d'un territoire délimité, d'une population et d'une reconnaissance internationale (siège à l'ONU). En Afrique, ce processus s'étend essentiellement de 1951 (indépendance de la Libye) à 1975 (indépendances des colonies portugaises), et s'achève symboliquement en 1990 avec l'indépendance de la Namibie.
\end{definition}

\begin{definition}[Le nationalisme africain]
Le \cle{nationalisme africain} est le courant de pensée et le mouvement politique qui affirme le droit des peuples africains à disposer d'eux-mêmes et à former des nations libres. Il se traduit par la création de partis politiques (Convention People's Party au Ghana, FLN en Algérie, PDG en Guinée, MPLA en Angola), de syndicats, de journaux, et par des revendications allant de l'égalité des droits à l'indépendance totale.
\end{definition}

\begin{attention}[Piège fréquent]
Ne confonds pas \emph{décolonisation} (le processus, vu du côté des colonisateurs et de l'histoire mondiale) et \emph{indépendance} (l'acte juridique précis, daté, par lequel un territoire devient un État souverain). La décolonisation est un long chemin ; l'indépendance est une étape de ce chemin. Autre piège : la décolonisation n'est pas un « don » des colonisateurs, elle est d'abord \emph{conquise} par la lutte des peuples africains.
\end{attention}

\courssub{Chronologie générale : de 1945 à 1990}

\textbf{Observation.} En 1945, à la fin de la Seconde Guerre mondiale, l'Afrique ne compte que quatre États indépendants : l'Égypte (indépendante formellement depuis 1922, mais sous forte tutelle britannique), l'Éthiopie, le Libéria et l'Afrique du Sud. Quarante-cinq ans plus tard, en 1990, le continent compte \cle{plus de 50 États indépendants}. Comment un tel bouleversement a-t-il pu se produire en une seule génération ?

La décolonisation africaine se déroule par vagues :
\begin{itemize}
\item \textbf{Avant 1956} : premières indépendances (Libye en 1951, Soudan, Maroc et Tunisie en 1956). Le Ghana, première colonie d'Afrique noire à s'émanciper, devient indépendant en 1957.
\item \textbf{1956--1960} : la grande vague. L'année \cle{1960}, surnommée « l'année de l'Afrique », voit 17 colonies accéder à l'indépendance, dont le Cameroun français le 1\ier{} janvier 1960, le Nigeria, le Congo, le Sénégal, la Côte d'Ivoire, Madagascar.
\item \textbf{Après 1960} : les indépendances se poursuivent (Algérie en 1962, Kenya en 1963), puis les dernières colonies portugaises en 1975 (Angola, Mozambique) après de longues guerres de libération. Le Zimbabwe en 1980 et la Namibie en 1990 ferment le cycle.
\end{itemize}

\begin{savaistu}[L'année 1960, « année de l'Afrique »]
L'année \cle{1960} est surnommée « l'année de l'Afrique » car \textbf{17 colonies} y accèdent à l'indépendance, dont le Cameroun français le 1\ier{} janvier 1960. En septembre 1960, seize de ces nouveaux États entrent massivement à l'ONU, transformant durablement l'équilibre de l'Assemblée générale. En février 1960, le Premier ministre britannique Harold Macmillan prononce au Parlement sud-africain du Cap son célèbre discours du « vent du changement » : « Le vent du changement souffle à travers ce continent. »
\end{savaistu}

\begin{attention}[Erreur fréquente : les dates du Cameroun]
Ne confonds pas la date d'\textbf{indépendance du Cameroun français} (\cle{1\ier{} janvier 1960}) avec celle de la \textbf{réunification} : le 1\ier{} octobre 1961, le Cameroun britannique du Sud (Southern Cameroons) se réunit à la République du Cameroun pour former la République fédérale du Cameroun, tandis que le Cameroun britannique du Nord rejoint le Nigeria. Au Baccalauréat technique, inverser ces deux dates coûte des points.
\end{attention}

\begin{popfigure}
\begin{tikzpicture}[scale=1.0]
% Contour très simplifié de l'Afrique
\draw[thick, popdark, fill=popcream] 
  (0.5,5.6) -- (1.2,6.3) -- (2.6,6.6) -- (4.2,6.4) -- (5.4,6.0) -- (6.2,5.4) --
  (6.9,4.9) -- (7.4,4.2) -- (7.1,3.4) -- (6.4,2.8) -- (6.0,2.0) -- (5.5,1.2) --
  (4.8,0.5) -- (4.2,0.1) -- (3.6,0.4) -- (3.2,1.1) -- (2.9,1.9) -- (2.6,2.6) --
  (2.2,3.1) -- (1.6,3.4) -- (1.0,3.6) -- (0.4,4.1) -- (0.1,4.8) -- cycle;
% Zones colorées (approximatives)
\fill[popgreenL, opacity=0.85] (1.2,6.3) -- (2.6,6.6) -- (4.2,6.4) -- (5.4,6.0) -- (5.6,5.2) -- (4.0,5.0) -- (2.4,5.2) -- (1.0,5.4) -- cycle; % Nord : avant 1956
\fill[poporangeL, opacity=0.85] (0.4,4.1) -- (1.0,3.6) -- (1.6,3.4) -- (2.2,3.1) -- (2.6,2.6) -- (4.6,2.6) -- (5.0,3.4) -- (5.6,4.0) -- (6.2,4.4) -- (6.9,4.9) -- (6.2,5.4) -- (5.6,5.2) -- (4.0,5.0) -- (2.4,5.2) -- (1.0,5.4) -- (0.1,4.8) -- cycle; % Centre/Ouest : 1956-1960
\fill[popblueL, opacity=0.85] (2.6,2.6) -- (2.9,1.9) -- (3.2,1.1) -- (3.6,0.4) -- (4.2,0.1) -- (4.8,0.5) -- (5.5,1.2) -- (6.0,2.0) -- (6.4,2.8) -- (6.2,2.6) -- (4.6,2.6) -- cycle; % Sud : après 1960
% Quelques points-villes / pays repères
\fill[popdark] (3.0,5.9) circle (0.06); \node[above, font=\scriptsize] at (3.0,5.95) {Maroc 1956};
\fill[popdark] (1.6,4.3) circle (0.06); \node[left, font=\scriptsize] at (1.55,4.3) {Ghana 1957};
\fill[popdark] (2.4,4.6) circle (0.06); \node[above, font=\scriptsize] at (2.4,4.7) {Guinée 1958};
\fill[popdark] (2.9,3.4) circle (0.06); \node[right, font=\scriptsize] at (2.95,3.4) {Cameroun 1960};
\fill[popdark] (4.9,5.7) circle (0.06); \node[right, font=\scriptsize] at (4.95,5.7) {Algérie 1962};
\fill[popdark] (3.4,1.7) circle (0.06); \node[left, font=\scriptsize] at (3.35,1.7) {Angola 1975};
% Flèche du nord
\draw[-{Stealth}, thick, popink] (7.6,5.6) -- (7.6,6.4); \node[font=\scriptsize\bfseries] at (7.6,6.6) {N};
% Légende
\fill[popgreenL] (0.2,-0.5) rectangle (0.5,-0.2); \node[right, font=\scriptsize] at (0.55,-0.35) {Indépendances avant 1956};
\fill[poporangeL] (3.0,-0.5) rectangle (3.3,-0.2); \node[right, font=\scriptsize] at (3.35,-0.35) {1956--1960 (grande vague)};
\fill[popblueL] (5.9,-0.5) rectangle (6.2,-0.2); \node[right, font=\scriptsize] at (6.25,-0.35) {Après 1960};
\end{tikzpicture}
\legende{Croquis schématique de l'Afrique : les grandes vagues d'indépendance. 1. Afrique du Nord (avant 1956 : Libye 1951, Soudan, Maroc, Tunisie 1956). 2. Afrique de l'Ouest et centrale (1956--1960 : Ghana 1957, Guinée 1958, 17 indépendances en 1960 dont le Cameroun). 3. Afrique australe et colonies portugaises (après 1960 : Algérie 1962, Angola 1975, Namibie 1990). Contour approximatif, sans échelle.}
\end{popfigure}

\begin{methode}[Analyser une carte historique]
Pour exploiter une carte historique à l'examen, suis toujours quatre étapes :
\begin{enumerate}
\item \textbf{Lire le titre} : quel sujet, quel espace, quelle période ?
\item \textbf{Lire la légende} : chaque couleur, symbole ou flèche correspond à une catégorie (ici, trois vagues d'indépendance). Cite-la dans ton analyse.
\item \textbf{S'orienter et situer} : repère le nord, l'échelle indicative et les lieux connus (Cameroun, Algérie, Ghana).
\item \textbf{Décrire puis interpréter} : décris d'abord ce que tu vois (répartition des couleurs), puis dégage l'enseignement historique (la décolonisation avance globalement du nord vers le sud, par vagues, entre 1951 et 1990). Mentionne la source si elle est donnée.
\end{enumerate}
\end{methode}

\begin{popfigure}
\begin{tikzpicture}[scale=1.0]
% Axe chronologique
\draw[-{Stealth}, very thick, popdark] (0,0) -- (13.4,0);
% Graduations
\foreach \x/\a in {0/1945, 2.2/1950, 4.4/1955, 6.6/1960, 8.8/1965, 11/1970, 13.2/1975}{
  \draw[thick, popdark] (\x,0.08) -- (\x,-0.08);
  \node[below, font=\scriptsize] at (\x,-0.12) {\a};
}
% Événements au-dessus
\draw[popblue, thick] (1.1,0) -- (1.1,0.7); \node[above, font=\scriptsize, text width=1.9cm, align=center] at (1.1,0.7) {1947 : indépendance de l'Inde};
\draw[popblue, thick] (3.3,0) -- (3.3,1.5); \node[above, font=\scriptsize, text width=2cm, align=center] at (3.3,1.5) {1954 : Diên Biên Phu ; début de la guerre d'Algérie};
\draw[popblue, thick] (4.4,0) -- (4.4,0.7); \node[above, font=\scriptsize, text width=1.9cm, align=center] at (4.4,0.7) {1955 : conférence de Bandung};
\draw[poporange, thick] (5.5,0) -- (5.5,1.5); \node[above, font=\scriptsize, text width=1.8cm, align=center] at (5.5,1.5) {1957 : indépendance du Ghana};
\draw[poporange, thick] (6.2,0) -- (6.2,0.7); \node[above, font=\scriptsize, text width=1.7cm, align=center] at (6.2,0.7) {1958 : le « non » guinéen};
\draw[poporange, thick] (6.6,0) -- (6.6,1.9); \node[above, font=\scriptsize, text width=2.1cm, align=center] at (6.6,1.9) {1960 : « année de l'Afrique » (17 indépendances ; Cameroun, 1\ier{} janvier) ; résolution 1514 de l'ONU};
\draw[poppurple, thick] (7.7,0) -- (7.7,0.7); \node[above, font=\scriptsize, text width=1.9cm, align=center] at (7.7,0.7) {1962 : indépendance de l'Algérie};
\draw[poppurple, thick] (13.2,0) -- (13.2,0.7); \node[above, font=\scriptsize, text width=2cm, align=center] at (12.9,0.7) {1975 : indépendance de l'Angola};
\end{tikzpicture}
\legende{Frise chronologique des grandes étapes de la décolonisation africaine (1945--1975). En bleu : les événements déclencheurs ; en orange : les premières indépendances d'Afrique noire ; en violet : les indépendances conquises par les armes.}
\end{popfigure}

\courssub{Les facteurs internes : l'éveil des peuples africains}

La décolonisation s'explique d'abord par des causes \emph{intérieures} à l'Afrique : les Africains eux-mêmes ont voulu, organisé et obtenu leur libération.

\begin{itemize}
\item \textbf{La montée du nationalisme.} Nourri par l'humiliation coloniale (travail forcé, impôts arbitraires, code de l'indigénat, discriminations), le sentiment national se transforme en volonté politique d'indépendance. Les discours de Kwame Nkrumah au Ghana ou de Sékou Touré en Guinée en incarnent l'esprit.
\item \textbf{L'émergence d'élites évoluées.} Les « évolués » — instituteurs, infirmiers, employés, commerçants, diplômés des universités coloniales ou européennes (Nkrumah, Senghor, Houphouët-Boigny, Sékou Touré, Um Nyobè au Cameroun) — maîtrisent la langue et les armes juridiques du colonisateur et les retournent contre lui au nom de l'égalité et de la liberté.
\item \textbf{Le rôle des partis politiques.} Ils encadrent les masses et structurent la revendication : la CPP de Nkrumah (Ghana, 1949), le RDA (Rassemblement démocratique africain, 1946), le FLN (Algérie, 1954), l'UPC au Cameroun (1948), le MPLA (Angola, 1956).
\item \textbf{Le rôle des syndicats.} Les syndicats de travailleurs (cheminots, dockers, fonctionnaires) organisent grèves et manifestations, comme la grande grève des cheminots de l'AOF en 1947--1948, qui prouve la capacité de mobilisation des Africains.
\item \textbf{Les anciens combattants.} Des centaines de milliers de soldats africains ont combattu pour la France et la Grande-Bretagne pendant les deux guerres mondiales. Revenus au pays, ils refusent de retourner au statut de sujets : ils ont versé leur sang pour la liberté des autres et réclament la leur.
\end{itemize}

\courssub{Les facteurs externes : un contexte mondial favorable}

\textbf{Observation.} Pourquoi les puissances coloniales, si fortes en 1939, cèdent-elles après 1945 ? Parce que le monde a changé.

\begin{itemize}
\item \textbf{L'affaiblissement de l'Europe après 1945.} Ruinées et épuisées par la Seconde Guerre mondiale, la France et la Grande-Bretagne n'ont plus les moyens militaires, financiers ni moraux de maintenir leurs empires. La charte de l'Atlantique (1941) et la charte de l'ONU (1945) proclament le droit des peuples à disposer d'eux-mêmes.
\item \textbf{L'appui des superpuissances.} Les États-Unis et l'URSS, adversaires dans la guerre froide, s'accordent paradoxalement pour condamner le colonialisme : Washington par anticolonialisme libéral et intérêt économique, Moscou par idéologie et pour attirer les jeunes États dans son camp. Chacun soutient, parfois en armes, les mouvements de libération.
\item \textbf{La pression de l'ONU.} L'organisation devient la tribune des peuples colonisés. Le 14 décembre 1960, l'Assemblée générale adopte la \cle{résolution 1514}, la « Déclaration sur l'octroi de l'indépendance aux pays et aux peuples coloniaux », qui proclame que « l'assujettissement des peuples à une domination étrangère constitue une négation des droits fondamentaux de l'homme » et demande la fin rapide du colonialisme.
\end{itemize}

\courssub{Les exemples déclencheurs}

Trois événements extérieurs à l'Afrique ont montré aux Africains que la libération était possible :

\begin{enumerate}
\item \textbf{L'indépendance de l'Inde (1947).} Le plus grand empire colonial, l'empire britannique des Indes, accède à l'indépendance grâce à la lutte non violente de Gandhi et du Congrès. L'exemple indien prouve qu'un peuple colonisé peut vaincre l'Empire.
\item \textbf{La défaite française de Diên Biên Phu (1954).} Au Vietnam, l'armée française est écrasée par le Vietminh. Le choc est immense : une armée coloniale européenne moderne a été battue par un peuple colonisé. La même année, la guerre d'Algérie éclate.
\item \textbf{La conférence de Bandung (1955).} 29 pays d'Asie et d'Afrique se réunissent en Indonésie pour affirmer leur solidarité anticolonialiste et leur refus des blocs. Bandung donne aux nationalistes africains une légitimité internationale et un sentiment de force collective.
\end{enumerate}

\courssub{Les trois voies de la décolonisation}

Toutes les indépendances ne se sont pas obtenues de la même manière. On distingue classiquement trois voies, qui structurent toute la leçon :

\begin{center}
\begin{tabularx}{\linewidth}{|l|X|X|}
\hline
\rowcolor{popblueL} \textbf{Voie} & \textbf{Principe} & \textbf{Exemple étudié} \\
\hline
Négociée & L'indépendance est obtenue par des discussions et des réformes progressives avec la métropole, sans guerre. & Ghana (1957) \\
\hline
Référendaire & Le peuple choisit l'indépendance par un vote (référendum). & Guinée Conakry (1958) \\
\hline
Armée & L'indépendance est conquise par une guerre de libération contre la puissance coloniale. & Algérie (1954--1962), Angola (1961--1975) \\
\hline
\end{tabularx}
\end{center}

\begin{exoresolu}[Identifier les voies de la décolonisation]
Énoncé : classe les situations suivantes selon la voie de décolonisation correspondante et justifie : (a) le Ghana de Nkrumah obtient l'indépendance en 1957 après des élections et des négociations avec Londres ; (b) la Guinée de Sékou Touré vote « non » au référendum du 28 septembre 1958 et proclame l'indépendance le 2 octobre 1958 ; (c) l'Algérie obtient l'indépendance en 1962 après huit ans de guerre.
\tcblower
Solution : (a) \textbf{Voie négociée} : l'indépendance résulte d'un processus électoral et de discussions avec la Grande-Bretagne, sans conflit armé. (b) \textbf{Voie référendaire} : le peuple guinéen choisit l'indépendance par un vote, en refusant la Communauté française proposée par de Gaulle. (c) \textbf{Voie armée} : l'indépendance algérienne est conquise par la lutte du FLN et de l'ALN de 1954 à 1962, avant d'être confirmée par les accords d'Évian et le référendum d'autodétermination.
\end{exoresolu}

\begin{aretenir}[L'essentiel de l'introduction]
\begin{itemize}
\item La décolonisation est le processus par lequel les territoires colonisés accèdent à la souveraineté nationale ; l'Afrique compte plus de 50 États indépendants en 1990.
\item 1960, « année de l'Afrique » : 17 indépendances, dont le Cameroun français le 1\ier{} janvier 1960 (réunification : 1\ier{} octobre 1961).
\item Facteurs internes : nationalisme, élites évoluées, partis, syndicats, anciens combattants. Facteurs externes : Europe affaiblie, USA et URSS, ONU (résolution 1514, 1960).
\item Déclencheurs : Inde (1947), Diên Biên Phu (1954), Bandung (1955).
\item Trois voies : négociée (Ghana), référendaire (Guinée), armée (Algérie, Angola).
\end{itemize}
\end{aretenir}

\begin{exercice}[Pour s'entraîner]
1. Définis la décolonisation et le nationalisme africain. \quad 2. Explique pourquoi l'année 1960 est appelée « année de l'Afrique ». \quad 3. Présente deux facteurs internes et deux facteurs externes de la décolonisation africaine. \quad 4. En quoi la défaite de Diên Biên Phu a-t-elle encouragé les peuples africains ?
\end{exercice}

\begin{corrige}[Exercice]
1. Décolonisation : processus par lequel les territoires colonisés accèdent à la souveraineté nationale. Nationalisme africain : mouvement affirmant le droit des peuples africains à disposer d'eux-mêmes et à former des nations libres. 2. Parce que 17 colonies y accèdent à l'indépendance, dont le Cameroun français le 1\ier{} janvier 1960. 3. Internes : montée du nationalisme et action des partis politiques (ou syndicats, élites, anciens combattants). Externes : affaiblissement de l'Europe après 1945 et pression de l'ONU (résolution 1514) ou appui des superpuissances. 4. Elle a prouvé qu'une armée coloniale européenne pouvait être vaincue par un peuple colonisé déterminé, rendant l'indépendance crédible et envisageable.
\end{corrige}

\coursec{L'émancipation des territoires britanniques : le cas du Ghana}

Dans la section précédente, nous avons vu que la décolonisation africaine résulte d'un faisceau de facteurs : l'affaiblissement des puissances coloniales après 1945, la montée des nationalismes, le contexte de la Guerre froide et la pression de l'ONU. Mais ces facteurs généraux ne disent pas \emph{comment}, concrètement, un peuple colonisé arrache son indépendance. Pour le comprendre, rien ne vaut l'étude d'un cas précis. Le Ghana — alors appelé \cle{Gold Coast} — offre le cas idéal : c'est la \textbf{première colonie d'Afrique noire} à accéder à l'indépendance, en 1957, et son leader Kwame Nkrumah est devenu le symbole de l'émancipation africaine. Sa démarche, pacifique et organisée, a servi de modèle à tout le continent.

\courssub{Une colonie mûre pour l'émancipation : la Gold Coast}

La \cle{Gold Coast} (littéralement « Côte de l'Or ») est une colonie britannique d'Afrique de l'Ouest, correspondant à l'actuel Ghana. Son nom même indique sa richesse : l'or y est exploité depuis des siècles. Au XX\textsuperscript{e} siècle, la colonie est devenue le \textbf{premier producteur mondial de cacao}, culture de rente qui enrichit à la fois la métropole et une partie des planteurs africains.

Trois atouts distinguent la Gold Coast des autres colonies africaines :

\begin{itemize}
\item Une \textbf{prospérité économique relative} : les revenus du cacao et de l'or permettent l'existence d'une classe de commerçants et de planteurs africains aisés, capables de financer l'action politique.
\item Une \textbf{élite éduquée} : les Britanniques y ont développé l'enseignement (écoles missionnaires, Achimota College). Avocats, enseignants, commerçants forment une bourgeoisie éclairée qui maîtrise la langue et le droit du colonisateur — et sait donc le combattre avec ses propres armes.
\item Une \textbf{presse dynamique} : des journaux africains, comme l'\emph{Accra Evening News} fondé par Nkrumah, diffusent les idées nationalistes auprès de la population urbaine.
\end{itemize}

Dès 1947, l'élite traditionnelle crée le premier parti, l'\cle{UGCC} (United Gold Coast Convention), dirigé par des notables modérés comme J. B. Danquah. Mais ces notables veulent une indépendance lente et négociée. C'est là qu'intervient un homme qui va tout changer.

\courssub{Kwame Nkrumah et la naissance du CPP}

Kwame \cle{Nkrumah} (1909-1972) est né dans une famille modeste de la Gold Coast. Après des études chez lui, il part se former aux \textbf{États-Unis} (Lincoln University) puis à \textbf{Londres}. Cette formation à l'étranger est décisive : il y découvre les idées socialistes, le panafricanisme (il participe au congrès panafricain de Manchester en 1945) et les méthodes de mobilisation des masses.

En \textbf{1947}, l'UGCC le rappelle au pays comme secrétaire général. Mais Nkrumah juge les notables trop timides : il veut l'« indépendance maintenant » (\emph{Self-Government Now}) et non demain. En \textbf{1949}, il rompt avec l'UGCC et fonde le \cle{CPP} (\emph{Convention People's Party}), un parti de masse ouvert aux jeunes, aux ouvriers, aux petits commerçants, aux femmes du marché — ceux que l'élite méprisait et que Nkrumah appelle les \emph{verandah boys}, les « garçons de la véranda » qui dorment dehors.

\begin{definition}[Action positive]
L'\cle{action positive} est la stratégie de lutte non violente imaginée par Nkrumah, inspirée de Gandhi en Inde : \textbf{grèves générales}, \textbf{boycott} des produits britanniques, \textbf{désobéissance civile} et manifestations pacifiques. Le principe : paralyser l'économie et l'administration coloniales \emph{sans violence}, afin de rendre le pays ingouvernable tout en gardant l'opinion internationale de son côté.
\end{definition}

En janvier \textbf{1950}, Nkrumah lance l'action positive : grève générale et boycott. Les Britanniques répriment et emprisonnent Nkrumah, condamné à trois ans de prison. Mais l'emprisonnement du leader le transforme en \textbf{martyr} et accroît sa popularité : erreur classique des puissances coloniales.

\courssub{Les étapes institutionnelles : des urnes à l'indépendance}

La Grande-Bretagne, contrairement à la France en Algérie, choisit la \textbf{négociation}. Elle organise des élections en février \textbf{1951} pour une nouvelle Assemblée législative.

\begin{exemplebox}[Une victoire électorale remportée depuis la prison]
Aux élections de février 1951, le CPP remporte \textbf{34 sièges sur 38} à pourvoir au suffrage, soit environ \textbf{89 \%} des sièges, alors que Nkrumah est toujours derrière les barreaux. Lui-même est élu député d'Accra. Devant ce verdict massif des urnes, le gouverneur britannique n'a d'autre choix que de le \textbf{libérer} et de le nommer \emph{Leader of Government Business} (chef des affaires gouvernementales), puis Premier ministre en 1952.
\end{exemplebox}

La marche vers l'indépendance suit ensuite un escalier institutionnel rigoureux :

\begin{popfigure}
\begin{center}
\begin{tikzpicture}[scale=1, transform shape]
% Escalier des étapes
\draw[fill=popblueL, draw=popblue] (0,0) rectangle (3,1);
\draw[fill=popblue!45, draw=popblue] (3,1) rectangle (6,2);
\draw[fill=popblue!65, draw=popblue] (6,2) rectangle (9,3);
\draw[fill=popblue, draw=popdark] (9,3) rectangle (12,4);
\node[align=center] at (1.5,0.5) {\textbf{1951}\\ Victoire du CPP\\ (34/38 sièges)};
\node[align=center] at (4.5,1.5) {\textbf{1952}\\ Nkrumah\\ Premier ministre};
\node[align=center, text=white] at (7.5,2.5) {\textbf{1954}\\ Nouvelle Constitution\\ autonomie interne};
\node[align=center, text=white] at (10.5,3.5) {\textbf{6 mars 1957}\\ \textbf{INDÉPENDANCE}\\ du Ghana};
% Flèches de progression
\draw[-{Stealth[length=3mm]}, very thick, poporange] (2.9,1.15) -- (3.2,1.15);
\draw[-{Stealth[length=3mm]}, very thick, poporange] (5.9,2.15) -- (6.2,2.15);
\draw[-{Stealth[length=3mm]}, very thick, poporange] (8.9,3.15) -- (9.2,3.15);
% Axe du temps
\draw[-{Stealth[length=3mm]}, thick, popdark] (0,-0.4) -- (12.5,-0.4);
\node[below] at (12.3,-0.4) {\footnotesize temps};
\end{tikzpicture}
\end{center}
\legende{Diagramme en escalier des étapes institutionnelles menant la Gold Coast à l'indépendance (1951-1957). Chaque marche correspond à un transfert de compétences de Londres vers Accra.}
\end{popfigure}

La \textbf{Constitution de 1954} accorde l'\cle{autonomie interne} : le gouvernement ghanéen dirige les affaires intérieures, tandis que Londres garde encore la défense et les affaires étrangères. Le CPP remporte à nouveau les élections de 1954 et de 1956, prouvant que sa légitimité n'est pas un hasard. Enfin, le \textbf{6 mars 1957}, l'indépendance est proclamée. La Gold Coast devient le \cle{Ghana}.

\begin{savaistu}[Pourquoi le nom « Ghana » ?]
Le nouveau pays abandonne le nom colonial de « Gold Coast » et choisit « Ghana » en référence à l'\textbf{ancien empire du Ghana} (IV\textsuperscript{e}-XIII\textsuperscript{e} siècle), brillant royaume d'Afrique de l'Ouest — qui ne se trouvait pourtant pas sur son territoire actuel ! Ce choix est un acte politique : il affirme que l'Afrique a connu une \textbf{grande civilisation précoloniale} et que l'indépendance n'est pas une naissance, mais une \emph{renaissance}.
\end{savaistu}

\begin{attention}[Erreur fréquente au Baccalauréat technique]
Le Ghana n'est \textbf{pas} le premier État indépendant d'Afrique : le \textbf{Libéria} (1847) et l'\textbf{Éthiopie} (jamais durablement colonisée) étaient déjà souverains, sans parler de l'Égypte. La formulation correcte est : le Ghana est la \textbf{première colonie d'Afrique subsaharienne (Afrique noire) à être décolonisée}, c'est-à-dire le premier territoire colonial à arracher son indépendance à une puissance européenne au sud du Sahara.
\end{attention}

\courssub{Analyser un document iconographique : la proclamation du 6 mars 1957}

La photographie de Nkrumah proclamant l'indépendance, la nuit du 5 au 6 mars 1957 au polo ground d'Accra, est l'un des documents les plus étudiés au Baccalauréat technique. Voyons comment l'exploiter.

\begin{methode}[Analyser une photographie historique]
\begin{enumerate}
\item \textbf{Décrire} objectivement : qui voit-on ? Où ? Quand ? Quels gestes, quels symboles (drapeaux, foule, micros) ?
\item \textbf{Contextualiser} : replacer l'image dans son époque et dans le processus historique étudié.
\item \textbf{Dégager l'intention} de l'auteur : pourquoi cette photo a-t-elle été prise et diffusée ? Quel message veut-elle transmettre ?
\item \textbf{Évaluer la portée} : quelle a été l'influence de cette image ? Que prouve-t-elle, et quelles sont ses limites (une photo est un témoignage orienté, pas la réalité entière) ?
\end{enumerate}
\end{methode}

\begin{popfigure}
\begin{center}
\begin{tikzpicture}
% Cadre de la photo
\draw[fill=popcream, draw=popdark, thick] (0,0) rectangle (8,5);
% Foule en bas
\foreach \x in {0.4,1.1,1.8,2.5,3.2,3.9,4.6,5.3,6.0,6.7,7.4}{
  \draw[fill=popmuted, draw=popmuted] (\x,0.55) circle (0.18);
}
\draw[fill=popline!40, draw=none] (0,0) rectangle (8,0.55);
% Podium
\draw[fill=popblueL, draw=popblue] (3.2,0.9) rectangle (4.8,2.0);
% Nkrumah
\draw[fill=popgold, draw=popdark] (4,2.6) circle (0.28);
\draw[draw=popdark, thick] (4,2.32) -- (4,1.95);
\draw[draw=popdark, thick] (4,2.2) -- (3.5,2.55);
\draw[draw=popdark, thick] (4,2.2) -- (4.55,2.6);
% Drapeau
\draw[draw=popdark] (5.6,0.9) -- (5.6,4.3);
\draw[fill=poporange, draw=popdark] (5.6,3.9) rectangle (7.2,4.3);
\draw[fill=popcream, draw=popdark] (5.6,3.5) rectangle (7.2,3.9);
\draw[fill=popgreen, draw=popdark] (5.6,3.1) rectangle (7.2,3.5);
% Micro
\draw[draw=popdark, thick] (3.6,2.0) -- (3.75,2.45);
\draw[fill=popdark] (3.75,2.48) circle (0.06);
% Légendes fléchées
\draw[-{Stealth}, poppurple, thick] (4,2.6) -- (4.1,4.7) node[right, align=left, text=popdark]{\footnotesize \textbf{1.} Nkrumah, bras levé,\\ \footnotesize au centre de la scène};
\draw[-{Stealth}, poppurple, thick] (6.4,3.5) -- (7.6,4.7) node[right, align=left, text=popdark]{\footnotesize \textbf{2.} Nouveau drapeau\\ \footnotesize rouge-or-vert};
\draw[-{Stealth}, poppurple, thick] (2,0.55) -- (1.2,-0.5) node[below, align=center, text=popdark]{\footnotesize \textbf{3.} Foule immense :\\ \footnotesize légitimité populaire};
\draw[-{Stealth}, poppurple, thick] (3.75,2.45) -- (1.6,3.4) node[left, align=right, text=popdark]{\footnotesize \textbf{4.} Micro : discours\\ \footnotesize retransmis à la radio};
\end{tikzpicture}
\end{center}
\legende{Schéma d'analyse de la photographie de la proclamation de l'indépendance du Ghana (Accra, 6 mars 1957). La mise en scène (leader au centre, drapeau neuf, foule, médias) construit l'image d'une indépendance à la fois populaire et moderne.}
\end{popfigure}

\begin{exoresolu}[Analyse de la photographie du 6 mars 1957]
\textbf{Énoncé.} À partir du schéma d'analyse ci-dessus, présentez cette photographie (nature, auteur, contexte), décrivez-la, puis dégagez son intérêt historique.
\tcblower
\textbf{Solution détaillée.} \emph{Nature et contexte} : il s'agit d'une photographie de presse prise à Accra, capitale de la Gold Coast, dans la nuit du 5 au 6 mars 1957, lors de la cérémonie de proclamation de l'indépendance du Ghana, première décolonisation d'Afrique noire. \emph{Description} : au centre, Kwame Nkrumah, Premier ministre, lève le bras devant un micro ; autour de lui, les membres du CPP ; un nouveau drapeau (rouge, or, vert frappé de l'étoile noire) remplace l'\emph{Union Jack} britannique ; en contrebas, une foule immense acclame. \emph{Intention} : le photographe et le pouvoir ghanéen veulent montrer une indépendance \textbf{légitime} (soutenue par le peuple), \textbf{pacifique} (cérémonie, pas guerre) et \textbf{moderne} (radio, symboles nationaux). \emph{Portée} : cette image, diffusée dans le monde entier, fait du Ghana le \textbf{phare de l'émancipation africaine} et de Nkrumah le héros du panafricanisme. \emph{Limite} : photo officielle, elle masque les oppositions internes (l'opposition du Nord, l'Ashanti) qui marqueront la suite.
\end{exoresolu}

\courssub{Le Ghana, phare du panafricanisme}

L'indépendance du Ghana a un effet d'entraînement immédiat : en 1960, dix-sept États africains deviennent indépendants (« année de l'Afrique »). Nkrumah veut aller plus loin : faire du Ghana le moteur de l'unité du continent.

\begin{definition}[Panafricanisme]
Le \cle{panafricanisme} est un mouvement d'idées et d'action politique qui affirme la \textbf{solidarité de tous les Africains} et des descendants d'Africains (diaspora), et qui vise l'\textbf{unité politique et économique du continent africain} pour achever sa libération et peser dans les affaires du monde. Nkrumah le résume : « L'indépendance du Ghana n'a pas de sens si elle n'est pas liée à la libération totale de l'Afrique. »
\end{definition}

Deux réalisations concrètes illustrent ce panafricanisme :

\begin{itemize}
\item La \textbf{conférence d'Accra (avril 1958)} : Nkrumah y réunit les huit États africains alors indépendants (Ghana, Égypte, Éthiopie, Libéria, Libye, Maroc, Soudan, Tunisie). C'est la première conférence des États africains indépendants \emph{sur le sol africain} ; elle affirme le soutien aux peuples encore colonisés.
\item La \textbf{création de l'OUA} : Nkrumah milite pour une union forte des États africains (il rêve même d'« États-Unis d'Afrique »). Le 25 mai \textbf{1963}, l'\cle{OUA} (Organisation de l'Unité Africaine) est fondée à Addis-Abeba. Nkrumah y joue un rôle majeur, même si sa vision d'un gouvernement continental est écartée au profit d'une simple coopération entre États souverains.
\end{itemize}

\courssub{Les limites : dérive autoritaire et chute de Nkrumah}

Le bilan du « père de l'indépendance » comporte une face sombre, qu'il faut connaître pour nuancer son jugement — qualité attendue au Baccalauréat technique.

\begin{itemize}
\item \textbf{Dérive autoritaire} : dès 1960, Nkrumah fait du Ghana une République dont il est le président tout-puissant. En \textbf{1964}, un référendum instaure le \textbf{parti unique} (le CPP) et Nkrumah est proclamé président à vie. La presse est muselée, les opposants emprisonnés (loi de détention préventive de 1958).
\item \textbf{Difficultés économiques} : les grands projets (barrage d'Akosombo, industrialisation forcée) endettent le pays ; la chute du cours du cacao ruine les finances publiques ; la vie chère provoque le mécontentement.
\item \textbf{La chute} : le \textbf{24 février 1966}, alors que Nkrumah est en voyage à Pékin, l'armée et la police le renversent par un \textbf{coup d'État}. Il meurt en exil en Guinée en 1972.
\end{itemize}

\begin{aretenir}[L'essentiel du cas ghanéen]
\begin{itemize}
\item La Gold Coast, colonie britannique riche (cacao, or) dotée d'une élite éduquée et d'une presse libre, est le laboratoire de la décolonisation négociée.
\item Nkrumah fonde le CPP en 1949 et impose l'« action positive » (grèves, boycotts, désobéissance civile non violente).
\item Étapes : victoire électorale de 1951 (34 sièges sur 38) $\rightarrow$ autonomie interne (Constitution de 1954) $\rightarrow$ indépendance le \textbf{6 mars 1957}, première décolonisation d'Afrique subsaharienne.
\item Nkrumah, champion du panafricanisme (Accra 1958, OUA 1963), dérive vers l'autoritarisme et est renversé en 1966.
\end{itemize}
\end{aretenir}

\begin{exercice}[Pour s'entraîner]
\begin{enumerate}
\item Définissez l'« action positive » et montrez en quoi elle s'inspire de Gandhi.
\item Expliquez pourquoi la victoire du CPP aux élections de 1951 a contraint les Britanniques à libérer Nkrumah.
\item Rédigez un paragraphe argumenté (15-20 lignes) répondant à la question : « Le Ghana de Nkrumah : un modèle réussi de décolonisation ? »
\end{enumerate}
\end{exercice}

\begin{corrige}[Exercice]
\textbf{1.} L'action positive est la stratégie de lutte non violente de Nkrumah : grèves générales, boycott des produits britanniques, désobéissance civile. Comme Gandhi en Inde, Nkrumah refuse la violence armée : il mise sur la paralysie économique et la mobilisation des masses pour rendre la colonie ingouvernable, tout en conservant la sympathie de l'opinion internationale.

\textbf{2.} En remportant 34 sièges sur 38 alors que son chef est en prison, le CPP démontre que le peuple soutient massivement l'indépendance immédiate. Maintenir Nkrumah en détention aurait rendu le pays ingouvernable et discrédité la démocratie britannique ; le gouverneur préfère négocier : Nkrumah est libéré et associé au pouvoir.

\textbf{3.} (Éléments attendus) Réussite : décolonisation pacifique et rapide, par les urnes et la négociation, première d'Afrique noire, effet d'entraînement sur le continent, rayonnement panafricain (Accra 1958, OUA 1963). Limites : dérive autoritaire (parti unique en 1964, présidence à vie), échecs économiques, coup d'État de 1966. Conclusion nuancée : modèle pour la \emph{conquête} de l'indépendance, contre-modèle pour son \emph{exercice}.
\end{corrige}

\coursec{L'émancipation des territoires français : le cas de l'Algérie}

\medskip
\noindent
Alors que le Ghana accède à la souveraineté par la voie négociée et le transfert pacifique des pouvoirs, l'Algérie offre un contraste saisissant : la décolonisation y prend la forme d'une \cle{guerre de libération nationale} longue de huit ans, d'une violence inouïe, qui ébranle la IV\textsuperscript{e} République française jusqu'à son effondrement. Cette spécificité tient à un fait structurel majeur : l'Algérie n'est pas une colonie classique, elle est une \cle{colonie de peuplement} intégrée administrativement à la métropole.

\courssub{Une colonie pas comme les autres : l'Algérie française}

\begin{definition}[Colonie de peuplement]
Territoire où la puissance coloniale a installé une population civile européenne nombreuse (les \cle{pieds-noirs}), venue s'y fixer durablement, s'accaparant les meilleures terres et dominant politiquement, économiquement et socialement la population autochtone. En Algérie, cette présence remonte à 1830.
\end{definition}

\begin{definition}[Pieds-noirs]
Terme désignant les Français d'Algérie (et par extension d'Afrique du Nord), nés sur place depuis plusieurs générations. Au moment de la guerre, ils représentent environ \num{1 000 000} de personnes sur \num{9 500 000} habitants (soit plus de 10\,\%), concentrés dans les villes et les plaines agricoles riches (Mitidja, Oranie, Constantinois).
\end{definition}

\medskip
\noindent
\textbf{Une intégration administrative fictive.} Depuis 1848, l'Algérie est divisée en trois départements français (Alger, Oran, Constantine). Théoriquement, les lois de la République s'y appliquent. Dans les faits, le \cle{code de l'indigénat} (aboli seulement en 1946 par la loi Lamine Guèye) et le statut personnel musulman maintiennent les Algériens dans une citoyenneté de seconde zone : ils sont « Français musulmans d'Algérie », soumis à une justice d'exception, exclus de la représentation politique réelle (le collège électoral musulman est minoritaire et truqué) et victimes de discriminations systémiques (accès à l'emploi, à l'éducation, à la santé).

\textbf{La spoliation foncière, moteur du ressentiment.} La colonisation s'est accompagnée d'une confiscation massive des terres : \emph{senatus-consulte} de 1863, lois de 1873 et 1874 (droit de propriété privée imposé, démantèlement de la propriété collective tribale), séquestre des biens \emph{habous}. Résultat : en 1954, les Européens (10\,\% de la population) détiennent près de \num{2 700 000} hectares des meilleures terres arables, tandis que la majorité paysanne algérienne survit sur des terres marginales, en montagne ou dans le steppique, ou forme une main-d'œuvre agricole sous-payée. Cette \cle{prolétarisation} paysanne alimente le vivier de recrutement du nationalisme radical.

\begin{savaistu}[L'Algérie, « partie intégrante de la France »]
Jusqu'en 1962, l'Algérie envoie des députés à l'Assemblée nationale à Paris. Mais le scrutin est biaisé : deux collèges électoraux (un pour les « Français de souche européenne », un pour les « Français musulmans ») aux poids inégaux. En 1954, 1,3 million d'Européens élisent 15 députés, tandis que 9 millions de « Français musulmans » n'en élisent que 15 également. Cette injustice flagrante convainc les nationalistes que la voie légale est bouchée.
\end{savaistu}

\courssub{Les prémices : le choc de Sétif et Guelma (mai 1945)}

\noindent
Le 8 mai 1945, jour de la victoire sur l'Allemagne nazie, des manifestations nationalistes éclatent à Sétif, Guelma et Kherrata pour réclamer l'indépendance, autorisées par le sous-préfet mais interdites au dernier moment. Le port de drapeaux algériens (interdits) et de pancartes « Libérez Messali Hadj » provoque des tirs de police. C'est l'étincelle : la foule s'en prend aux Européens (une centaine de morts). La répression française est féroce, disproportionnée et dure un mois : bombardements aviation, tirs d'artillerie, ratissages, exécutions sommaires dans les villages du Constantinois.

\begin{exemplebox}[Le chiffre impossible : les morts de Sétif-Guelma]
L'écart entre les bilans officiels illustre la violence coloniale et le déni d'État :
\begin{itemize}
    \item \textbf{Bilan officiel français (1945) :} \num{1 500} à \num{2 000} morts « indigènes ».
    \item \textbf{Bilan nationaliste algérien / historiens (Benjamin Stora, etc.) :} entre \num{15 000} et \num{45 000} morts.
\end{itemize}
Ce massacre traumatise la jeunesse algérienne. Les modérés (Ferhat Abbas, UDMA) qui croyaient encore à une évolution dans le cadre français basculent vers l'idée que seule la lutte armée peut faire plier la France. Messali Hadj (PPA/MTLD) est arrêté. Les futurs chefs du FLN (Ben Bella, Bitat, Khider, Aït Ahmed, Boudiaf) tirent la leçon : l'indépendance ne se négociera pas, elle s'arrachera.
\end{exemplebox}

\begin{attention}[Piège chronologique fréquent]
\emph{Ne pas confondre le début de la guerre d'Algérie (1er novembre 1954) avec les massacres de Sétif (8 mai 1945).} Sétif est un \textbf{prélude}, un tournant radicalisateur, mais la guerre d'indépendance à proprement parler commence près de dix ans plus tard. Au Bac, écrire « La guerre d'Algérie commence en 1945 » est une erreur grave de repérage temporel.
\end{attention}

\courssub{Le déclenchement : la Toussaint rouge et la naissance du FLN (1er novembre 1954)}

\begin{definition}[FLN (Front de Libération Nationale)]
Organisation nationaliste algérienne créée en 1954 (proclamation du 1er novembre) par le CRUA (Comité Révolutionnaire d'Unité et d'Action). Elle rassemble des tendances diverses (arabo-islamistes, berbéristes, marxistes) autour d'un objectif unique : l'indépendance totale par la lutte armée. Sa branche armée est l'ALN (Armée de Libération Nationale). Le FLN impose son hégémonie sur le mouvement nationaliste en éliminant physiquement ses rivaux (MNA de Messali Hadj).
\end{definition}

\medskip
\noindent
Dans la nuit du 31 octobre au 1er novembre 1954 (la \cle{Toussaint rouge}), une soixantaine d'attentats coordonnés visent des symboles de la présence française : casernes, gendarmeries, postes de garde, bâtiments publics, domiciles de colons. Le bilan est modeste militairement (7 morts européens, 5 nationalistes), mais l'impact politique est immense. Le FLN diffuse un \emph{Manifeste} appelant le peuple algérien à la lutte pour « la restauration de l'État algérien souverain, démocratique et social dans le cadre des principes islamiques ».

La réponse du gouvernement Mendès France (puis Edgar Faure) est sans ambiguïté : « Les départements algériens font partie de la République française... La sécession est impossible. » C'est le choix de la \cle{guerre} (euphémisée en « opérations de maintien de l'ordre », « pacification », « événements d'Algérie ») plutôt que de la négociation. Le contingent est rappelé, les effectifs passent de \num{50 000} à \num{400 000} hommes en 1956.

\courssub{Les grandes phases d'une guerre sans nom}

\begin{popfigure}
\begin{tikzpicture}[scale=0.75, every node/.style={font=\small}]
    % Contour très simplifié de l'Algérie (polygone approximatif)
    \draw[popdark, thick, fill=popcream] 
        (0,0) -- (10,0) -- (10.5,1) -- (11,3) -- (10.5,5) -- (9,7) -- (7,8) -- (5,7.5) -- (3,8) -- (1.5,7) -- (0.5,5) -- (0,3) -- cycle;
    % Frontières voisines (traits pointillés)
    \draw[dashed, popmuted] (0,0) -- (-1,-0.5) node[left] {Maroc};
    \draw[dashed, popmuted] (10,0) -- (11.5,-0.5) node[right] {Tunisie};
    \draw[dashed, popmuted] (0,0) -- (-0.5,2) node[left] {Mauritanie};
    \draw[dashed, popmuted] (0,0) -- (-0.5,-2) node[below] {Mali/Niger};
    % Villes principales
    \node[circle, fill=popblue, inner sep=2pt, label=above:{Alger}] at (3.5, 6.5) {};
    \node[circle, fill=poporange, inner sep=2pt, label=right:{Sétif}] at (5.5, 5.2) {};
    \node[circle, fill=poporange, inner sep=2pt, label=right:{Constantine}] at (7.5, 5.5) {};
    \node[circle, fill=poporange, inner sep=2pt, label=below:{Oran}] at (1.5, 3.5) {};
    \node[circle, fill=poporange, inner sep=2pt, label=below:{Batna}] at (6.5, 4) {};
    % Ligne Morice (frontière est, vers Tunisie)
    \draw[poppurple, very thick, decorate, decoration={snake, amplitude=1pt, segment length=4pt}] 
        (9,7) -- (10.5,5) -- (11,3) -- (10.5,1);
    \node[poppurple, rotate=-30] at (10.5, 4) {\textbf{Ligne Morice}};
    \node[poppurple, rotate=-30] at (10.5, 3.3) {(barbelés, mines, radar)};
    % Ligne Challe (frontière ouest, vers Maroc) - plus au sud
    \draw[poppurple, thick, decorate, decoration={snake, amplitude=1pt, segment length=4pt}] 
        (0,0) -- (1.5,7);
    \node[poppurple, rotate=30] at (1, 3.5) {\textbf{Ligne Challe}};
    % Zones de maquis (ombrage)
    \fill[popgreen!30] (5,4) circle (1.2); % Aurès
    \node[popgreen!80!black] at (5,4) {Aurès};
    \fill[popgreen!30] (3,6) ellipse ({1.5} and {0.8}); % Kabylie
    \node[popgreen!80!black] at (3,6) {Kabylie};
    \fill[popgreen!30] (7,5.5) circle (0.8); % Constantinois
    \node[popgreen!80!black] at (7,5.5) {Constantinois};
    \fill[popgreen!30] (1.5,3.5) circle (0.8); % Oranie
    \node[popgreen!80!black] at (1.5,3.5) {Oranie};
    % Légende
    \begin{scope}[shift={(11.5, 7)}]
        \draw[popdark, thick, fill=popcream] (0,0) rectangle (3.5, -3.5);
        \node[anchor=west] at (0, 0.2) {\textbf{Légende}};
        \fill[popblue] (0.2, -0.5) circle (2pt); \node[anchor=west] at (0.6, -0.5) {Capitale (Alger)};
        \fill[poporange] (0.2, -1.0) circle (2pt); \node[anchor=west] at (0.6, -1.0) {Villes clés};
        \draw[poppurple, very thick, decorate, decoration={snake}] (0.2, -1.5) -- (1.2, -1.5); \node[anchor=west] at (1.4, -1.5) {Lignes Morice / Challe};
        \fill[popgreen!30] (0.2, -2.0) rectangle (1.0, -2.4); \node[anchor=west] at (1.2, -2.2) {Zones de maquis (ALN)};
        \draw[->, thick, popdark] (1.5, -3.0) -- (2.5, -3.0); \node[anchor=west] at (2.7, -3.0) {Nord};
    \end{scope}
    % Mer Méditerranée
    \node[popblue!50, rotate=90] at (-1, 4) {Méditerranée};
\end{tikzpicture}
\legende{Carte schématique de l'Algérie pendant la guerre (1954-1962) : localisation des principaux foyers de guérilla (maquis), des lignes de barrage électrifiées (Morice à l'est, Challe à l'ouest) et des villes symboliques (Sétif, Alger).}
\end{popfigure}

\medskip
\noindent
La guerre se structure en phases marquées par l'adaptation mutuelle des stratégies :

\textbf{1. La guérilla et la contre-guérilla (1954-1956).} Le FLN structure le territoire en six \cle{wilayas} (zones militaires), chacune commandée par un chef (ex: Krim Belkacem en Wilaya III Kabylie, Ben Boulaïd en Wilaya I Aurès, Ben M'hidi en Wilaya V Oranie). L'ALN pratique la guérilla : embuscades, sabotage, harcèlement des postes isolés. L'armée française répond par le \cle{quadrillage} (création de secteurs, postes de surveillance), les \cle{regroupements de populations} (déplacement forcé de plus de \num{2 000 000} de ruraux vers des camps surveillés pour couper le FLN de sa base logistique) et l'usage massif de la torture (systématisée comme outil de renseignement, notamment par les paras du général Massu).

\textbf{2. La bataille d'Alger (janvier-octobre 1957).} Le FLN lance une offensive urbaine : pose de bombes dans les lieux publics fréquentés par les Européens (stade, cafés, dancing). Le général Massu reçoit les pleins pouvoirs (pouvoirs de police). La 10\textsuperscript{e} DP (division parachutiste) quadrille la Casbah, utilise la torture à grande échelle, « tourne » des militants, décapite l'organisation FLN à Alger (arrestation de Larbi Ben M'hidi, Ali La Pointe, Hassiba Ben Bouali). Victoire tactique française, mais défaite stratégique : l'opinion métropolitaine et internationale découvre l'horreur de la « sale guerre », la torture devient un scandale d'État.

\textbf{3. Le retour de De Gaulle et la bascule (mai 1958 - septembre 1959).} L'insurrection du 13 mai 1958 à Alger (comité de salut public, juntes militaires et pieds-noirs) provoque la chute de la IV\textsuperscript{e} République. De Gaulle revient au pouvoir, acclamé à Alger (« Je vous ai compris », 4 juin 1958). Mais il comprend vite que la guerre est intenable : coût financier exorbitant (20\,\% du budget), isolement diplomatique, fracture morale. Le 16 septembre 1959, il prononce le mot tabou : « \cle{Autodétermination} ». Le FLN crée le GPRA (Gouvernement Provisoire de la République Algérienne) à Tunis (septembre 1958), reconnu par de nombreux pays.

\textbf{4. La fin de partie : putsch des généraux et répression parisienne (1960-1961).}
\begin{itemize}
    \item \textbf{Janvier 1960 :} « Semaine des barricades » à Alger (ultra-colons, mouvements précurseurs de l'OAS).
    \item \textbf{Avril 1961 :} \cle{Putsch des généraux} (Challe, Zeller, Salan, Jouhaud) à Alger pour maintenir l'Algérie française. Échec par la loyauté de la troupe et l'appel de De Gaulle à la télévision.
    \item \textbf{17 octobre 1961 :} Manifestation pacifique du FLN à Paris contre le couvre-feu raciste. Répression brutale par la police de Papon : dizaines, peut-être centaines de morts (jetés dans la Seine), milliers d'arrestations. Longtemps nié, ce crime d'État est reconnu par la France en 2021.
\end{itemize}

\courssub{La fin de la guerre : accords d'Évian, indépendance et drames humains}

\noindent
Les négociations s'ouvrent à Évian (Haute-Savoie) en mars 1962. Le FLN (GPRA) et le gouvernement français signent les \cle{accords d'Évian} le \textbf{18 mars 1962}. Ils prévoient : cessez-le-feu immédiat (19 mars), autodétermination dans 3 mois, garanties pour les Européens (droits civils, propriété), libération des prisonniers, coopération future.

\begin{itemize}
    \item \textbf{8 avril 1962 :} Référendum d'autodétermination en France : 91\,\% de « Oui ».
    \item \textbf{1er juillet 1962 :} Référendum d'autodétermination en Algérie : 99,7\,\% de « Oui ».
    \item \textbf{3 juillet 1962 :} Proclamation officielle de l'indépendance par De Gaulle.
    \item \textbf{5 juillet 1962 :} Fête de l'indépendance (anniversaire de la prise d'Alger en 1830). Ben Bella devient le premier chef de gouvernement (président du Conseil), puis premier Président de la République en septembre 1963.
\end{itemize}

\medskip
\noindent
Mais la joie de l'indépendance est mêlée à deux tragédies humaines majeures :

\begin{aretenir}[L'exode des Pieds-noirs et le sort des Harkis]
\begin{itemize}
    \item \textbf{L'exode des Pieds-noirs :} Craignant les représailles, ne faisant pas confiance aux garanties d'Évian, près de \num{900 000} Européens fuient l'Algérie en quelques mois (été 1962), abandonnant biens, maisons, souvenirs. Accueillis souvent dans la précarité en métropole, ils forment un traumatisme collectif durable (« la valise ou le cercueil »).
    \item \textbf{Le drame des Harkis :} Environ \num{200 000} Algériens ont servi dans l'armée française (harkis, moghaznis, gardes mobiles). Malgré les accords d'Évian interdisant les représailles, des dizaines de milliers sont massacrés sommairement par le FLN ou des populations vengeresses après le cessez-le-feu. Seuls \num{~40 000} à \num{60 000} (avec leurs familles) sont rapatriés en France, parqués dans des camps (Rivesaltes, Bias) dans des conditions indignes. La reconnaissance de ce « devoir de mémoire » ne commence qu'en 2001 (loi sur la reconnaissance nationale).
\end{itemize}
\end{aretenir}

\begin{popfigure}
\begin{tikzpicture}
\begin{axis}[
    ybar,
    bar width=22pt,
    width=\linewidth,
    height=9cm,
    ylabel={Effectifs / Pertes (en milliers)},
    symbolic x coords={Effectifs français (pic 1958), Pertes militaires françaises, Pertes algériennes (est. basse), Pertes algériennes (est. haute)},
    xtick=data,
    xticklabel style={rotate=30, anchor=east, font=\small},
    ymin=0,
    ymax=1600,
    ymajorgrids=true,
    grid style=dashed,
    legend style={at={(0.5,-0.25)}, anchor=north, legend columns=-1, font=\small},
    nodes near coords,
    nodes near coords align={vertical},
    every node near coord/.append style={font=\tiny, rotate=90, anchor=west, yshift=2pt},
]
\addplot[fill=popblue!70] coordinates {
    (Effectifs français (pic 1958), 470)
    (Pertes militaires françaises, 27)
    (Pertes algériennes (est. basse), 300)
    (Pertes algériennes (est. haute), 1500)
};
\legend{Effectifs / Pertes}
\end{axis}
\end{tikzpicture}
\legende{Comparaison des effectifs militaires français au pic de l'engagement (1958) et des pertes humaines. Sources : Histoire de la guerre d'Algérie (B. Stora), Ministère des Armées. Les pertes algériennes (civils + militaires) font l'objet d'estimations très variables (300 000 à 1 500 000).}
\end{popfigure}

\courssub{Bilan et portée internationale}

\begin{aretenir}[Bilan humain et matériel catastrophique]
\begin{itemize}
    \item \textbf{Morts algériens :} Entre \num{300 000} et \num{1 500 000} (civils et combattants), soit 3\,\% à 15\,\% de la population musulmane de 1954. Destructions massives de villages, cultures brûlées, empoisonnement des puits, défoliants (agent orange testé dès 1959).
    \item \textbf{Morts français :} \num{~27 000} militaires tués, \num{~6 000} civils européens tués.
    \item \textbf{Traumatismes :} Torture (des deux côtés, mais systématique côté français), viols, déplacements forcés, exil. La mémoire reste vive et conflictuelle des deux rives de la Méditerranée.
\end{itemize}
\end{aretenir}

\medskip
\noindent
\textbf{Portée internationale.} La guerre d'Algérie devient une cause emblématique de la \cle{décolonisation} et de la \cle{Guerre froide} :
\begin{itemize}
    \item \textbf{Bloc socialiste (URSS, Chine, pays de l'Est) :} Soutien diplomatique, financier, livraison d'armes (via Tunisie, Maroc, Yougoslavie), formation de cadres ALN. Vote constant à l'ONU pour le droit à l'autodétermination.
    \item \textbf{Pays non-alignés (Bandung 1955, Belgrade 1961) :} Soutien politique fort (Nehru, Nasser, Tito, Nkrumah). L'Algérie devient le symbole de la lutte anti-impérialiste.
    \item \textbf{Isolement de la France :} Critiquée par ses alliés de l'OTAN (USA inquiets de la stabilité méditerranéenne), condamnée à l'ONU (commission d'enquête 1957), la France perd son aura de patrie des Droits de l'Homme. De Gaulle utilise l'indépendance pour réorienter la politique étrangère française (indépendance nationale, dialogue Est-Ouest, francophonie).
\end{itemize}

\courssub{Méthode : Construire un paragraphe argumenté à partir de documents}

\begin{methode}[Rédaction d'un paragraphe historique argumenté (type Bac)]
Pour réussir un exercice de commentaire de documents ou une dissertation, appliquez la structure \textbf{IDEE - PREUVE - EXPLICATION - LIEN} (IPEL) à chaque paragraphe :
\begin{enumerate}
    \item \textbf{Idée directrice (Thèse) :} Une phrase affirmative qui répond directement à la question posée. \emph{Ex : « La guerre d'Algérie est d'abord une guerre civile au sein même de la population européenne d'Algérie. »}
    \item \textbf{Preuve (Document / Connaissance) :} Citez explicitement le document (Doc 1, ligne 3...) ou mobilisez un fait précis (chiffre, date, acteur, citation). \emph{Ex : « Le document 3, extrait du discours de De Gaulle du 16 septembre 1959, montre qu'il propose l'autodétermination, provoquant la colère des ultras (Doc 4 : affiche OAS « De Gaulle au poteau »). »}
    \item \textbf{Explication (Analyse) :} Montrez \emph{pourquoi} cette preuve valide votre idée. Analysez le sens, l'intention, le contexte. \emph{Ex : « Cette rupture avec la doctrine de l'intégration (Algérie française) révèle que l'État français accepte enfin la réalité nationaliste, mais divise son propre camp : l'OAS naît de ce sentiment de trahison. »}
    \item \textbf{Lien (Transition) :} Une phrase qui relie ce paragraphe au suivant ou qui nuance. \emph{Ex : « Cependant, cette division ne profite pas seulement au FLN ; elle accélère aussi l'exode des Pieds-noirs... »}
\end{enumerate}
\textbf{Conseil :} Évitez le résumé successif des documents (« Le doc 1 dit... Le doc 2 dit... »). Croisez-les : « Si le doc 1 (rapport militaire) minimise les pertes civiles, le doc 2 (témoignage FLN) les dénonce, ce qui illustre la guerre de l'information... »
\end{methode}

\begin{exoresolu}[Exemple d'application IPEL]
\textbf{Sujet :} « En vous appuyant sur les documents et vos connaissances, montrez que la guerre d'Algérie a été une guerre totale. »
\medskip
\noindent
\textbf{Paragraphe 1 (Dimension militaire/population) :}
\textbf{Idée :} La guerre touche l'ensemble de la société algérienne par la terreur et le déplacement forcé. \\
\textbf{Preuve :} Doc 1 (rapport général Salan 1958) mentionne « 2 millions de personnes regroupées » ; Connaissance : regroupements de peuplement, interdiction de zones, mines antipersonnel (ligne Morice). \\
\textbf{Explication :} L'armée française vise à couper l'ALN de sa base logistique (le peuple) selon la doctrine de la « guerre révolutionnaire » (Lacheroy, Trinquier). La population devient l'enjeu et la cible. \\
\textbf{Lien :} Cette violence de masse s'accompagne d'une répression judiciaire et policière sans précédent en métropole même.
\medskip
\noindent
\textbf{Paragraphe 2 (Dimension métropolitaine/juridique) :}
\textbf{Idée :} L'état d'exception s'étend à la France métropolitaine, suspendant l'État de droit. \\
\textbf{Preuve :} Doc 2 (témoignage Maurice Papon 1961) ; Connaissance : loi de 1956 (pouvoirs spéciaux), censure, torture à Paris (17 oct. 1961), OAS terrorisme (1961-62). \\
\textbf{Explication :} La guerre « sans nom » n'a pas de front clair ; l'ennemi est intérieur. La République utilise des moyens contraires à ses principes (torture, exécutions extra-judiciaires) pour tenter de la gagner. \\
\textbf{Lien :} Cette totalisation de la guerre explique l'effondrement du régime parlementaire en 1958 et le coût moral durable pour la France.
\end{exoresolu}

\coursec{Dossier 2 : L'indépendance de la Guinée Conakry (activité guidée)}

La lutte de l'Algérie montre que l'émancipation des territoires français peut passer par une guerre longue et meurtrière. Mais ce n'est pas la seule voie possible. En Afrique occidentale française (AOF), un territoire va conquérir son indépendance sans effusion de sang, par les urnes : la Guinée. Ce dossier guidé vous apprend à analyser un ensemble de documents comme à l'examen, tout en retenant les faits essentiels d'un événement capital : le « non » guinéen du 28 septembre 1958.

\courssub{Le contexte : de la loi-cadre à la Communauté française}

Après la Seconde Guerre mondiale, la France affaiblie doit composer avec la montée des nationalismes africains. Plutôt que d'accorder tout de suite l'indépendance, elle cherche des solutions intermédiaires.

\begin{definition}[La loi-cadre Defferre (1956)]
La \cle{loi-cadre}, votée le 23 juin 1956 à l'initiative du ministre Gaston Defferre, accorde une large autonomie interne aux territoires d'Afrique noire : suffrage universel, gouvernements locaux dirigés par des Africains, assemblées territoriales élues. Mais la France conserve les domaines dits « régaliens » : la défense, la diplomatie, la monnaie et les grandes orientations économiques. Ce n'est donc pas encore l'indépendance.
\end{definition}

\begin{definition}[La Communauté française (1958)]
La \cle{Communauté française} est une institution créée par la Constitution de la Ve République, voulue par le général de Gaulle en 1958. Elle associe la France et ses territoires d'outre-mer : ceux-ci deviennent des États autonomes, mais la France garde le contrôle de la politique étrangère, de la défense, de la monnaie et des matières premières stratégiques. Le président de la République française est aussi président de la Communauté.
\end{definition}

\begin{definition}[Le référendum]
Un \cle{référendum} est une consultation par laquelle l'ensemble des électeurs répond directement, par oui ou par non, à une question posée par le pouvoir. Le 28 septembre 1958, les peuples de l'empire français sont appelés à approuver la Constitution de la Ve République : voter « oui » signifie entrer dans la Communauté ; voter « non » signifie quitter immédiatement la France et accéder à l'indépendance.
\end{definition}

\begin{attention}[Ne pas confondre les trois Guinées]
Erreur fréquente à l'examen : il existe trois pays appelés « Guinée ». La \cle{Guinée Conakry} (notre sujet) est une ancienne colonie \textbf{française}, indépendante le 2 octobre 1958. La \cle{Guinée-Bissau} est une ancienne colonie \textbf{portugaise} (indépendante en 1974). La \cle{Guinée équatoriale} est une ancienne colonie \textbf{espagnole} (indépendante en 1968). Précisez toujours « Guinée Conakry » dans vos copies.
\end{attention}

\courssub{Sékou Touré, le PDG et le discours du 25 août 1958}

En Guinée, le mouvement nationaliste est dominé par le \cle{Parti Démocratique de Guinée} (PDG), section guinéenne du Rassemblement Démocratique Africain (RDA). Son chef, \cle{Ahmed Sékou Touré}, syndicaliste de formation, est devenu en 1957 vice-président du conseil de gouvernement grâce à la loi-cadre. Il incarne une génération de leaders issus non des élites traditionnelles, mais des syndicats et des masses populaires.

En août 1958, le général de Gaulle entreprend une tournée africaine pour défendre la Communauté. Le 25 août 1958, à Conakry, Sékou Touré prononce devant lui un discours resté célèbre. Il y affirme que l'Afrique ne demande pas la charité mais la dignité, et conclut par une phrase historique : « \emph{Nous préférons la liberté dans la pauvreté à la richesse dans l'esclavage.} »

\begin{experience}[Lecture guidée d'un extrait du discours de Conakry (25 août 1958)]
\textbf{Document :} « Nous préférons la liberté dans la pauvreté à la richesse dans l'esclavage [...] Nous exercerons notre droit à l'indépendance, mais nous resterons proches de la France. » (Sékou Touré, discours devant le général de Gaulle, Conakry, 25 août 1958.)
\begin{enumerate}
\item \textbf{Observer :} qui parle ? devant qui ? à quelle date ? (Un leader nationaliste africain s'adresse au chef de l'État français, un mois avant le référendum.)
\item \textbf{Dégager l'idée principale :} Sékou Touré annonce le refus de la Communauté au nom de la liberté, tout en souhaitant garder des relations amicales avec la France.
\item \textbf{Interpréter :} la formule oppose deux couples (liberté/pauvreté contre richesse/esclavage) : la dignité politique prime sur les avantages matériels de la colonisation.
\item \textbf{Conclure :} le discours prépare l'opinion guinéenne au vote « non » et annonce la rupture à venir.
\end{enumerate}
\end{experience}

\courssub{Le référendum du 28 septembre 1958 : le triomphe du « non »}

Le 28 septembre 1958, les électeurs guinéens se prononcent massivement. Le PDG a fait campagne pour le « non », c'est-à-dire pour l'indépendance immédiate, tandis que la France et ses alliés locaux défendent le « oui ».

\begin{exemplebox}[Des résultats sans appel]
En Guinée, le « non » l'emporte avec environ \num{1136000} voix contre environ \num{57000} voix pour le « oui », soit près de \textbf{95 \%} des suffrages exprimés. La Guinée est le \cle{seul territoire de l'AOF} (et de toute l'Afrique française subsaharienne) à refuser la Communauté. Partout ailleurs — Sénégal, Côte d'Ivoire, Gabon, etc. — le « oui » l'emporte largement.
\end{exemplebox}

\begin{popfigure}
\begin{center}
\begin{tikzpicture}
% Diagramme circulaire : Non 95 %, Oui 5 %
\fill[popgreen] (0,0) -- (90:2.6) arc (90:-252:2.6) -- cycle;
\fill[poporange] (0,0) -- (-252:2.6) arc (-252:-270:2.6) -- cycle;
\node[white,font=\bfseries] at (-35:1.5) {NON 95 \%};
\node[popdark,font=\bfseries] at (270:3.15) {OUI 5 \%};
\node[font=\small] at (0,-3.4) {Environ \num{1136000} voix « non » contre \num{57000} voix « oui ».};
\end{tikzpicture}
\end{center}
\legende{Résultats du référendum du 28 septembre 1958 en Guinée : le « non » à la Communauté française (indépendance immédiate) recueille près de 95 \% des suffrages.}
\end{popfigure}

\courssub{La réaction française : une rupture brutale}

De Gaulle avait prévenu : le « non » entraînerait des « conséquences ». La France applique une séparation d'une rare brutalité, destinée à punir la Guinée et à décourager les autres territoires tentés de l'imiter.

\begin{itemize}
\item \textbf{Rupture immédiate :} toute coopération est suspendue dès l'annonce des résultats ; la Guinée est exclue de la zone franc et des aides françaises.
\item \textbf{Retrait précipité :} environ \num{3000} fonctionnaires français (administrateurs, médecins, enseignants, techniciens) quittent le pays en quelques semaines, emportant le matériel administratif.
\item \textbf{Sabotages :} des témoignages concordants décrivent le démontage des lignes téléphoniques, la destruction ou l'enlèvement d'archives, de plans et de dossiers, voire de matériel hospitalier et scolaire.
\end{itemize}

\begin{savaistu}[Une vengeance jusqu'aux ampoules]
En représailles, certains fonctionnaires français quittant la Guinée emportèrent tout ce qui pouvait l'être : on rapporte qu'ils arrachèrent jusqu'aux \cle{ampoules} des bâtiments publics et emportèrent les \cle{médicaments} des hôpitaux. Cet épisode, resté dans les mémoires guinéennes, illustre la volonté de faire de la Guinée un « exemple » de ce qu'il en coûte de dire non à la France.
\end{savaistu}

\courssub{L'indépendance du 2 octobre 1958 et ses conséquences}

Le 2 octobre 1958, quatre jours après le référendum, Sékou Touré proclame l'indépendance de la République de Guinée et en devient le premier président. La jeune nation, abandonnée par la France, cherche aussitôt des appuis à l'étranger.

\begin{itemize}
\item \textbf{Reconnaissances immédiates :} l'\cle{URSS}, la \cle{Chine} populaire et le \cle{Ghana} de Kwame Nkrumah reconnaissent aussitôt la Guinée. Nkrumah accorde même un prêt substantiel (environ 10 millions de livres sterling) pour aider le pays à survivre.
\item \textbf{Isolement diplomatique :} la France fait pression pour retarder l'admission de la Guinée à l'ONU (obtenue en décembre 1958) et la tient à l'écart des circuits économiques occidentaux.
\item \textbf{Rapprochement avec le bloc de l'Est :} privée de l'aide française, la Guinée se tourne vers l'URSS, la Chine et les pays socialistes pour l'équipement, la formation et le commerce.
\item \textbf{Radicalisation du régime :} le PDG devient parti unique ; l'économie est étatisée ; l'opposition est réprimée. Le régime de Sékou Touré évolue vers un pouvoir autoritaire qui durera jusqu'à sa mort en 1984.
\end{itemize}

\begin{popfigure}
\begin{center}
\begin{tikzpicture}[mindmap, grow cyclic,
  every node/.style={concept, text=white},
  root concept/.append style={concept color=popblue, font=\bfseries},
  level 1/.append style={level distance=3.6cm, sibling angle=120},
  level 1 concept/.append style={font=\small\bfseries},
  level 2/.append style={level distance=2.4cm, sibling angle=45},
  level 2 concept/.append style={font=\scriptsize}]
\node[root concept]{Le « non » guinéen du 28 sept. 1958}
  child[concept color=popgreen]{ node{Politiques}
    child{ node{Indépendance le 2 oct. 1958} }
    child{ node{Parti unique (PDG)} }
    child{ node{Régime autoritaire} }
  }
  child[concept color=poporange]{ node{Économiques}
    child{ node{Retrait des aides et fonctionnaires} }
    child{ node{Sabotages, exclusion de la zone franc} }
    child{ node{Prêt du Ghana ; étatisation} }
  }
  child[concept color=poppurple]{ node{Diplomatiques}
    child{ node{Isolement face à la France} }
    child{ node{Reconnaissance URSS, Chine, Ghana} }
    child{ node{ONU en déc. 1958} }
  };
\end{tikzpicture}
\end{center}
\legende{Carte mentale des conséquences du « non » guinéen : politiques (indépendance, parti unique), économiques (rupture des aides, appuis nouveaux), diplomatiques (isolement puis ouverture vers le bloc de l'Est).}
\end{popfigure}

\courssub{Travail guidé : exploiter un dossier documentaire au Bac}

\begin{methode}[Analyser un dossier documentaire en quatre étapes]
\begin{enumerate}
\item \textbf{Dégager le thème :} lire les titres, sources et dates de chaque document pour identifier le sujet commun (ici : les conditions de l'indépendance de la Guinée).
\item \textbf{Problématiser :} formuler une question directrice. Exemple : « Pourquoi et comment la Guinée a-t-elle accédé seule à l'indépendance en 1958, et à quel prix ? »
\item \textbf{Croiser les documents :} confronter les points de vue (discours de Sékou Touré, chiffres du référendum, témoignages sur le départ français) pour montrer ce qu'ils confirment ou nuancent mutuellement.
\item \textbf{Conclure :} répondre à la problématique en une phrase synthétique qui hiérarchise les causes et les conséquences.
\end{enumerate}
\end{methode}

\begin{exoresolu}[Étude de documents : le prix de la liberté]
\textbf{Énoncé.} Document 1 : « Nous préférons la liberté dans la pauvreté à la richesse dans l'esclavage » (Sékou Touré, 25 août 1958). Document 2 : tableau des résultats du référendum : « non » \num{1136000} voix ; « oui » \num{57000} voix. Document 3 : témoignage : « Les Français sont partis en emportant les dossiers, les téléphones, jusqu'aux médicaments des hôpitaux. » Question : montrez que l'indépendance de la Guinée fut un choix libre mais coûteux.
\tcblower
\textbf{Solution détaillée.} \emph{Choix libre :} le document 1 montre que le leader guinéen revendique la liberté avant tout avantage matériel ; le document 2 prouve que ce choix est massivement populaire (environ 95 \% de « non », soit un rapport d'environ 20 voix contre 1). La Guinée est le seul territoire de l'AOF à rejeter la Communauté, ce qui souligne le caractère volontaire et singulier de sa décision. \emph{Choix coûteux :} le document 3 révèle la riposte française : départ brutal des fonctionnaires, sabotages, exclusion de la zone franc. La Guinée indépendante (2 octobre 1958) doit alors survivre grâce au prêt du Ghana et à l'aide du bloc de l'Est, ce qui l'isole et pousse son régime vers la radicalisation. \emph{Conclusion :} la Guinée a payé sa liberté d'un isolement économique et diplomatique qui a durablement marqué son histoire.
\end{exoresolu}

\begin{aretenir}[L'essentiel du dossier]
La loi-cadre de 1956 donne l'autonomie interne ; la Communauté de 1958 propose une indépendance surveillée. Le 28 septembre 1958, la Guinée de Sékou Touré vote « non » à près de 95 \% : seule colonie française d'Afrique noire à choisir l'indépendance immédiate, proclamée le 2 octobre 1958. La France riposte par une rupture brutale (retrait des fonctionnaires, sabotages). Soutenue par le Ghana, l'URSS et la Chine, la Guinée s'isole, se rapproche du bloc socialiste et son régime se radicalise.
\end{aretenir}

\begin{exercice}[Pour s'entraîner]
1. Définis la loi-cadre de 1956 et la Communauté française en deux phrases chacune. 2. Explique la signification du discours de Sékou Touré du 25 août 1958. 3. À partir des chiffres du référendum, calcule le pourcentage exact du « non » et commente. 4. Rédige un paragraphe argumenté (10 à 15 lignes) répondant à la question : « L'indépendance de la Guinée : victoire ou piège ? »
\end{exercice}

\begin{corrige}[Exercice]
1. La loi-cadre (1956) accorde l'autonomie interne aux territoires africains (suffrage universel, gouvernements locaux) sans leur donner la souveraineté. La Communauté (1958) associe ces territoires à la France, qui conserve défense, diplomatie et monnaie. 2. Le discours affirme la priorité de la liberté sur les avantages matériels et annonce le vote « non ». 3. $\dfrac{1136000}{1136000+57000}\times 100 \approx 95{,}2\,\%$ : rejet massif de la Communauté. 4. Attendu : victoire (dignité, indépendance immédiate, exemple pour l'Afrique) mais coût lourd (sabotages, isolement, dépendance envers le bloc de l'Est, autoritarisme) ; conclusion nuancée exigée.
\end{corrige}

\coursec{L'émancipation des autres territoires : le cas de l'Angola}

Après avoir analysé le cas particulier de la Guinée, qui a accédé à l'indépendance dès 1958 par la voie du « non » au référendum gaulliste, nous nous tournons maintenant vers l'Afrique australe et l'Angola. Ce territoire portugais illustre une tout autre dynamique : celle d'une décolonisation tardive, arrachée par la lutte armée face à un pouvoir métropolitain refusant obstinément toute évolution, avant de basculer dans une guerre civile internationalisée par la guerre froide.

\courssub{Le contexte portugais : une dictature coloniale figée}

Contrairement à la France et au Royaume-Uni, qui ont engagé (plus ou moins volontairement) un processus de décolonisation à partir des années 1950, le Portugal maintient sous la dictature d'António de Oliveira \cle{Salazar} (au pouvoir de 1932 à 1968, puis Marcelo Caetano jusqu'en 1974) une conception impériale intransigeante. La Constitution de 1933 (État nouveau) proclame l'indivisibilité de la patrie portugaise : les colonies ne sont pas des territoires distincts mais des « provinces d'outre-mer ».

\begin{definition}[Province d'outre-mer et assimilation forcée]
En 1951, la révision constitutionnelle remplace le terme « colonie » par « province d'outre-mer » pour l'Angola, le Mozambique, la Guinée-Bissau, le Cap-Vert et São Tomé-et-Principe. Cette qualification juridique vise à contourner la Charte des Nations unies (Chapitre XI sur les territoires non autonomes) et à nier le droit des peuples à l'autodétermination. Le régime impose l'assimilation culturelle (langue portugaise, catholicisme) et maintient le statut d'\cle{indigénat} (travail forcé, code de l'indigénat) pour la majorité africaine, tandis qu'une minorité d'\cle{assimilados} accède à la citoyenneté portugaise sous conditions strictes.
\end{definition}

\begin{savaistu}[L'illusion du « luso-tropicalisme »]
Le régime salazariste promeut l'idéologie du « luso-tropicalisme » (théorisée par Gilberto Freyre) : le Portugal serait un colonisateur « doux », métissé, non raciste, capable d'une harmonisation multiraciale. Cette propagande masque une réalité de ségrégation spatiale (quartiers européens vs \emph{musseques} africains), d'exploitation économique (cultures de rente : café, coton, diamants au profit de la métropole) et de répression policière (PIDE, police politique).
\end{savaistu}

\courssub{Les mouvements de libération : une opposition fragmentée et internationalisée}

Face au refus de tout dialogue, trois mouvements nationalistes émergent, portés par des bases ethno-régionales et des soutiens extérieurs distincts, ce qui préfigure les conflits post-indépendance.

\begin{center}
\begin{tabularx}{\linewidth}{|l|X|X|X|}
\hline
\rowcolor{popblueL}
\textbf{Mouvement} & \textbf{Dirigeant historique} & \textbf{Base sociologique / régionale} & \textbf{Appuis extérieurs principaux} \\
\hline
\cle{MPLA} (Mouvement populaire de libération de l'Angola) & Agostinho \cle{Neto} & Urbaine, mestiços, ambundu (région de Luanda, Malanje) & URSS, bloc de l'Est, \cle{Cuba} \\
\hline
\cle{FNLA} (Front national de libération de l'Angola) & Holden \cle{Roberto} & Bakongo (nord, région de l'Uíge, Zaïre frontalier) & USA (via CIA), \cle{Zaïre} (Mobutu), Chine (initialement) \\
\hline
\cle{UNITA} (Union nationale pour l'indépendance totale de l'Angola) & Jonas \cle{Savimbi} & Ovimbundu (centre, plateau du Bié) & \cle{Afrique du Sud} (apartheid), USA (après 1975) \\
\hline
\end{tabularx}
\end{center}

\begin{aretenir}[Fragmentation nationaliste]
L'absence de front unique (contrairement au FLN algérien ou au PAIGC en Guinée-Bissau) s'explique par des rivalités ethno-régionales, des divergences idéologiques (marxisme-léninisme pour le MPLA, nationalisme conservateur pour FNLA/UNITA) et la logique de la guerre froide : chaque parrain extérieur finance son client, durcissant les clivages.
\end{aretenir}

\courssub{Le déclenchement de la lutte armée (1961)}

La répression d'une grève des cotonculteurs à Baixa de Cassanje (janvier 1961) et le massacre qui s'ensuit (des milliers de morts) radicalisent les nationalistes. Deux événements marquent le début de la guerre :

\begin{enumerate}
\item \textbf{4 février 1961 (Luanda) :} Attaque du MPLA contre la prison de São Paulo, le poste de police et la caserne de Catete. Échec militaire, mais rupture symbolique : début de la guérilla urbaine.
\item \textbf{15 mars 1961 (Nord, Uíge) :} Insurrection paysanne du café dirigée par la FNLA (alors UPA, União das Populações de Angola). Attaques de fermes portugaises, massacres de colons et de populations africaines « loyales ». Répression massive de l'armée portugaise (bombardements au napalm, regroupements de villages). Des dizaines de milliers de réfugiés fuient vers le Zaïre.
\end{enumerate}

\begin{definition}[Guérilla et guerre de libération nationale]
La \cle{guérilla} (petite guerre) est une forme de guerre asymétrique menée par des unités mobiles, légères, s'appuyant sur la population civile, évitant les combats frontaux contre une armée régulière supérieure. La \cle{guerre de libération nationale} désigne le conflit armé mené par un peuple colonisé pour conquérir son indépendance politique, reconnu comme légitime par l'ONU (Résolution 1514, 1960) et le droit international humanitaire (Protocole I additionnel aux Conventions de Genève, 1977).
\end{definition}

\begin{popfigure}
\begin{tikzpicture}[scale=0.85]
% Contour simplifié de l'Angola (polygone approximatif)
\draw[popdark, thick, fill=popcream] 
(0,0) -- (5,0.5) -- (6,2) -- (5.5,4) -- (4,5) -- (2,5.5) -- (0.5,4.5) -- (-0.5,3) -- (-0.5,1) -- cycle;

% Villes
\fill[popblue] (2.5, 1.2) circle (3pt) node[right, popdark] {Luanda (MPLA)};
\fill[poporange] (4.5, 3.8) circle (3pt) node[right, popdark] {Uíge (FNLA)};
\fill[popgreen] (3.5, 2.5) circle (3pt) node[above, popdark] {Huambo (UNITA)};

% Zones d'influence (hachures légères)
\pattern[pattern=north east lines, pattern color=popblue!40] (1,0.5) -- (3.5,1.5) -- (3,3) -- (1,3.5) -- cycle;
\node[popblue, font=\bfseries\small] at (1.5, 2) {Zone MPLA};

\pattern[pattern=north west lines, pattern color=poporange!40] (4,3) -- (6,2.5) -- (5.5,4.5) -- (3.5,4.5) -- cycle;
\node[poporange, font=\bfseries\small] at (4.8, 3.8) {Zone FNLA};

\pattern[pattern=crosshatch dots, pattern color=popgreen!40] (2.5,2) -- (4.5,2.5) -- (4,4) -- (2,4) -- cycle;
\node[popgreen, font=\bfseries\small] at (3.2, 3) {Zone UNITA};

% Frontières voisines
\node[popmuted, font=\small] at (-1.2, 2) {Congo};
\node[popmuted, font=\small] at (6.5, 2) {Zaïre};
\node[popmuted, font=\small] at (3, -0.5) {Namibie (Afq. Sud)};
\node[popmuted, font=\small] at (3, 5.5) {Zambie};

% Flèches appuis extérieurs
\draw[->, thick, popblue] (0.5, 0.5) -- (-1, 0) node[left, popblue, font=\small] {URSS / Cuba};
\draw[->, thick, poporange] (5.5, 3.5) -- (7, 3.5) node[right, poporange, font=\small] {USA / Zaïre};
\draw[->, thick, popgreen] (3.5, 0.5) -- (3.5, -0.5) node[below, popgreen, font=\small] {Afrique du Sud};

% Légende intégrée
\begin{scope}[shift={(-1.8, -0.8)}]
\draw[popblue, thick, pattern=north east lines, pattern color=popblue!40] (0,0) rectangle (0.5,0.3) node[right, font=\small] {MPLA};
\draw[poporange, thick, pattern=north west lines, pattern color=poporange!40] (0,-0.5) rectangle (0.5,-0.2) node[right, font=\small] {FNLA};
\draw[popgreen, thick, pattern=crosshatch dots, pattern color=popgreen!40] (0,-1) rectangle (0.5,-0.7) node[right, font=\small] {UNITA};
\end{scope}
\end{tikzpicture}
\legende{Carte schématique de l'Angola : zones d'influence des trois mouvements nationalistes et appuis extérieurs (vers 1974-1975).}
\end{popfigure}

\courssub{La guerre coloniale (1961-1974) : un fardeau insoutenable pour le Portugal}

Pendant treize ans, le Portugal mène simultanément trois guerres coloniales (Angola, Guinée-Bissau, Mozambique). L'effort militaire mobilise jusqu'à 150 000 hommes (conscription obligatoire, service de 4 ans dont 2 au front).

\begin{exemplebox}[Le coût économique de l'empire]
Au début des années 1970, l'ensemble des trois guerres coloniales absorbe environ \textbf{40 \% du budget de l'État portugais} et 7 à 8 \% du PIB. Ce fardeau asphyxie l'économie métropolitaine (la plus pauvre d'Europe occidentale), provoque l'émigration massive de jeunes Portugais (fuite du service militaire) et alimente le mécontentement au sein même de l'armée (mouvements des \cle{capitaines d'avril}).
\end{exemplebox}

\begin{attention}[Piège chronologique fréquent]
\emph{Erreur :} « L'Angola devient indépendant en 1960 comme les autres pays africains. »
\emph{Correction :} L'indépendance de l'Angola est proclamée le \textbf{11 novembre 1975}, soit 15 ans après « l'année de l'Afrique » (1960). La décolonisation portugaise est la \textbf{dernière} grande vague en Afrique (avec le Mozambique, la Guinée-Bissau, le Cap-Vert, São Tomé). Ne pas généraliser la date de 1960 à l'ensemble du continent.
\end{attention}

\courssub{La révolution des Œillets (25 avril 1974) : le tournant de Lisbonne}

L'usure de la guerre provoque un coup d'État militaire presque sans violence. Le \cle{Mouvement des Forces Armées (MFA)} renverse Marcelo Caetano.

\begin{savaistu}[Pourquoi « Révolution des Œillets » ?]
Le nom vient des fleurs (\cle{œillets}) que la population de Lisbonne offrit aux soldats insurgés, glissant les tiges dans les canons des fusils et les canons des chars. Ce geste symbolisa le caractère pacifique et populaire du renversement de la dictature la plus ancienne d'Europe (42 ans). C'est la fin de l'État nouveau et l'ouverture immédiate des négociations de décolonisation.
\end{savaistu}

\courssub{Les accords d'Alvor et l'indépendance (janvier - novembre 1975)}

Le nouveau pouvoir portugais (gouvernement de transition) convoque les trois mouvements à \cle{Alvor} (Portugal) en janvier 1975.

\begin{propriete}[Accords d'Alvor (15 janvier 1975)]
\begin{itemize}
\item Gouvernement de transition tripartite (MPLA, FNLA, UNITA + Portugal) jusqu'aux élections d'octobre 1976.
\item Cessez-le-feu, intégration des guérillas dans une armée unique.
\item Indépendance fixée au \textbf{11 novembre 1975}.
\end{itemize}
\end{propriete}

Mais la confiance est inexistante. Dès mars 1975, les combats reprennent à Luanda entre MPLA et FNLA. L'Afrique du Sud intervient militairement (opération \cle{Savannah}) au sud en appui à l'UNITA (août-octobre 1975). Le MPLA, menacé, appelle Cuba à l'aide. Début novembre, l'opération \cle{Carlota} déploie des milliers de soldats cubains (jusqu'à 36 000 en 1976).

\begin{popfigure}
\begin{tikzpicture}
% Frise chronologique 1961-1975
\draw[popdark, thick] (0,0) -- (16,0);
% Années
\foreach \x/\year in {0/1961, 2/1963, 4/1965, 6/1967, 8/1969, 10/1971, 12/1973, 14/1975} {
    \draw (\x, -0.2) -- (\x, 0.2);
    \node[below, font=\small] at (\x, -0.4) {\year};
}
% Événements au-dessus
\node[above, align=center, font=\small, popblue] at (0, 1.2) {4 fév: Attaque\\ Luanda (MPLA)};
\node[above, align=center, font=\small, poporange] at (0.5, 2) {15 mars: Insurrection\\ café Nord (FNLA)};
\node[above, align=center, font=\small, popdark] at (4, 1.2) {Guerre coloniale\\ (3 fronts)};
\node[above, align=center, font=\small, poppurple] at (10, 2) {1973: Assassinat\\ d'Amílcar Cabral\\(PAIGC)};
\node[above, align=center, font=\small, popgreen] at (13.5, 1.2) {25 avril 74:\\ Révolution\\ des Œillets};
\node[above, align=center, font=\small, popblue] at (14.5, 2) {Jan 75:\\ Accords\\ d'Alvor};
\node[above, align=center, font=\small, popdark] at (15.8, 1.2) {11 nov 75:\\ Indépendance};
% Événements en-dessous
\node[below, align=center, font=\small, popmuted] at (2, -1.5) {Guérilla,\\ répression};
\node[below, align=center, font=\small, popmuted] at (6, -1.5) {Coût 40\% budget\\ Portugal (3 guerres)};
\node[below, align=center, font=\small, popmuted] at (11, -1.5) {Chute Salazarisme\\ Négociations};
\node[below, align=center, font=\small, popmuted] at (15, -1.5) {Guerre civile\\ + Cuba / Afq Sud};
\end{tikzpicture}
\legende{Frise chronologique : de la révolte de 1961 à l'indépendance de 1975.}
\end{popfigure}

\courssub{Conséquences : une indépendance confisquée par la guerre civile}

Le 11 novembre 1975, le MPLA proclame unilatéralement la \cle{République populaire d'Angola} à Luanda (reconnue par le bloc de l'Est). La FNLA et l'UNITA proclament une république rivale à Huambo. La guerre civile s'internationalise :

\begin{itemize}
\item \textbf{Intervention cubaine massive} (soutien soviétique) : sécurise Luanda, repousse l'Afrique du Sud et la FNLA/UNITA.
\item \textbf{Intervention sud-africaine} : vise à installer un régime ami à sa frontière (Namibie occupée), détruit les infrastructures angolaises.
\item \textbf{Géopolitique de la guerre froide} : L'Angola devient un « champ de bataille par procuration » (proxy war) Est-Ouest jusqu'en 1988 (accords de New York / tripartite) et 1991 (accords de Bicesse), puis reprise jusqu'en 2002 (mort de Savimbi).
\end{itemize}

\begin{exoresolu}[Analyse d'un document : Extrait des Accords d'Alvor (Art. 1 et 3)]
\textbf{Énoncé :} « Les Parties conviennent que l'indépendance de l'Angola sera proclamée le 11 novembre 1975... Le gouvernement de transition sera composé de représentants des trois mouvements... » Expliquez pourquoi cet accord n'a pas empêché la guerre civile.

\tcblower
\textbf{Solution détaillée :}
\begin{enumerate}
\item \textbf{Contexte de méfiance :} Les trois mouvements se sont battus entre eux pendant la guerre anti-coloniale (ex: massacres de 1974-75 à Luanda). L'accord est imposé par le Portugal pressé de partir, non négocié librement.
\item \textbf{Partage du pouvoir illusoire :} Le gouvernement tripartite (Premier ministre tournant, ministères partagés) paralyse la décision. Chaque mouvement conserve son armée, son territoire, son appareil de propagande.
\item \textbf{Interventionnisme extérieur :} L'Afrique du Sud envahit le sud (octobre 75) avant l'indépendance ; Cuba intervient massivement (novembre 75). Les parrains extérieurs ne respectent pas l'esprit d'Alvor (non-ingérence).
\item \textbf{Problème de la légitimité :} Le MPLA contrôle la capitale (Luanda) et l'appareil d'État au jour J. Il proclame la République populaire, marginalisant les deux autres qui refusent de reconnaître cette proclamation unilatérale.
\end{enumerate}
\emph{Conclusion :} L'accord d'Alvor a réglé le départ du Portugal, mais pas le partage du pouvoir entre rivaux armés et parrainés par des puissances étrangères antagonistes.
\end{exoresolu}

\courssub{Méthode : Construire un croquis de synthèse sur la décolonisation de l'Angola}

\begin{methode}[Croquis de synthèse : Décolonisation et guerre civile en Angola]
\textbf{Objectif :} Restituer visuellement les acteurs, les territoires, la chronologie et les enjeux géopolitiques.
\begin{enumerate}
\item \textbf{Titre précis, daté, orienté :} Ex: « Angola : de la guerre de libération à la guerre civile (1961-1975) ».
\item \textbf{Fond de carte simplifié :} Contour du pays, capitale (Luanda), frontières (Zaïre, Zambie, Namibie/Afrique du Sud, Atlantique).
\item \textbf{Zones d'influence (hachures/couleurs) :} MPLA (côte/centre-nord), FNLA (nord/Zaire), UNITA (centre/sud-est). Légende claire.
\item \textbf{Flux extérieurs (flèches proportionnelles) :} 
\begin{itemize}
    \item Flèches entrantes : URSS/Cuba $\to$ MPLA ; USA/Zaïre $\to$ FNLA ; Afq Sud/USA $\to$ UNITA.
    \item Flèches sortantes : Réfugiés, retrait portugais (1975).
\end{itemize}
\item \textbf{Repères chronologiques (jalons) :} 1961 (début), 1974 (Œillets), Jan 75 (Alvor), Nov 75 (Indépendance + intervention Cuba/Afq Sud).
\item \textbf{Légende organisée :} 
\begin{itemize}
    \item \textit{Acteurs internes} (symboles points/carrés pour mouvements).
    \item \textit{Acteurs externes} (drapeaux ou codes pays sur flèches).
    \item \textit{Conflits} (éclairs ou zones hachurées pour combats 1975).
\end{itemize}
\item \textbf{Échelle et Nord :} Indiqués (même approximatifs).
\end{enumerate}
\end{methode}

\begin{exercice}[Croquis guidé]
À partir de la carte et de la frise ci-dessus, réalisez sur une feuille A4 un croquis de synthèse intitulé : \emph{« L'Angola : décolonisation tardive et internationalisation du conflit (1961-1975) »}. Respectez la méthode en 6 étapes. La légende doit faire apparaître : les 3 mouvements, leurs parrains, les 2 interventions militaires majeures de 1975 (Cuba, Afrique du Sud), la date d'indépendance.
\end{exercice}

\begin{corrige}[Exercice n°1 - Éléments attendus dans le croquis]
\begin{itemize}
\item \textbf{Titre} centré, daté, orienté (flèche Nord).
\item \textbf{Contour} Angola + pays limitrophes nommés.
\item \textbf{3 zones colorées/hachurées} (MPLA bleu Luanda/centre, FNLA orange Nord/Uíge, UNITA vert Centre-Sud/Huambo).
\item \textbf{Flèches entrantes} : Cuba/URSS vers MPLA (grosse flèche) ; Zaïre/USA vers FNLA ; Afq Sud vers UNITA.
\item \textbf{Symboles de combats} (éclairs) sur Luanda (mars 75), Sud (opération Savannah oct 75), Nord (janv 75).
\item \textbf{Jalon « 11 nov 1975 : Indépendance »} sur Luanda.
\item \textbf{Légende structurée} en 3 blocs : Mouvements / Appuis / Événements 1975.
\item \textbf{Échelle graphique} (0 - 500 km).
\end{itemize}
\end{corrige}

\begin{aretenir}[Bilan : L'Angola, décolonisation par les armes et guerre froide]
L'Angola illustre le \textbf{modèle portugais} : décolonisation tardive (1975), imposée par l'effondrement du régime métropolitain (Révolution des Œillets) après 13 ans de guerre coloniale ruineuse. La fragmentation nationaliste (MPLA, FNLA, UNITA), encouragée par les rivalités de la guerre froide, a transformé l'accession à l'indépendance en guerre civile internationale (Cuba vs Afrique du Sud), dont le pays ne sortira vraiment qu'en 2002. Ce cas contraste avec les décolonisations négociées (Ghana) ou référendaires (Guinée), et même avec la guerre d'Algérie (un seul front nationaliste dominant, le FLN).
\end{aretenir}

\coursec{L'Afrique sous l'emprise du néocolonialisme}

La proclamation solennelle de l'indépendance de l'Angola en novembre 1975, après quatorze ans de guerre de libération, marque la fin formelle de l'empire colonial portugais en Afrique. Pourtant, pour l'immense majorité des nouveaux États africains, l'accession à la souveraineté internationale ne s'est pas traduite par une autonomie effective. Derrière le drapeau et l'hymne national, d'anciens liens de subordination se sont recomposés sous des formes plus subtiles, mais tout aussi contraignantes. C'est ce phénomène que Kwame Nkrumah, premier président du Ghana, a théorisé dès 1965 sous le terme de \cle{néocolonialisme} : la survie du système colonial sans la responsabilité coloniale.

\courssub{Définition et concept du néocolonialisme}

\begin{definition}[Néocolonialisme]
Le \textbf{néocolonialisme} désigne l'ensemble des mécanismes par lesquels d'anciennes puissances coloniales (ou de nouvelles puissances hégémoniques) maintiennent une domination économique, politique, militaire et culturelle sur des États formellement indépendants. Le terme a été popularisé par \textbf{Kwame Nkrumah} dans son ouvrage \emph{Néocolonialisme, dernier stade de l'impérialisme} (1965). Il y écrit : « L'essence du néocolonialisme est que l'État qui en est victime est, en théorie, indépendant et possède tous les attributs extérieurs de la souveraineté internationale. En réalité, son système économique et par conséquent sa politique sont dirigés de l'extérieur. »
\end{definition}

\begin{aretenir}[Les trois piliers de la domination néocoloniale]
\begin{itemize}
    \item \textbf{Économique :} contrôle des matières premières, des circuits commerciaux, de la monnaie et de l'endettement.
    \item \textbf{Politico-militaire :} accords de défense, bases militaires, ingérence dans le choix des dirigeants, réseaux d'influence (Françafrique).
    \item \textbf{Culturel :} maintien des langues officielles, systèmes éducatifs calqués sur le modèle métropole, médias sous influence.
\end{itemize}
\end{aretenir}

\begin{attention}[Piège à éviter : indépendance $\neq$ souveraineté totale]
L'erreur fréquente consiste à confondre \emph{indépendance politique} (siège à l'ONU, drapeau, armée) et \emph{souveraineté effective} (capacité à décider seul de sa politique économique, monétaire, diplomatique). Le néocolonialisme ne nie pas l'indépendance juridique ; il en vide la substance. Au Bac, une dissertation sur ce thème doit impérativement distinguer le \textbf{formel} (juridique) du \textbf{réel} (économique et stratégique).
\end{attention}

\courssub{Les instruments économiques : monnaie, dette et multinationales}

Le levier le plus puissant reste l'économie. Dès les indépendances, les accords de coopération ont verrouillé les structures extractives : l'Afrique exporte des matières premières brutes (cacao, café, coton, pétrole, minerais) et importe des produits manufacturés, dans un échange inégal dont les termes de l'échange se dégradent structurellement (théorie de Prebisch-Singer).

\begin{definition}[Franc CFA]
Le \textbf{franc CFA} (initialement \emph{Franc des Colonies Françaises d'Afrique}, devenu \emph{Franc de la Communauté Financière Africaine} pour la zone UEMOA et \emph{Franc de la Coopération Financière en Afrique Centrale} pour la zone CEMAC) est une monnaie commune à 14 pays africains (dont le Cameroun), arrimée à l'euro (anciennement au franc français) à parité fixe (1 € = 655,957 FCFA). Sa convertibilité est \textbf{garantie par le Trésor français}, contre dépôt d'une partie des réserves de change (historiquement 50 \%, ramené à 20 \% en 2019 pour l'UEMOA ; la réforme pour la zone CEMAC est en cours depuis 2023) sur un compte d'opérations auprès de la Banque de France.
\end{definition}

\begin{exemplebox}[Exemple chiffré : le franc CFA au Cameroun]
Le Cameroun, membre de la zone CEMAC (BEAC), utilise le franc CFA depuis 1945. En 2023, ses réserves de change logées au Trésor français représentent une garantie de convertibilité, mais limitent la marge de manœuvre de la politique monétaire nationale (taux d'intérêt alignés sur la BCE, impossibilité de dévaluation compétitive unilatérale). Le débat sur la sortie du franc CFA (projet \emph{Eco} en zone UEMOA) illustre la tension entre stabilité monétaire et souveraineté économique.
\end{exemplebox}

\begin{propriete}[Les sociétés transnationales : acteurs majeurs du néocolonialisme]
Des groupes comme \textbf{TotalEnergies} (pétrole, gaz), \textbf{Orange} (télécoms) ou \textbf{Olam} (agro-industrie) contrôlent des segments entiers d'économies africaines. Au Cameroun, \textbf{Africa Global Logistics (AGL, ex-Bolloré Africa Logistics, racheté par MSC en 2022)} gère le port autonome de Douala (concession 2008-2028) et la filiale Camrail (chemin de fer) ; TotalEnergies domine la distribution pétrolière ; Orange est leader de la téléphonie mobile. Ces positions dominantes s'accompagnent souvent d'optimisation fiscale et de rapatriement massif des bénéfices.
\end{propriete}

\begin{popfigure}
\begin{tikzpicture}
\begin{axis}[
    width=\linewidth,
    height=9cm,
    ybar,
    bar width=18pt,
    enlarge x limits=0.25,
    ylabel={Part de marché estimée (\%)\strut},
    symbolic x coords={Ports/Logistique (AGL), Pétrole aval (TotalEnergies), Télécoms (Orange), Cacao/Export (Olam/Cargill), Coton (Sodecoton/Geocoton)},
    xtick=data,
    x tick label style={rotate=30, anchor=east, font=\small},
    ymin=0, ymax=100,
    ymajorgrids=true,
    grid style=dashed,
    legend style={at={(0.5,-0.35)}, anchor=north, legend columns=-1, font=\small},
    nodes near coords={\pgfmathprintnumber\pgfplotspointmeta\,\%},
    nodes near coords align={vertical},
    every node near coord/.append style={font=\footnotesize, rotate=90, anchor=west, yshift=2pt}
]
\addplot[fill=popblue!70] coordinates {
    (Ports/Logistique (AGL), 85)
    (Pétrole aval (TotalEnergies), 65)
    (Télécoms (Orange), 55)
    (Cacao/Export (Olam/Cargill), 70)
    (Coton (Sodecoton/Geocoton), 60)
};
\addplot[fill=poporange!70] coordinates {
    (Ports/Logistique (AGL), 15)
    (Pétrole aval (TotalEnergies), 35)
    (Télécoms (Orange), 45)
    (Cacao/Export (Olam/Cargill), 30)
    (Coton (Sodecoton/Geocoton), 40)
};
\legend{Entreprises étrangères (groupe dominant), Acteurs locaux / autres}
\end{axis}
\end{tikzpicture}
\legende{Part de marché estimée de groupes étrangers dominants dans cinq secteurs stratégiques d'Afrique de l'Ouest et centrale (données indicatives, sources : rapports annuels des groupes, Banque mondiale, 2020-2023). La domination dépasse souvent 50 \% dans les infrastructures portuaires, l'énergie et l'export agricole.}
\end{popfigure}

\courssub{Les instruments politiques, militaires et la « Françafrique »}

La domination économique s'appuie sur un dispositif politico-militaire dense, hérité des accords de coopération signés au moment des indépendances (1960-1961).

\begin{definition}[Françafrique]
La \textbf{Françafrique} désigne le réseau opaque de relations personnelles, financières et militaires entre l'Élysée, le Quai d'Orsay, les services secrets (DGSE, DST historique) et les présidents africains francophones, mis en place dès 1960 par \textbf{Jacques Foccart}, secrétaire général à l'Élysée pour les affaires africaines et malgaches. Ce système assure le soutien à des régimes « amis » (stabilité, alignement diplomatique, accès aux ressources) en échange de loyauté politique et de contrats pour les entreprises françaises.
\end{definition}

\begin{aretenir}[Dispositif militaire et interventions (historique et actualité)]
\begin{itemize}
    \item \textbf{Accords de défense :} signés avec une dizaine de pays (Cameroun, Gabon, Côte d'Ivoire, Sénégal, Tchad, etc.), ils prévoient l'assistance technique, la formation, l'intervention militaire en cas d'agression extérieure \emph{ou} de menace à l'ordre constitutionnel (clause souvent invoquée pour justifier des interventions).
    \item \textbf{Bases militaires historiques :} Dakar (Sénégal), Abidjan (Côte d'Ivoire), Libreville (Gabon), N'Djamena (Tchad), Djibouti. Prépositionnement de forces (opérations \emph{Barkhane}, \emph{Sangaris}, \emph{Licorne}).
    \item \textbf{Interventions notables :} Gabon (1964, rétablissement de M'ba), Tchad (1968-1972, 1983-1984 Opération Manta, 1986 Opération Épervier), Côte d'Ivoire (2002-2011 Opération Licorne), Centrafrique (2013-2016 Sangaris), Sahel (2014-2022 Barkhane).
    \item \textbf{Ruptures récentes (2022-2025) :} Retrait du Mali (août 2022), du Burkina Faso (janvier 2023), du Niger (décembre 2023), annonce du retrait du Tchad (janvier 2025) et de la Côte d'Ivoire (janvier 2025), fermeture de la base de Dakar (2024). Redéploiement vers le Tchad, Djibouti, Gabon.
\end{itemize}
\end{aretenir}

\begin{savaistu}[L'assassinat de Patrice Lumumba (1961) : un crime d'État néocolonial]
Premier ministre du Congo indépendant (juin 1960), \textbf{Patrice Lumumba} refuse la tutelle belge et occidentale, appelle à l'unité africaine et sollicite l'URSS. Il est renversé par un coup d'État du colonel Mobutu (soutenu par la CIA et les services belges), transféré au Katanga sécessionniste et exécuté le 17 janvier 1961 sous supervision d'officiers belges. En 2002, la Belgique a reconnu sa « responsabilité morale » ; en 2022, une dent de Lumumba, seule relique conservée par un policier belge, a été restituée à la RDC. Ce crime illustre l'élimination physique des leaders nationalistes refusant le néocolonialisme.
\end{savaistu}

\courssub{Le cas camerounais : la liquidation de l'UPC et les accords de coopération}

Le Cameroun offre un cas d'école de la transition « gérée » vers un néocolonialisme structurel. L'indépendance (1er janvier 1960) est arrachée dans la violence : l'Union des Populations du Cameroun (UPC), parti nationaliste radical fondé en 1948, mène une guérilla (maquis) dès 1955. La France refuse de négocier avec l'UPC (interdite en 1955) et soutient Ahmadou Ahidjo, perçu comme modéré et fiable.

\begin{experience}[Chronologie de la neutralisation de l'UPC (1958-1971)]
\begin{itemize}
    \item \textbf{13 septembre 1958 :} Assassinat de \textbf{Ruben Um Nyobè}, secrétaire général de l'UPC, par l'armée française près de Boumnyébel (Sanaga-Maritime). La France présente sa mort comme un « accident » de maquis.
    \item \textbf{3 novembre 1960 :} Empoisonnement au \textbf{thallium} de \textbf{Félix-Roland Moumié}, successeur d'Um Nyobè à la tête de l'UPC, à Genève (Suisse), par \textbf{William Bechtel}, agent du SDECE (services secrets français) infiltré. Bechtel sera identifié, arrêté, puis relâché sous pression diplomatique.
    \item \textbf{1964 :} Mort suspecte d'\textbf{Abel Kingué} au Caire.
    \item \textbf{1966 :} Mort d'\textbf{Osende Afana}, économiste brillant, cadre de l'UPC, tué au maquis. Instauration du parti unique (UNC) par Ahidjo.
    \item \textbf{Janvier 1971 :} Exécution publique d'\textbf{Ernest Ouandié} (dernier grand chef maquisard) à Bafoussam.
    \item \textbf{1971 :} Reddition du dernier maquisard, \textbf{Wambo le Roi}. L'UPC est décapitée.
\end{itemize}
\end{experience}

\begin{savaistu}[Félix Moumié : l'empoisonnement au thallium à Genève]
Le 3 novembre 1960, \textbf{Félix Moumié}, président de l'UPC en exil, déjeune à l'hôtel Platini à Genève avec William Bechtel, un journaliste suisse en réalité agent du SDECE. Bechtel verse du thallium (poison lent, indétectable à l'époque) dans son verre. Moumié meurt le 15 novembre à l'hôpital cantonal. L'autopsie révèle le thallium. Bechtel est arrêté, avoue sous contrainte, puis bénéficie d'un non-lieu suisse « pour raison d'État » et quitte la Suisse sous protection française. Ce crime d'État n'a jamais été jugé.
\end{savaistu}

\begin{aretenir}[Accords de coopération Cameroun-France (1960-1961)]
Signés dès janvier 1960, ils couvrent : \textbf{défense} (intervention française possible), \textbf{monnaie} (maintien du franc CFA, compte d'opérations au Trésor), \textbf{affaires étrangères} (consultation préalable), \textbf{justice} (magistrats français détachés), \textbf{éducation} (programmes, bourses), \textbf{technique} (coopérants). Ces accords, renouvelés et complétés (1974, 1994), structurent encore aujourd'hui la relation bilatérale.
\end{aretenir}

\courssub{Les réponses africaines : de l'OUA à la ZLECAF}

Face à la dépendance, les États africains ont tenté des réponses collectives et nationales.

\begin{aretenir}[Principales réponses historiques]
\begin{itemize}
    \item \textbf{OUA (25 mai 1963, Addis-Abeba) :} Charte affirmant l'intangibilité des frontières héritées (principe \emph{uti possidetis juris}) pour éviter les guerres, soutien aux mouvements de libération (Angola, Mozambique, Afrique du Sud), non-ingérence dans les affaires intérieures (frein à la sanction des dictatures).
    \item \textbf{Mouvement des Non-Alignés (Belgrade 1961) :} Refus des blocs Est/Ouest, revendication d'un \textbf{Nouvel Ordre Économique Mondial (NOEM)} aux Nations Unies (années 1970) : souveraineté sur ressources, régulation des transnationales, transfert de technologie, annulation de la dette.
    \item \textbf{Nationalisations :} Algérie (pétrole 1971), Zaïre (Zaïrianisation 1973-1974), Tanzanie (villagisation Ujamaa), Ghana (cacao). Résultats mitigés : manque de cadres, chute de production, endettement accru.
    \item \textbf{Plans de redressement structurel (FMI/Banque mondiale, années 1980-1990) :} Imposés comme condition d'aide, ils ont démantelé les services publics, privatisé à bas prix, aggravé la pauvreté — souvent perçus comme instruments néocoloniaux.
\end{itemize}
\end{aretenir}

\begin{savaistu}[Complément utile au Bac : la ZLECAF (2019/2021)]
La \textbf{Zone de Libre-Échange Continentale Africaine (ZLECAF / AfCFTA)}, entrée en vigueur le 30 mai 2019 (signature 2018 à Kigali), démarrage effectif des échanges le 1er janvier 2021, vise à créer un marché unique de 1,3 milliard de consommateurs, PIB cumulé $\approx$ 3 400 milliards \$. Objectifs : supprimer les droits de douane sur 90 \% des lignes tarifaires, harmoniser les règles d'origine, faciliter la circulation des personnes et des capitaux. C'est la réponse panafricaine la plus ambitieuse au néocolonialisme économique : \emph{produire et échanger en Afrique pour l'Afrique}. Le Cameroun a ratifié en 2020 ; le secrétariat est à Accra (Ghana). Au Bac, citer la ZLECAF comme « perspective d'émancipation économique » valorise la conclusion.
\end{savaistu}

\courssub{Bilan nuancé et enjeux contemporains}

\begin{propriete}[Bilan : indépendances politiques réelles, souveraineté économique inachevée]
\begin{itemize}
    \item \textbf{Acquis indéniables :} États membres de l'ONU, constitutions, armées, élites formées, classe moyenne émergente, démarches démocratiques (années 1990), croissance soutenue depuis 2000 (hors chocs Covid/guerre Ukraine).
    \item \textbf{Dépendances persistantes :} Monnaie (CFA), dette extérieure (service de la dette $\approx$ 15-25 \% des recettes budgétaires), balance commerciale structurellement déficitaire (export brut / import manufacturé), armées formées/équipées par l'ancien colonisateur, élites francophones socialisées en France.
    \item \textbf{Nouveaux acteurs :} Chine (infrastructures contre ressources, dette opaque), Russie (Wagner/Africa Corps, appui militaire aux juntes au Sahel), Turquie, Émirats, Inde. Le néocolonialisme n'est plus l'apanage exclusif de la France.
\end{itemize}
\end{propriete}

\begin{methode}[Rédiger une conclusion de dissertation sur le néocolonialisme]
\textbf{Étape 1 : Rappel de la problématique.} Reprenez la question posée (ex. : « L'indépendance africaine est-elle une réalité ou une illusion ? »).
\textbf{Étape 2 : Bilan synthétique de l'argumentation.} Deux phrases : « Sur le plan politique, les indépendances des années 1960 sont une réalité juridique et diplomatique incontestable (sièges à l'ONU, frontières reconnues). Sur le plan économique et stratégique, la souveraineté reste tronquée : monnaie arrimée, dette structurante, mainmise des multinationales, accords de défense asymétriques. »
\textbf{Étape 3 : Ouverture pertinente.} Évitez les lieux communs. Privilégiez : « L'émergence de la ZLECAF (2021), la diversification des partenariats (Chine, Turquie, Brésil) et les mobilisations citoyennes pour la sortie du franc CFA (mouvements \emph{France Dégage}, \emph{Y'en a marre}) dessinent les contours d'une \emph{deuxième indépendance}, cette fois économique, dont l'aboutissement reste le grand défi du XXIe siècle africain. »
\end{methode}

\begin{popfigure}
\begin{tikzpicture}[
    node distance=2.2cm and 3cm,
    box/.style={rectangle, rounded corners, draw=popdark, thick, fill=popblueL, text width=2.8cm, align=center, font=\small, minimum height=1.6cm},
    centerbox/.style={rectangle, rounded corners, draw=poporange, very thick, fill=poporangeL, text width=3.2cm, align=center, font=\small\bfseries, minimum height=2cm},
    arrow/.style={-Stealth, thick, popdark}
]
\node[centerbox, align=center] (center){NÉOCOLONIALISME\\MÉCANISMES};
\node[box, align=center] (eco) [above left=of center]{ÉCONOMIQUE\\Franc CFA, dette,\\multinationales,\\matières premières};
\node[box, align=center] (mon) [left=of center]{MONÉTAIRE\\Parité fixe euro,\\réserves au Trésor,\\pas de banque centrale\\souveraine};
\node[box, align=center] (mil) [below left=of center]{MILITAIRE\\Bases, accords\\de défense,\\interventions,\\formation officiers};
\node[box, align=center] (pol) [below right=of center]{POLITIQUE\\Françafrique,\\soutien régimes\\amis, ingérences\\électorales};
\node[box, align=center] (cul) [above right=of center]{CULTUREL\\Langue, éducation,\\médias, élites\\formées à l'étranger};
\foreach \n in {eco,mon,mil,pol,cul}
    \draw[arrow] (center) -- (\n);
\draw[arrow, poporange, dashed] (eco) -- node[midway, above, sloped, font=\tiny, poporange] {financement} (mon);
\draw[arrow, poporange, dashed] (mil) -- node[midway, below, sloped, font=\tiny, poporange] {garantie} (pol);
\draw[arrow, poporange, dashed] (cul) -- node[midway, above, sloped, font=\tiny, poporange] {légitimation} (pol);
\end{tikzpicture}
\legende{Schéma en étoile des cinq mécanismes interconnectés du néocolonialisme en Afrique francophone. Le centre (néocolonialisme) irrigue et relie les cinq piliers ; les flèches orange pointillées montrent les rétroactions (ex. : le militaire garantit le politique, le culturel légitime le politique, l'économique finance le monétaire).}
\end{popfigure}

\begin{exoresolu}[Commentaire de document : extrait de Kwame Nkrumah, \emph{Néocolonialisme, dernier stade de l'impérialisme}, 1965]
\textbf{Document :} « Le néocolonialisme est la forme la plus aboutie de l'impérialisme. Il est l'impérialisme sans responsabilité. Ceux qui le pratiquent ont le pouvoir sans la responsabilité ; ceux qui le subissent ont la responsabilité sans le pouvoir. [...] Le résultat du néocolonialisme est que le capital étranger est utilisé pour l'exploitation plutôt que pour le développement des parties moins développées du monde. L'investissement, sous le néocolonialisme, augmente l'écart entre les pays riches et les pays pauvres. »

\textbf{Question :} En vous appuyant sur le document et vos connaissances, montrez comment Nkrumah définit le néocolonialisme et illustrez par deux exemples précis (monétaire et militaire) en Afrique francophone.

\tcblower

\textbf{Solution détaillée :}

\textbf{1. Analyse du document (intuition + rigueur).}
\begin{itemize}
    \item \textbf{Auteur/contexte :} Kwame Nkrumah, premier président du Ghana (indépendance 1957), figure du panafricanisme, renversé par un coup d'État soutenu par la CIA en 1966. Ouvrage publié en 1965, au cœur de la Guerre froide et des premières indépendances.
    \item \textbf{Thèse centrale :} Le néocolonialisme = « impérialisme sans responsabilité ». Deux formules clés : « pouvoir sans responsabilité » (pour l'ex-puissance) / « responsabilité sans pouvoir » (pour l'État africain) ; « capital étranger utilisé pour l'exploitation, pas le développement » ; « l'investissement creuse l'écart riches/pauvres ».
    \item \textbf{Portée :} Définition structurelle, systémique. Ce n'est pas une anecdote, c'est un \emph{système}.
\end{itemize}

\textbf{2. Illustration monétaire : le franc CFA (zone CEMAC/UEMOA).}
\begin{itemize}
    \item \textbf{Mécanisme :} Parité fixe garantie par le Trésor français contre dépôt de réserves (compte d'opérations). Taux d'intérêt alignés sur la BCE. Impossible de dévaluer pour relancer l'export ou de financer le déficit budgétaire par la planche à billets.
    \item \textbf{Exemple camerounais :} En 2023, le Cameroun dépose une partie de ses réserves à la Banque de France. La BEAC (Banque des États de l'Afrique Centrale) n'a pas de mandat de plein emploi, seulement la stabilité des prix. Résultat : politique monétaire \emph{procyclique} (hausse des taux en 2022-2023 pour suivre la BCE alors que l'inflation est importée), crédit à l'économie étouffé. « Pouvoir sans responsabilité » : la France garantit la convertibilité mais n'assume pas les conséquences sociales de la rigueur monétaire. « Responsabilité sans pouvoir » : Yaoundé gère le chômage et la pauvreté sans levier monétaire.
\end{itemize}

\textbf{3. Illustration militaire : l'opération \emph{Barkhane} au Sahel (2014-2022) et les accords de défense.}
\begin{itemize}
    \item \textbf{Cadre :} Accords de défense signés dans les années 1960 (Niger, Mali, Tchad, Burkina Faso, Mauritanie). France maintient 5 100 hommes au pic (2020), bases à N'Djamena, Gao, Niamey.
    \item \textbf{Dynamique néocoloniale :} Intervention décidée à Paris (Élysée), non par les parlements africains. Objectifs flous (contre-terrorisme, mais aussi protection des intérêts Areva/Orano uranium au Niger, Total au Mali). Soutien à des régimes autoritaires (Déby au Tchad, Bazoum au Niger) au nom de la « stabilité ».
    \item \textbf{Rupture 2022-2024 :} Junte malienne (2020) exige le départ de Barkhane (août 2022) ; junte burkinabè (2022) dénonce l'accord de défense (janvier 2023) ; junte nigérienne (juillet 2023) exige le retrait français (décembre 2023). La France retire ses troupes, mais conserve des capacités de projection (Tchad, Djibouti, Gabon). Illustration parfaite : « pouvoir sans responsabilité » (France décide intervention/retrait) ; « responsabilité sans pouvoir » (États sahéliens subissent insécurité, flux réfugiés, sans maîtrise des opérations).
\end{itemize}

\textbf{4. Conclusion synthétique.}
Nkrumah avait vu juste : le néocolonialisme est un \emph{système} où la monnaie (CFA) et l'armée (accords de défense/Barkhane) sont les deux piliers techniques de la « responsabilité sans pouvoir ». Les ruptures actuelles au Sahel (sortie du franc CFA annoncée par l'AES — Alliance des États du Sahel — en 2024, diversification vers la Russie) tentent de briser ce système, mais le risque est de substituer une dépendance (française) à une autre (russe/chinoise). La vraie issue, comme l'esquisse la ZLECAF, est l'intégration africaine par le bas.
\end{exoresolu}

\begin{exercice}[Entraînement Bac : Sujet de dissertation]
\textbf{Sujet :} « L'indépendance politique de l'Afrique (années 1960) a-t-elle été confisquée par le néocolonialisme ? »
\textit{Consignes :} Définir les termes, problématiser, structurer en trois parties (I. Une indépendance formelle acquise ; II. Une souveraineté économique et stratégique confisquée par des mécanismes néocoloniaux ; III. Des réponses africaines inabouties mais renouvelées aujourd'hui). Utiliser les exemples du Ghana, de l'Algérie, de l'Angola, du Cameroun et du Sahel actuel. Intégrer la ZLECAF en ouverture.
\end{exercice}

\begin{corrige}[Exercice n°1 : Éléments de corrigé]
\textbf{Introduction :} Accroche (citation Nkrumah ou Senghor « Nous avons l'indépendance, nous n'avons pas la liberté »). Définitions : indépendance politique (souveraineté juridique) vs néocolonialisme (domination réelle). Problématique : l'indépendance a-t-elle été \emph{vidée} de sa substance ou \emph{reconfigurée} ? Annonce du plan.
\textbf{I. Une indépendance formelle acquise (1957-1975).} Ghana (1957, négociation), Algérie (1962, guerre), Angola (1975, guerre), Cameroun (1960, violence/UPC). Sièges ONU, frontières, constitutions, armées.
\textbf{II. Une souveraineté confisquée : les mécanismes néocoloniaux.} Monétaire (Franc CFA, 14 pays, Trésor français), Économique (multinationales AGL/TotalEnergies/Orange, matières premières bradées, dette), Militaire (bases, accords défense, Barkhane, Françafrique/Foccart), Politique (coups d'État téléguidés, assassinats Lumumba/Moumié/Um Nyobè), Culturel (langue, élites).
\textbf{III. Réponses africaines : de l'OUA à la ZLECAF.} OUA 1963 (frontières, non-ingérence), Non-alignés/NOEM (échec années 1980), Nationalisations (mitigé), PAS/FMI (aggravation), ZLECAF 2021 (marché unique, espoir), Ruptures Sahel 2023-2024 (sortie CFA, départ France, pivot Russie/Chine — nouvelle dépendance ?).
\textbf{Conclusion :} Bilan : indépendance réelle mais tronquée. Ouverture : la « deuxième indépendance » économique (ZLECAF, industrialisation, maîtrise monétaire) est le chantier du siècle. Le néocolonialisme n'est pas une fatalité, mais sa défaite exige l'unité africaine que Nkrumah appelait de ses vœux.
\end{corrige}

\newpage

\coursec{Activité d'intégration}

\courssub{Objectifs de l'activité}

À l'issue de cette activité d'intégration, l'élève doit être capable de :
\begin{itemize}
\item construire une carte historique conforme aux règles cartographiques (titre, légende, orientation, échelle indicative) à partir d'un tableau de données ;
\item mobiliser les savoirs de la leçon (cas du Ghana, de l'Algérie, de la Guinée Conakry et de l'Angola) pour caractériser les différentes voies d'accès à l'indépendance ;
\item analyser des documents historiques variés (discours, données économiques, article de presse, photographie) et en extraire des arguments ;
\item rédiger et justifier une conclusion argumentée répondant à une problématique : \emph{« L'indépendance du Cameroun est-elle achevée ? »} ;
\item appliquer les notions de décolonisation et de néocolonialisme à une situation concrète de la vie nationale.
\end{itemize}

\courssub{Situation}

Dans le cadre de la semaine de la jeunesse, le club d'Histoire du Lycée technique de Nkongsamba prépare une exposition intitulée \emph{« L'Afrique : des indépendances à la souveraineté réelle »}. Le proviseur a demandé aux élèves de la classe de Terminale technique de réaliser deux productions : une grande carte murale de la décolonisation africaine et un texte de synthèse qui sera lu lors du vernissage devant les parents d'élèves et les autorités locales.

Pour construire la carte, le club dispose du tableau de données suivant :

\begin{center}
\begin{tabularx}{\linewidth}{|l|X|X|X|}\hline
\rowcolor{popblueL} \textbf{Pays} & \textbf{Puissance coloniale} & \textbf{Date d'indépendance} & \textbf{Voie d'accès} \\ \hline
Ghana (ex-Gold Coast) & Grande-Bretagne & 6 mars 1957 & Négociée (lutte politique non violente de Nkrumah) \\ \hline
Guinée Conakry & France & 2 octobre 1958 & Référendaire (« Non » au référendum du 28 septembre 1958) \\ \hline
Cameroun & France et Grande-Bretagne & 1\textsuperscript{er} janvier 1960 & Négociée (après répression de l'UPC) \\ \hline
Algérie & France & 5 juillet 1962 & Armée (guerre de 1954 à 1962, accords d'Évian) \\ \hline
Angola & Portugal & 11 novembre 1975 & Armée (guerre de libération de 1961 à 1975) \\ \hline
\end{tabularx}
\end{center}

Pour le texte de synthèse, le club dispose d'un dossier documentaire de quatre documents :
\begin{itemize}
\item \textbf{Document 1} : extrait de la réponse de Sékou Touré au général de Gaulle, Conakry, 25 août 1958 : « Nous préférons la liberté dans la pauvreté à la richesse dans l'esclavage. »
\item \textbf{Document 2} : données économiques : le franc CFA, utilisé au Cameroun, est arrimé à l'euro avec une parité fixe garantie par le Trésor français ; une partie des réserves de change de la zone franc est déposée auprès du Trésor français.
\item \textbf{Document 3} : article de presse : « Au Cameroun, de grandes entreprises françaises (Bolloré, Total, Orange, Société générale) occupent des positions dominantes dans les ports, l'énergie, les télécommunications et la banque, soixante ans après l'indépendance. »
\item \textbf{Document 4} : photographie de la base militaire française de Djibouti, avec la légende : « La France conserve plusieurs bases militaires en Afrique et intervient régulièrement sur le continent. »
\end{itemize}

La classe est divisée en quatre groupes. Chaque groupe construit sa carte, analyse le dossier, puis participe au débat contradictoire final : deux camps s'opposent, l'un soutenant que l'indépendance du Cameroun est achevée, l'autre qu'elle demeure inachevée. Le travail se termine par une rédaction individuelle.

\courssub{Consignes}

\begin{enumerate}
\item À partir du tableau de données, construis par groupe une carte intitulée « Les voies de la décolonisation en Afrique ». Respecte les règles cartographiques : titre, légende en trois couleurs (une couleur par voie d'accès : négociée, référendaire, armée), flèche du nord et échelle indicative.
\item Pour chacun des quatre pays étudiés dans la leçon (Ghana, Algérie, Guinée, Angola), indique sur la carte la date d'indépendance et la puissance coloniale, puis justifie en une phrase la couleur choisie.
\item Compare les voies d'accès à l'indépendance : quels pays ont obtenu leur indépendance sans guerre ? Lesquels ont dû recourir à la lutte armée ? Propose une explication à cette différence.
\item Analyse le document 1 : que révèle le choix de Sékou Touré sur la volonté d'émancipation des peuples africains ? Quelles en ont été les conséquences immédiates pour la Guinée ?
\item Analyse les documents 2, 3 et 4 : quels domaines de la vie du Cameroun (monnaie, économie, sécurité) restent marqués par la présence de l'ancienne puissance coloniale ? À quel concept étudié dans la leçon cela renvoie-t-il ?
\item Participe au débat contradictoire en défendant la position de ton camp à l'aide d'au moins deux arguments tirés des documents et de la leçon.
\item Rédige individuellement une conclusion argumentée (15 à 20 lignes) répondant à la question : « L'indépendance du Cameroun est-elle achevée ? » Ta conclusion doit être justifiée par des faits précis et datés.
\end{enumerate}

\courssub{Ressources à mobiliser}

\begin{aretenir}[Notions et repères indispensables]
\begin{itemize}
\item \cle{Décolonisation} : processus par lequel les territoires colonisés accèdent à la souveraineté internationale. Trois voies : \cle{négociée} (Ghana, 1957), \cle{référendaire} (Guinée, « Non » du 28 septembre 1958, indépendance le 2 octobre 1958), \cle{armée} (Algérie, guerre 1954-1962 ; Angola, guerre 1961-1975).
\item Acteurs clés : Kwame \cle{Nkrumah} (Ghana), le \cle{FLN} et Ahmed Ben Bella (Algérie), Sékou \cle{Touré} (Guinée), le \cle{MPLA} et Agostinho Neto (Angola).
\item \cle{Néocolonialisme} : domination économique, monétaire, culturelle et militaire exercée par les anciennes puissances coloniales (ou de nouvelles puissances) sur des États pourtant politiquement indépendants. Manifestations : franc CFA, accords de coopération, entreprises étrangères, bases militaires.
\item Règles cartographiques : tout document cartographique comporte un titre, une légende, une orientation (flèche du nord) et une échelle indicative.
\end{itemize}
\end{aretenir}

\courssub{Démarche et corrigé commenté}

\begin{exoresolu}[Corrigé commenté de l'activité]
\textbf{Étape 1 : construction de la carte (consignes 1 et 2).}
Le groupe classe d'abord les pays selon la voie d'accès, ce qui donne la légende à trois couleurs :
\begin{itemize}
\item Voie négociée (vert) : Ghana (Grande-Bretagne, 6 mars 1957), Cameroun (France et Grande-Bretagne, 1\textsuperscript{er} janvier 1960) ;
\item Voie référendaire (orange) : Guinée Conakry (France, 2 octobre 1958) ;
\item Voie armée (rouge) : Algérie (France, 5 juillet 1962), Angola (Portugal, 11 novembre 1975).
\end{itemize}
Justifications : le Ghana est coloré en vert car Nkrumah a obtenu l'indépendance par la lutte politique non violente (boycotts, grèves, victoire électorale du CPP) et la négociation avec Londres. La Guinée est en orange car son indépendance résulte du « Non » massif au référendum du 28 septembre 1958 sur la Communauté française. L'Algérie et l'Angola sont en rouge car elles ont conquis leur indépendance par de longues guerres de libération.

\begin{popfigure}
\begin{center}
\begin{tikzpicture}[scale=0.9, font=\small]
\draw[thick, popdark, fill=popcream] (0,5) -- (2,6.2) -- (4.5,6.5) -- (6,5.5) -- (6.5,4) -- (5.5,2.5) -- (4.5,0.5) -- (3.5,-0.5) -- (2.5,0.5) -- (1.5,2) -- (0,3) -- cycle;
\fill[popgreen!60] (1.2,3.4) circle (0.35) node[right=2pt, popdark] {Ghana 1957};
\fill[poporange!80] (0.6,3.9) circle (0.25) node[left=2pt, popdark] {Guinée 1958};
\fill[poppink!70] (3.4,4.8) circle (0.5) node[above=2pt, popdark] {Algérie 1962};
\fill[poppink!70] (3.6,1.2) circle (0.35) node[right=2pt, popdark] {Angola 1975};
\fill[popgreen!60] (3.2,2.9) circle (0.3) node[right=2pt, popdark] {Cameroun 1960};
\draw[->, thick, popdark] (6.8,5.5) -- (6.8,6.3) node[above] {N};
\draw[thick, popdark] (0.2,-0.8) -- (1.7,-0.8) node[midway, below] {\footnotesize échelle indicative};
\end{tikzpicture}
\end{center}
\legende{Croquis simplifié de la carte attendue : vert = voie négociée, orange = voie référendaire, rouge = voie armée. Le contour de l'Afrique est volontairement très schématique.}
\end{popfigure}

\tcblower

\textbf{Étape 2 : comparaison des voies (consigne 3).}
Le Ghana et le Cameroun ont obtenu leur indépendance sans guerre de libération, par la négociation ; la Guinée par un simple vote. En revanche, l'Algérie et l'Angola ont dû mener de longues luttes armées. Explication : là où la présence coloniale était forte et ancienne (Algérie, départements français peuplés d'environ un million de colons européens ; Angola, colonie de peuplement portugaise dirigée par la dictature de Salazar), la métropole a refusé tout départ négocié, ce qui a rendu la guerre inévitable.

\textbf{Étape 3 : analyse du document 1 (consigne 4).}
La phrase de Sékou Touré montre que les peuples africains placent la liberté politique au-dessus des avantages matériels offerts par la colonisation. Conséquences immédiates : la Guinée proclame son indépendance le 2 octobre 1958 ; la France se retire brutalement, emporte ses équipements, rappelle ses cadres et rompt toute aide, tentant ainsi de sanctionner l'exemple guinéen.

\textbf{Étape 4 : analyse des documents 2 à 4 (consigne 5).}
Le document 2 montre que la monnaie du Cameroun reste liée au Trésor français (domaine monétaire) ; le document 3 montre que des entreprises françaises dominent des secteurs stratégiques de l'économie (domaine économique) ; le document 4 montre que la France conserve une présence militaire en Afrique (domaine sécuritaire). Ces trois faits renvoient au concept de \cle{néocolonialisme} : la souveraineté politique est formelle, mais la dépendance économique, monétaire et militaire persiste.

\textbf{Étape 5 : débat (consigne 6).}
Camp A (indépendance achevée) : le Cameroun possède un drapeau, un hymne, un siège à l'ONU depuis 1960, choisit ses dirigeants et signe ses propres accords internationaux ; le « Non » guinéen de 1958 prouve qu'un peuple peut rompre totalement s'il le veut. Camp B (indépendance inachevée) : le franc CFA arrimé au Trésor français, la domination des entreprises étrangères et les accords de coopération militaire limitent concrètement les choix nationaux, comme le montre le dossier.

\textbf{Étape 6 : modèle de conclusion individuelle (consigne 7).}
« L'indépendance du Cameroun, proclamée le 1\textsuperscript{er} janvier 1960, est réelle sur le plan politique : l'État dispose de ses institutions, de sa diplomatie et de son armée. Cependant, l'analyse du dossier montre que cette souveraineté reste inachevée sur les plans économique et monétaire : le franc CFA demeure garanti par le Trésor français et de grandes entreprises étrangères contrôlent des secteurs clés comme les ports et les télécommunications. Comme l'Algérie ou l'Angola, qui ont payé leur liberté au prix de la guerre, le Cameroun est politiquement libre ; mais à l'image de la Guinée de Sékou Touré, qui a voulu une rupture totale en 1958, la pleine souveraineté suppose aussi l'émancipation économique. L'indépendance du Cameroun est donc acquise juridiquement, mais demeure à parachever économiquement. »

\textbf{Commentaire méthodologique :} la bonne conclusion ne tranche pas de façon absolue ; elle distingue les plans (politique, économique, monétaire, militaire), cite des faits datés et mobilise les quatre cas de la leçon. C'est exactement ce qu'attend l'épreuve d'Histoire au Probatoire et au Baccalauréat technique.
\end{exoresolu}

\courssub{Bilan de l'activité}

Cette activité a d'abord permis de transformer un tableau de données en carte historique : les élèves ont vérifié concrètement que la décolonisation africaine n'a pas suivi un modèle unique, mais trois voies distinctes — négociée (Ghana, 1957), référendaire (Guinée, 1958) et armée (Algérie, 1954-1962 ; Angola, 1961-1975) — et que la radicalité de la lutte dépend de la nature de la colonisation (colonie d'exploitation ou de peuplement). Ensuite, l'analyse croisée du dossier documentaire a rendu tangible la notion de néocolonialisme : derrière les indépendances politiques proclamées autour de 1960 subsistent des dépendances monétaires (franc CFA), économiques (entreprises étrangères) et militaires (bases, accords de coopération). Enfin, le débat contradictoire et la rédaction finale ont entraîné les élèves à construire une argumentation nuancée et justifiée, compétence centrale du savoir-faire « rédiger et justifier une démarche, une réponse ou une conclusion » exigée au Probatoire et au Baccalauréat technique.

\newpage

\coursec{Méthodes et exercices résolus}

\courssub{Méthodes et exercices résolus}

\begin{methode}[Analyser un document historique : commentaire composé type Bac]
\textbf{Objectif :} Restituer le sens d'un document (texte, discours, loi, témoignage) en le resituant dans son contexte et en en dégagent la portée historique.
\par \textbf{Démarche en 4 étapes :}
\begin{enumerate}
    \item \textbf{Identification (Nature, Auteur, Date, Source, Contexte immédiat) :} Préciser la nature du document (discours, article de presse, traité, lettre), son auteur (statut, rôle), la date précise (contexte chronologique), la source et le contexte immédiat (événement déclencheur).
    \item \textbf{Analyse du contenu (Vocabulaire, Idées forces, Structure) :} Définir les termes clés (\cle{vocabulaire spécifique}). Repérer les idées principales (souvent 2 ou 3) et la logique du raisonnement de l'auteur (argumentation, justification, revendication).
    \item \textbf{Mise en perspective (Contexte général, Antécédents, Conséquences) :} Mobiliser les connaissances du cours pour expliquer \emph{pourquoi} ce document est produit à ce moment-là (causes profondes) et \emph{ce qu'il change} (conséquences immédiates et à long terme). C'est l'apport du savoir personnel.
    \item \textbf{Synthèse / Portée historique (Conclusion) :} Résumer l'apport du document à l'histoire du thème étudié (ex: basculement vers l'indépendance, rupture ou continuité). Rédiger une phrase de conclusion qui répond à la problématique implicite : « En quoi ce document illustre-t-il... ? »
\end{enumerate}
\par \textbf{Reflexes Bac :} Toujours citer le document (\guillemotleft ... \guillemotright) pour appuyer son propos. Ne pas faire de paraphrase (résumé), mais de l'explication. Distinguer ce que le document \emph{dit} (explicite) de ce qu'il \emph{révèle} (implicite/contexte).
\end{methode}

\begin{exoresolu}[Commentaire de document : L'indépendance du Ghana (1957)]
\textbf{Énoncé :} À l'aide du document ci-dessous et de vos connaissances, montrez en quoi la proclamation d'indépendance du Ghana par Kwame Nkrumah le 6 mars 1957 marque une rupture dans l'histoire de la décolonisation africaine.

\textbf{Document :} Extrait du discours de Kwame Nkrumah, Accra, 6 mars 1957.
\guillemotleft \emph{Enfin, la bataille est terminée ! Et ainsi, le Ghana, votre pays bien-aimé, est libre pour toujours ! [...] Notre indépendance est sans signification à moins qu'elle ne soit liée à la libération totale de l'Afrique. [...] Nous allons créer notre propre personnalité et notre propre identité africaine.} \guillemotright

\tcblower
\textbf{Solution détaillée :}

\textbf{1. Identification}
\begin{itemize}
    \item \textbf{Nature :} Discours politique officiel, proclamation solennelle.
    \item \textbf{Auteur :} \cle{Kwame Nkrumah}, leader du \cle{Convention People's Party (CPP)}, Premier ministre de la colonie de la Côte-de-l'Or, premier chef de gouvernement africain noir au sud du Sahara.
    \item \textbf{Date :} 6 mars 1957 (minuit). Contexte immédiat : Cérémonie officielle de remise du pouvoir par le Royaume-Uni (Duchesse de Kent), fin du statut de colonie.
    \item \textbf{Source :} Archives officielles ghanéennes / presse internationale.
\end{itemize}

\textbf{2. Analyse du contenu}
\begin{itemize}
    \item \textbf{Idée force 1 : La fin de la tutelle coloniale.} \guillemotleft La bataille est terminée ! \guillemotright \guillemotleft le Ghana est libre pour toujours \guillemotright. Le vocabulaire militaire (\cle{bataille}) et l'affirmation définitive (\cle{libre pour toujours}) marquent la rupture juridique et politique : le Ghana devient un État souverain, membre du Commonwealth.
    \item \textbf{Idée force 2 : La dimension panafricaine.} \guillemotleft Notre indépendance est sans signification à moins qu'elle ne soit liée à la libération totale de l'Afrique \guillemotright. Nkrumah dépasse le cadre national. Il pose le principe de la \cle{solidarité africaine} et fait du Ghana le \cle{foyer de la révolution panafricaine}.
    \item \textbf{Idée force 3 : La construction identitaire.} \guillemotleft Créer notre propre personnalité \guillemotright. Volonté de rompre avec l'aliénation culturelle coloniale (renommage du pays : Gold Coast $\to$ Ghana, référence à l'empire ancien du Ghana/Wagadou).
\end{itemize}

\textbf{3. Mise en perspective (Connaissances mobilisées)}
\begin{itemize}
    \item \textbf{Antécédents / Causes :} Lutte politique non-violente (\cle{Positive Action}), grève générale (1950), victoire électorale massive du CPP (1951, 1954, 1956). Politique britannique de \cle{transfert progressif de pouvoirs} (contrairement à la France en Algérie).
    \item \textbf{Portée immédiate :} Effet \cle{boule de neige} (domino) : Guinée (1958), Nigeria (1960), etc. Le Ghana devient le \cle{modèle de la décolonisation négociée} britannique en Afrique noire.
    \item \textbf{Portée à long terme :} Organisation de la \cle{Conférence des États africains indépendants} (Accra, 1958), puis rôle moteur dans la création de l'\cle{OUA} (1963).
\end{itemize}

\textbf{4. Synthèse / Conclusion}
Le document acte la naissance du premier État d'Afrique subsaharienne (ancienne colonie) à accéder à l'indépendance par la voie politique et pacifique. Il révèle que pour Nkrumah, la souveraineté nationale n'est qu'une \cle{étape} vers l'unité continentale. Ce discours fonde le \cle{panafricanisme politique} comme idéologie directrice de la décolonisation des années 1960.

\par \textbf{Phrase de conclusion type Bac :} \emph{En proclamant l'indépendance du Ghana comme le point de départ de la libération totale du continent, Nkrumah transforme un événement national en acte fondateur du panafricanisme, faisant du Ghana le symbole et le moteur de la décolonisation de l'Afrique subsaharienne.}
\end{exoresolu}

\begin{methode}[Rédiger une dissertation historique : structure et argumentation]
\textbf{Objectif :} Répondre à une problématique historique par une démonstration structurée, nuancée et étayée d'exemples précis.
\par \textbf{Démarche en 5 étapes :}
\begin{enumerate}
    \item \textbf{Analyse du sujet :} Souligner les mots-clés (dates, bornes chronologiques, concepts : \cle{néocolonialisme}, \cle{émancipation}, \cle{guérilla}). Définir les termes. Formuler la \cle{problématique} (question ouverte à laquelle le plan répond).
    \item \textbf{Recherche d'idées / Plan détaillé :} Classer les arguments en 2 ou 3 grandes parties équilibrées (souvent : I. Contexte/Causes $\to$ II. Déroulement/Modalités $\to$ III. Conséquences/Bilan OU I. Facteurs internes $\to$ II. Facteurs externes $\to$ III. Héritages). Pour chaque partie, lister 2 arguments + 1 exemple précis (date, chiffre, nom, loi).
    \item \textbf{Rédaction de l'introduction :} \textbf{Accroche} (contexte général) $\to$ \textbf{Définitions} $\to$ \textbf{Problématique} $\to$ \textbf{Annonce du plan} (mots de liaison : \emph{Dans un premier temps... Ensuite... Enfin...}).
    \item \textbf{Rédaction du développement :} Une partie = un paragraphe (ou deux). Structure du paragraphe : \textbf{Phrase thèse} (idée directrice) $\to$ \textbf{Arguments expliqués} (causalité) $\to$ \textbf{Exemples précis} (preuves) $\to$ \textbf{Mini-synthèse/Transition}.
    \item \textbf{Rédaction de la conclusion :} \textbf{Bilan} (réponse synthétique à la problématique) $\to$ \textbf{Ouverture} (prolongement chronologique, thématique ou comparatif).
\end{enumerate}
\par \textbf{Critères de réussite au Bac :} Respect de la chronologie, vocabulaire précis (\cle{protectorat}, \cle{assimilation}, \cle{loi-cadre}, \cle{FNLA}, \cle{MPLA}), nuance (pas de manichéisme), lisibilité (saut de ligne entre parties).
\end{methode}

\begin{exoresolu}[Dissertation : L'émancipation des territoires français : le cas de l'Algérie (1954-1962)]
\textbf{Énoncé :} \guillemotleft La guerre d'Algérie (1954-1962) est-elle le produit exclusif de la volonté nationaliste algérienne ou l'échec de la politique d'intégration française ? \guillemotright

\tcblower
\textbf{Solution détaillée :}

\textbf{1. Analyse \& Problématique}
\emph{Mots-clés :} Guerre d'Algérie, 1954-1962, nationalisme (FLN), intégration (départements français), échec.
\emph{Problématique :} Dans quelle mesure l'indépendance de l'Algérie résulte-t-elle de la conjonction d'une revendication nationaliste radicalisée par la répression et de l'incapacité structurelle de la France à réformer son système colonial d'assimilation ?

\textbf{2. Plan détaillé}
\begin{itemize}
    \item \textbf{I. L'échec structurel de l'\cle{assimilation} et de la \cle{loi-cadre} (1945-1954) :} La France refuse l'égalité réelle (code de l'indigénat, statut personnel, inégalités électorales). Répression de Sétif/Guelma (mai 1945) $\to$ radicalisation. Loi-cadre Defferre (1956) arrive trop tard, perçue comme un leurre.
    \item \textbf{II. La guerre : expression de la volonté nationaliste (FLN) face à la répression française (1954-1958) :} Insurrection du 1er novembre 1954 (\cle{Toussaint rouge}). Stratégie FLN : internationalisation (ONU), maquis, terrorisme urbain (bataille d'Alger 1957). Réponse française : \cle{quadrillage}, torture, \cle{plan de Constantine}, rappel du général de Gaulle (mai 1958).
    \item \textbf{III. L'issue négociée : des Accords d'Évian à l'indépendance (1962) :} Revirement de Gaulle (autodétermination 1959, référendum 1961). Négociations Évian (mars 1962). Indépendance 5 juillet 1962. Conséquences : Exode \cle{Pieds-Noirs} / \cle{Harkis}, instabilité algérienne, rupture franco-algérienne.
\end{itemize}

\textbf{3. Introduction (modèle)}
[Accroche] L'Algérie, « département français » depuis 1848, est le joyau de l'empire colonial. [Définitions] L'assimilation promet l'égalité ; le nationalisme algérien (FLN) réclame la souveraineté. [Problématique] Reprise ci-dessus. [Plan] Annoncé.

\textbf{4. Extrait de développement (Partie I)}
\emph{Phrase thèse :} L'origine du conflit réside dans le refus persistant de la France d'appliquer en Algérie les principes républicains d'égalité, transformant une revendication réformiste en revendication séparatiste.
\emph{Arguments :} Depuis 1848, l'Algérie est divisée en 3 départements. Mais la population \cle{musulmane} (terme administratif) est soumise au \cle{code de l'indigénat} (justice d'exception) et au \cle{statut personnel} (droit coranique pour la famille), exclus de la citoyenneté pleine (sauf naturalisation rare). Le scrutin est inégalitaire (collèges électoraux séparés, poids du vote européen supérieur).
\emph{Exemple précis :} Massacres de \cle{Sétif, Guelma, Kherrata} (mai 1945) : manifestation nationaliste réprimée dans le sang (estimations : 1 000 à 45 000 morts). Tournant : l'espoir réformiste (Manifeste du peuple algérien, Ferhat Abbas) s'effondre.
\emph{Mini-synthèse :} L'intégration est un leurre ; la violence d'État valide le choix de la lutte armée par le FLN.

\textbf{5. Conclusion (modèle)}
[Bilan] L'indépendance est le fruit de la détermination du FLN (interne/extérieur) ET de l'impasse française (coût humain/financier, isolement diplomatique, crise de la IVe République). Ce n'est pas l'un \emph{ou} l'autre, mais l'interaction des deux. [Ouverture] Héritage conflictuel : mémoire blessée (Harkis, disparus), relations franco-algériennes complexes, modèle de \cle{guérilla} pour Angola/Mozambique.

\par \textbf{Note de correction :} L'élève doit éviter de raconter la bataille d'Alger sans l'analyser comme un tournant médiatique/diplomatique. Il doit citer \cle{Ferhat Abbas}, \cle{Ben Bella}, \cle{de Gaulle}, \cle{Salan}, \cle{OAS}.
\end{exoresolu}

\begin{methode}[Comparer des processus historiques : le tableau comparatif argumenté]
\textbf{Objectif :} Mettre en parallèle plusieurs cas d'études (pays, colonies, périodes) pour faire ressortir des similitudes (constantes) et des différences (spécificités) structurantes.
\par \textbf{Démarche en 4 étapes :}
\begin{enumerate}
    \item \textbf{Définir les critères de comparaison pertinents :} Ne pas comparer « tout ». Choisir 4 à 5 critères d'analyse imposés par le sujet : \cle{Nature du colonisateur} (métropole), \cle{Mode d'administration} (assimilation/association/indirect rule), \cle{Type de résistance} (politique/pacifique/armée), \cle{Rôle des élites/Partis}, \cle{Issu/Indépendance} (négociée/imposée/guérilla), \cle{Héritage immédiat}.
    \item \textbf{Remplir le tableau par cases :} Une case = un fait précis + une date/chiffre. Éviter les phrases vagues. Utiliser des abréviations claires.
    \item \textbf{Rédiger la synthèse comparative (paragraphe) :} Ne pas se contenter de lire le tableau. Croiser les lignes : « Contrairement au Ghana (britannique, négocié), l'Angola (portugais) connaît une guérilla prolongée car... ». Utiliser les connecteurs logiques : \emph{alors que, contrairement à, de même que, en revanche, par ailleurs}.
    \item \textbf{Tirer une conclusion historiographique :} Que nous apprend la comparaison sur le thème général (décolonisation) ? Ex: Le mode de colonisation détermine largement le mode de décolonisation.
\end{enumerate}
\par \textbf{Reflexes Bac :} Le tableau est un outil de brouillon (ou annexé si demandé). La copie rendue doit contenir le \emph{texte comparatif}. Ne pas faire deux dissertations côte à côte (plan parallèle), mais un plan \emph{thématique croisé}.
\end{methode}

\begin{exoresolu}[Comparaison : Décolonisation britannique (Ghana) vs Française (Algérie) vs Portugaise (Angola)]
\textbf{Énoncé :} À l'aide d'un tableau comparatif et d'un paragraphe de synthèse, mettez en évidence les spécificités des trois grands types de décolonisation en Afrique (britannique, française, portugaise) à travers les cas du Ghana, de l'Algérie et de l'Angola.

\tcblower
\textbf{Solution détaillée :}

\textbf{1. Tableau comparatif (Brouillon / Annexe)}

\begin{center}
\begin{tabularx}{\linewidth}{|l|X|X|X|}\hline
\rowcolor{popblueL}
\textbf{Critères} & \textbf{Ghana (Britannique)} & \textbf{Algérie (Française)} & \textbf{Angola (Portugaise)} \\ \hline
\textbf{Statut colonial} & Colonie de la Couronne (Indirect Rule) & 3 Départements français (Assimilation théorique) & Province d'outre-mer (Intégration tardive, travail forcé) \\ \hline
\textbf{Métropole / Contexte} & UK affaibli (1945), pragmatique, Commonwealth & IVe République instable, puis Ve République (De Gaulle) & Dictature salazariste (Estado Novo), refus catégorique \\ \hline
\textbf{Acteurs nationalistes} & \cle{CPP} (Nkrumah), parti de masse, élection & \cle{FLN} (Front unique), lutte armée, GPRA (gouvernement provisoire) & \cle{MPLA} (Neto, marxiste), \cle{FNLA} (Roberto), \cle{UNITA} (Savimbi) - Guerre civile immédiate \\ \hline
\textbf{Modalités de lutte} & \cle{Action positive} (grèves, boycott), élections, négociations & \cle{Guerre totale} (1954-62), guérilla, terrorisme, torture, OAS & \cle{Guerre de libération} (1961-74), guérilla rurale, fronts multiples, intervention étrangère (Cuba, RSA, URSS, USA) \\ \hline
\textbf{Déclencheur de l'indépendance} & Victoire électorale 1956 $\to$ Indépendance négociée 1957 & Referendum autodétermination (1961) + Accords d'Évian (1962) & \cle{Révolution des Œillets} (Lisbonne, avril 1974) $\to$ Accords d'Alvor (1975) \\ \hline
\textbf{Date Indépendance} & 6 mars 1957 (1er Afrique subsaharienne) & 5 juillet 1962 & 11 novembre 1975 \\ \hline
\textbf{Héritage immédiat} & Stabilité relative (puis coup d'État 1966), modèle panafricain & Exode Pieds-Noirs/Harkis, pouvoir militaire (Boumediene), guerre froide & \cle{Guerre civile} (1975-2002), enjeu pétrolier/diamantifère, proxy war Guerre Froide \\ \hline
\end{tabularx}
\end{center}

\textbf{2. Paragraphe de synthèse comparative (Rédaction attendue)}

La comparaison des trois cas révèle que \textbf{la nature du pouvoir métropolitain détermine le calendrier et la violence de la décolonisation}.
\par \textbf{Similitude :} Dans les trois cas, le nationalisme africain s'affirme après 1945 (conférence de Brazzaville 1944, Bandung 1955) et s'appuie sur des partis structurés (CPP, FLN, MPLA/FNLA/UNITA).
\par \textbf{Différence majeure 1 : La négociation vs la guerre.} Le Royaume-Uni, pragmatique, accepte le \cle{transfert de pouvoirs} par étapes (élections, autonomie interne) dès les années 1950. La France tente l'\cle{assimilation} (Algérie française) puis la répression militaire, ne négociant (Évian) que sous la contrainte de la crise de la IVe République. Le Portugal (dictature) refuse tout dialogue jusqu'à l'effondrement du régime en 1974 (\cle{Révolution des Œillets}).
\par \textbf{Différence majeure 2 : L'unité nationaliste.} Au Ghana, le CPP hégémonique assure une transition unitaire. En Algérie, le FLN élimine ses rivaux (MNA) pour imposer l'État unitaire. En Angola, la \cle{fractionnement} des mouvements (MPLA pro-soviétique, FNLA pro-chinois/occidental, UNITA pro-sud-africain/US) transforme la décolonisation en \cle{guerre civile par procuration} (Guerre Froide), empêchant la construction étatique immédiate.
\par \textbf{Conclusion :} Le Ghana incarne la \cle{décolonisation par le haut} (négociée), l'Algérie la \cle{décolonisation par le bas} (violente, imposée), l'Angola la \cle{décolonisation par l'effondrement métro} (tardive, internationalisée, conflictuelle).

\par \textbf{Note :} L'élève doit maîtriser le vocabulaire : \cle{Indirect Rule}, \cle{Assimilation}, \cle{Estado Novo}, \cle{Guerre Froide}, \cle{Proxy war}.
\end{exoresolu}

\begin{methode}[Expliquer la causalité historique : l'analyse multicausale]
\textbf{Objectif :} Expliquer un événement historique majeur (indépendance, guerre, basculement) en articulant des causes de nature différente (structurelles, conjoncturelles, déclenchantes) et des échelles différentes (internes, externes).
\par \textbf{Démarche en 3 temps (Modèle « Oignon » ou « Arbre ») :}
\begin{enumerate}
    \item \textbf{Les causes structurelles (Long terme / Racines) :} Facteurs profonds, lents, souvent économiques, sociaux, démographiques, idéologiques. Ex: \cle{Exploitation coloniale}, \cle{naissance d'une élite occidentalisée}, \cle{idéologie panafricaine}, \cle{faiblesse économique métropole post-1945}.
    \item \textbf{Les causes conjoncturelles (Moyen terme / Branches) :} Événements, réformes, crises qui préparent le terrain. Ex: \cle{Seconde Guerre mondiale} (participation tirailleurs, affaiblissement Europe), \cle{Guerre froide} (blocs US/URSS anti-coloniaux), \cle{Charte de l'ONU} (droit des peuples), \cle{Conférence de Bandung} (1955).
    \item \textbf{Les causes déclenchantes / immédiates (Court terme / Feuille) :} L'étincelle, la décision, l'événement précis qui fait basculer. Ex: \cle{1er Novembre 1954} (Algérie), \cle{Référendum 1958} (Guinée), \cle{Révolution des Œillets 1974} (Angola), \cle{Grève générale 1950} (Ghana).
\end{enumerate}
\par \textbf{Rédiger l'explication :} Utiliser les connecteurs de causalité : \emph{C'est parce que... que... ; Le facteur décisif est... ; Sans ..., ... n'aurait pas eu lieu ; La conjonction de ... et ... explique...}. Hiérarchiser : « Si la cause structurelle est nécessaire, elle n'est pas suffisante sans le déclencheur... ».
\end{methode}

\begin{exoresolu}[Analyse causale : Pourquoi le « Non » de la Guinée au référendum de 1958 ?]
\textbf{Énoncé :} Expliquez pourquoi la Guinée, seule parmi les territoires d'Afrique occidentale française, a voté « Non » au référendum du 28 septembre 1958 sur la Constitution de la Ve République, accédant ainsi immédiatement à l'indépendance.

\tcblower
\textbf{Solution détaillée :}

\textbf{1. Causes structurelles (Long terme)}
\begin{itemize}
    \item \textbf{Tradition de résistance :} Histoire de l'\cle{Imamat du Fouta-Djalon} (Almamy Bokar Biro, vaincu 1896) et de l'empire \cle{Wassoulou} de \cle{Samory Touré} (résistance 1880-1898). Mémoire collective de souveraineté.
    \item \textbf{Structure syndicale forte :} Le \cle{RDA} (Rassemblement Démocratique Africain) est très implanté via les syndicats (CGAT). \cle{Sékou Touré}, leader syndical (secrétaire général CGAT), a une base ouvrière/paysanne solide, contrairement aux élites « assimilées » d'autres territoires.
    \item \textbf{Refus de l'assimilation :} Rejet du statut de « citoyen français » perçu comme aliénant. Volonté de \cle{dignité} et de \cle{souveraineté pleine}.
\end{itemize}

\textbf{2. Causes conjoncturelles (Moyen terme : 1945-1958)}
\begin{itemize}
    \item \textbf{Loi-cadre Defferre (1956) :} Donne l'autonomie interne mais maintient la tutelle française (diplomatie, défense, monnaie). Perçue comme un piège (\cle{néocolonialisme} avant l'heure) par l'aile radicale du RDA guinéen.
    \item \textbf{Élections de 1957 :} Victoire écrasante du \cle{PDG-RDA} (56 sièges sur 60). Légitimité démocratique incontestable pour exiger la souveraineté.
    \item \textbf{Contexte international :} Indépendance du Ghana (mars 1957) : effet de démonstration immédiat. Conférence de Bandung (1955) : légitimité du non-alignement.
\end{itemize}

\textbf{3. Cause déclenchante (Court terme : Septembre 1958)}
\begin{itemize}
    \item \textbf{Le référendum du 28 septembre 1958 :} De Gaulle propose la \cle{Communauté française} : autonomie large mais pas d'indépendance (pas d'armée, pas de monnaie, pas de diplomatie propre). Le « Oui » = Communauté. Le « Non » = Indépendance immédiate + rupture aide française.
    \item \textbf{Le discours de De Gaulle à Conakry (août 1958) :} Menaces voilées : « L'indépendance, c'est la misère ». Cela braque l'opinion et Sékou Touré.
    \item \textbf{La campagne du « Non » :} Slogan : \guillemotleft \cle{Préférer la liberté dans la pauvreté à la richesse dans l'esclavage} \guillemotright. Mobilisation massive (femmes, jeunes, syndicats).
\end{itemize}

\textbf{4. Hiérarchisation et Conclusion}
Le « Non » guinéen n'est pas un accident. Il est la \textbf{convergence} d'une \cle{identité historique de résistance} (structurel), d'une \cle{organisation politique/syndicale hégémonique} (conjoncturel) et d'une \cle{offre politique métropolitaine jugée inacceptable} (déclencheur : la Communauté). Le courage politique de Sékou Touré (risque économique majeur) transforme le référendum en acte fondateur de la \cle{souveraineté nationale}.

\par \textbf{Phrase clé pour le Bac :} \emph{Le vote du 28 septembre 1958 est l'aboutissement d'un long processus d'affirmation nationale (syndicale et politique) qui a saisi l'opportunité offerte par la crise de la IVe République pour refuser un statut d'autonomie tronquée, faisant de la Guinée le symbole de la rupture radicale avec l'ordre colonial français.}
\end{exoresolu}

\begin{methode}[Caractériser un concept historique : le néocolonialisme]
\textbf{Objectif :} Définir, identifier les mécanismes et illustrer un concept clé du programme (néocolonialisme, impérialisme, développement, sous-développement) pour l'utiliser comme outil d'analyse.
\par \textbf{Démarche en 4 étapes (Fiche concept) :}
\begin{enumerate}
    \item \textbf{Définition rigoureuse :} Citation d'auteur de référence (Kwame Nkrumah, Samir Amin, Jean-Marc Ela) + définition propre : \guillemotleft Situation où un État formellement indépendant voit son économie, sa politique, sa culture orientées par des puissances extérieures (anciennes métropoles, multinationales, institutions financières) \guillemotright.
    \item \textbf{Les 4 piliers (Mécanismes) :} Les classer par domaine.
        \begin{itemize}
            \item \textbf{Économique :} \cle{Spécialisation primaire} (exportation matières premières brutes / importation produits finis), \cle{Termes de l'échange} défavorables, \cle{Endettement} / \cle{PPTE} / \cle{FMI/Banque Mondiale} (Plans d'Ajustement Structurel - PAS), \cle{Franc CFA} (zone monétaire), \cle{Accords de coopération} inégaux.
            \item \textbf{Politique / Militaire :} \cle{Bases militaires}, \cle{Accords de défense}, \cle{Ingérence} (coups d'État, soutien dictateurs), \cle{Françafrique} (réseaux opaques).
            \item \textbf{Culturel / Scientifique :} \cle{Langue officielle} unique, programmes scolaires exogènes, \cle{Fuite des cerveaux} (brain drain), dépendance technologique.
            \item \textbf{Démographique / Migratoire :} Main-d'œuvre peu qualifiée vers le Nord, remises d'argent (transferts) mais perte de capital humain.
        \end{itemize}
    \item \textbf{Acteurs :} Anciennes métropoles (France, UK, Portugal), \cle{Nouveaux acteurs} (Chine, USA, Russie, Turquie, Émirats), \cle{Firmes multinationales} (TotalEnergies, Areva/Orano, Bolloré, miniers), \cle{Institutions de Bretton Woods} (FMI, BM), \cle{Compradores} (élites locales corrompues/cooptées).
    \item \textbf{Exemples concrets (Cameroun / Afrique centrale) :} \cle{Pétrole / Bois / Cacao / Coton} exportés bruts. \cle{Franc CFA} (garantie Trésor français, parité fixe euro). \cle{Accords de défense} (bases de Douala, Yaoundé, N'Djamena). \cle{PAS} (années 80-90) : dévaluation 1994, privatisations, baisse dépenses santé/éducation.
\end{enumerate}
\par \textbf{Attention :} Ne pas confondre \cle{néocolonialisme} (structure systémique) et simple \cle{coopération} ou \cle{aide au développement}. Le néocolonialisme implique une \cle{asymétrie de pouvoir} structurelle reproduisant le schéma colonial.
\end{methode}

\begin{exoresolu}[Analyse de situation : Le Franc CFA, instrument de néocolonialisme monétaire ?]
\textbf{Énoncé :} En vous appuyant sur vos connaissances, montrez en quoi le Franc CFA (Zone Franc) peut être analysé comme un mécanisme néocolonial, tout en nuanceant ce diagnostic par les arguments de ses défenseurs.

\tcblower
\textbf{Solution détaillée :}

\textbf{1. Thèse : Le Franc CFA comme pilier du néocolonialisme économique (Arguments critiques)}
\begin{itemize}
    \item \textbf{Deux zones, une logique :} La zone Franc comprend l'\cle{UEMOA} (BCEAO, Afrique de l'Ouest) et la \cle{CEMAC} (BEAC, Afrique Centrale, dont le Cameroun). Historiquement, les deux banques centrales devaient déposer une partie de leurs réserves (65\% puis 50\%) sur un \cle{compte d'opérations} au \cle{Trésor français}, qui garantissait la convertibilité.
    \item \textbf{Réformes récentes (2019-2024) :} Pour l'UEMOA (2019) : suppression du compte d'opérations, retrait des représentants français du conseil d'administration, changement de nom prévu (Eco) mais reporté. Pour la CEMAC (BEAC, 2019-2023) : suppression du compte d'opérations et retrait des administrateurs français actés, mais \textbf{parité fixe avec l'Euro maintenue} (1 Euro = 655,957 FCFA) et \textbf{garantie de convertibilité illimitée par la France conservée}.
    \item \textbf{Perte de souveraineté monétaire effective :} Taux de change immuable $\to$ pas de dévaluation possible pour relancer les exportations (compétitivité-prix). Politique monétaire importée de la BCE (cible inflation zone euro), inadaptée aux économies africaines (besoin de croissance, emploi).
    \item \textbf{Libre transfert des capitaux :} Garanti par les accords. Facilite la \cle{fuite des capitaux} (évasion fiscale, rapatriement bénéfices multinationales) vers la France/UE. Pas de contrôle des changes.
    \item \textbf{Crédit orienté :} Historiquement, le crédit bancaire finance l'importation et l'exportation (secteur primaire), peu l'industrie locale ou l'agriculture vivrière.
\end{itemize}

\textbf{2. Antithèse : Les arguments de la « garantie » et de la stabilité (Arguments défenseurs / officiels)}
\begin{itemize}
    \item \textbf{Stabilité macroéconomique :} Inflation maîtrisée (moyenne ~2-3\% vs 10-20\% dans pays hors zone). Confiance des investisseurs (convertibilité garantie par l'État français).
    \item \textbf{Facilitation des échanges :} Monnaie commune au sein de la CEMAC/UEMOA (intégration régionale). Pas de risque de change intra-zone.
    \item \textbf{Filet de sécurité :} Garantie de conversion illimitée par le Trésor français (crédibilité). Aide à l'importation de biens d'équipement/médicaments.
\end{itemize}

\textbf{3. Synthèse nuancée (Niveau Bac « Excellence »)}
Le Franc CFA est un \cle{dispositif néocolonial structurel} (conçu à la Libération pour lier les colonies à la métropole, assurer matières premières bon marché et marchés captifs). Les réformes 2019-2023 en ont \textbf{atténué la forme} (fin du compte d'opérations, retrait représentants français) mais \textbf{pas le fond} : la parité fixe avec l'Euro et la garantie française maintiennent une \cle{dépendance monétaire} qui prive les États (notamment du Cameroun via la BEAC) d'un instrument majeur de politique économique (taux de change, taux d'intérêt directeurs).
\par \textbf{Conclusion :} C'est un \cle{frein à l'industrialisation} (surévaluation du change = importations moins chères que production locale) et un \cle{levier d'influence politique} français (droit de regard sur politiques budgétaires via FMI/zone franc). La sortie (projet \cle{Eco} CEDEAO) bute sur les divergences économiques (Nigeria vs zone CFA) et la réticence française à perdre son pré-carré monétaire.

\par \textbf{Mots-clés à placer :} \cle{Souveraineté monétaire}, \cle{Parité fixe}, \cle{Compte d'opérations}, \cle{Trésor français}, \cle{Convertibilité}, \cle{Compétitivité-prix}, \cle{Désindustrialisation}, \cle{Réforme 2019}, \cle{Eco}, \cle{BEAC}, \cle{BCEAO}, \cle{CEMAC}, \cle{UEMOA}.
\end{exoresolu}

\begin{methode}[Exploiter une carte / un croquis schématique : commentaire cartographique]
\textbf{Objectif :} Transformer une information visuelle (carte, croquis, flux) en démonstration écrite structurée.
\par \textbf{Démarche en 4 étapes :}
\begin{enumerate}
    \item \textbf{Lecture technique (Légende + Titre + Échelle + Projection) :} Lire \emph{tous} les items de la légende. Identifier le thème (Indépendances, Néocolonialisme, Flux migratoires, Bases militaires). Noter l'échelle (planétaire, continental, national) et la date de la situation cartographiée.
    \item \textbf{Description structurée (Le « Quoi » et le « Où ») :} Organiser la description par \textbf{grands ensembles spatiaux} (régions, pays, façades maritimes) et non par item de légende. Utiliser le vocabulaire spatial : \cle{Discontinuité}, \cle{Gradient}, \cle{Pôle}, \cle{Axe}, \cle{Interface}, \cle{Enclavement}. Citer les noms propres (villes, pays, ressources).
    \item \textbf{Explication / Analyse (Le « Pourquoi » et le « Comment ») :} Mobiliser le cours pour expliquer les figures observées. Relier les flux aux ressources, les bases aux accords de défense, les dates d'indépendance aux métropoles. Montrer les \cle{logiques d'acteurs} (États, firmes, organisations internationales).
    \item \textbf{Synthèse / Problématisation :} Dégager la dynamique principale (ex: \guillemotleft La carte révèle la permanence des liens Nord-Sud asymétriques malgré les indépendances politiques \guillemotright). Proposer un titre de croquis amélioré si demandé.
\end{enumerate}
\par \textbf{Code couleur mental :} \textbf{Rouge} = Conflits / Tensions / Ruptures. \textbf{Bleu} = Flux dominants / Dépendances. \textbf{Vert} = Résistances / Alternatives / Intégration régionale. \textbf{Noir} = Structures lourdes (frontières, capitales).
\end{methode}

\begin{exoresolu}[Commentaire de croquis : L'Afrique dans la mondialisation néocoloniale (années 2020)]
\textbf{Énoncé :} Le croquis ci-dessous représente les principaux flux économiques et militaires liant l'Afrique aux puissances extérieures. En vous appuyant sur la légende et vos connaissances, commentez ce croquis en montrant qu'il illustre la persistance de rapports néocoloniaux.

\textbf{Légende simulée (à décrire mentalement) :}
\begin{itemize}
    \item Flèches épaisses \textbf{Rouge} (Exportations) : Pétrole (Golfe de Guinée, Angola), Minerais (Cu/Co Katanga, Uranium Niger, Or Sahel), Cacao/Café (Côte d'Ivoire, Ghana, Cameroun) $\to$ Europe, Chine, USA.
    \item Flèches fines \textbf{Bleu} (Importations) : Produits manufacturés, céréales, médicaments, hydrocarbures raffinés $\leftarrow$ Europe, Chine, Inde.
    \item \textbf{Étoiles Jaunes} : Bases militaires françaises (Djibouti, Gabon, Sénégal, Côte d'Ivoire, Tchad), US (Djibouti, Niger - retrait 2024), Chinoises (Djibouti), Russes (RCA, Mali, Burkina).
    \item \textbf{Trait discontinu Vert} : Nouvelles routes de la Soie (Chine) : Ports (Mombasa, Djibouti, Lagos, Abidjan), Chemins de fer (Addis-Abeba/Djibouti, Lagos/Kano).
    \item \textbf{Cercles Orange} : Sièges institutions financières (FMI/BM - Washington/Paris), Sièges multinationales (TotalEnergies, Bolloré, Huawei, Rosatom).
\end{itemize}

\tcblower
\textbf{Solution détaillée :}

\textbf{1. Introduction}
Ce croquis met en lumière la \cle{structure extravertie} des économies africaines et la \cle{concurrence géopolitique} pour l'accès aux ressources, confirmant la thèse du néocolonialisme : indépendance politique (années 60) mais dépendance économique et stratégique renforcée.

\textbf{2. Description structurée : Une Afrique « perforée » par les flux sortants}
\begin{itemize}
    \item \textbf{L'exportation de la rente primaire (Flèches rouges massives) :} L'Afrique reste le \cle{réservoir de matières premières} du monde. Trois pôles exportateurs majeurs :
        \begin{itemize}
            \item \textbf{Golfe de Guinée / Afrique Centrale :} Pétrole (Nigeria, Angola, Gabon, Congo, Guinée Eq., Cameroun, Tchad) + Bois/Minerais. Flux vers Europe (TotalEnergies, ENI, Shell) et Chine (CNPC, Sinopec).
            \item \textbf{Afrique Australe / Cuivre :} Cuivre/Cobalt (RDC, Zambie) $\to$ Chine (quasi-monopole traitement).
            \item \textbf{Sahel / Afrique de l'Ouest :} Uranium (Niger - ex-Areva/Orano), Or (Mali, Burkina, Ghana), Cacao (Côte d'Ivoire, Ghana, Cameroun - 60\% marché mondial).
        \end{itemize}
    \item \textbf{L'importation de biens finis et alimentaires (Flèches bleues diffuses) :} Déséquilibre structurel de la balance commerciale. Dépendance alimentaire (blé, riz, huile) et technologique.
    \item \textbf{La militarisation des enjeux (Étoiles) :} Densification des bases \textbf{Françaises} (historique, accords défense) et \textbf{Américaines} (lutte anti-terroriste Sahel) au Sahel et Golfe de Guinée. \textbf{Nouveaux entrants :} Base chinoise (Djibouti - 2017), présence \textbf{Russe} (Wagner/Africa Corps) au Sahel (Mali, Burkina, Niger) et RCA (ressources minières contre sécurité régime).
\end{itemize}

\textbf{3. Explication : Les mécanismes néocoloniaux spatialisés}
\begin{itemize}
    \item \textbf{Spécialisation bloquée (Héritage colonial) :} Infrastructures (chemins de fer, ports) conçues \textbf{vers la côte} (export), pas \textbf{vers l'intérieur} (intégration régionale). Ex: Chemin de fer Benguela (Angola), Transcamerounais, Dakar-Niger.
    \item \textbf{Accords inégaux :} \cle{Conventions d'établissement} (droit fiscal avantageux pour multinationales), \cle{Accords de défense} (droit de préposition militaire), \cle{Franc CFA} (zone monétaire unique facilitant rapatriement profits).
    \item \textbf{Nouvelle donne : La Chine.} Partenaire commercial n°1 de l'Afrique. Modèle \cle{Infrastructures contre Ressources} (prêts garantis sur ressources). Crée de la dette (risque \cle{piège à dette} : port de Mombasa/Kenya, chemin de fer Addis-Abeba/Djibouti). Mais offre alternative au monopole occidental (prix plus bas, non-conditionnalité politique).
\end{itemize}

\textbf{4. Synthèse / Conclusion}
La carte révèle une \cle{triple dépendance} : \textbf{Économique} (spécialisation primaire, termes de l'échange défavorables), \textbf{Financière} (dette, FMI, Franc CFA), \textbf{Sécuritaire} (bases extérieures, mercenaires).
\par Cependant, elle montre aussi des \textbf{marges de manœuvre} : multiplication des partenaires (Chine, Turquie, Russie, Émirats, Inde) affaiblit le \cle{pré-carré} français/européen. Emergence de \cle{pôles régionaux} (Nigeria, Afrique du Sud, Éthiopie, Kenya, Maroc, Égypte) et projets d'intégration (\cle{ZLECAf} - Zone de Libre-Échange Continentale Africaine) tentent de réorienter les flux \textbf{intra-africains} (actuellement ~15-18\% seulement).

\par \textbf{Titre amélioré pour le croquis :} \emph{« Afrique : l'enjeu des ressources face à la multipolarité mondiale – Permanence de l'extraversion et recomposition des dépendances (2020) »}.
\end{exoresolu}

\begin{methode}[Rédiger une biographie historique / un portrait d'acteur : le « Grand Homme » dans son temps]
\textbf{Objectif :} Présenter un acteur majeur (chef d'État, résistant, intellectuel) sans tomber dans l'hagiographie (culte du héros) ni la psychologisation, en l'insérant dans les structures et conjonctures.
\par \textbf{Structure type (5 paragraphes) :}
\begin{enumerate}
    \item \textbf{Origines \& Formation :} Milieu social, éducation (coranique, coloniale, métropole), prise de conscience politique (syndicat, parti, armée).
    \item \textbf{Accession au pouvoir / Leadership :} Contexte de la lutte (élections, guérilla, coup d'État, héritage). Style de pouvoir (charismatique, bureaucratique, militaire, idéologique).
    \item \textbf{Action politique majeure (Bilan) :} 2-3 réalisations concrètes (Indépendance, Unité nationale, Industrialisation, Éducation, Santé, Droits des femmes) + 2-3 échecs/ombres (Dictature, Économie ruinée, Guerre civile, Corruption, Droits de l'homme).
    \item \textbf{Politique étrangère / Panafricanisme :} Positionnement Guerre Froide / Non-alignement. Rôle OUA/UA. Relations ex-métropole / Nouvelles puissances.
    \item \textbf{Héritage \& Historiographie :} Comment est-il perçu aujourd'hui ? Figure tutélaire / Contestée. « Père de la nation » vs « Dictateur ». Nuance : \emph{Produit de son temps (décolonisation, Guerre Froide) autant que faiseur d'histoire.}
\end{enumerate}
\par \textbf{Vocabulaire d'analyse :} \cle{Charisme}, \cle{Légitimité} (traditionnelle, charismatique, rationnelle-légale - Weber), \cle{Patrimonialisme}, \cle{Clientélisme}, \cle{État-parti}, \cle{Parti unique}, \cle{Démocratie guidée}.
\end{methode}

\begin{exoresolu}[Portrait croisé : Kwame Nkrumah (Ghana) vs Sékou Touré (Guinée) vs Agostinho Neto (Angola)]
\textbf{Énoncé :} Présentez brièvement le parcours et le bilan de ces trois « Pères de l'indépendance ». En quoi leurs trajectoires illustrent-elles la diversité des chemins vers la souveraineté ?

\tcblower
\textbf{Solution détaillée :}

\textbf{1. Kwame Nkrumah (Ghana, 1909-1972) : Le Panafricaniste visionnaire}
\begin{itemize}
    \item \textbf{Origines :} Fils d'orpailleur, éducation missionnaire, études USA (Lincoln Univ, Penn), Londres (LSE, droit). Marxisme modéré, panafricanisme (Padmore).
    \item \textbf{Accession :} \cle{Action Positive} (grèves, boycott). Victoire électorale 1951 (premier ministre « en prison »). Indépendance négociée 1957.
    \item \textbf{Bilan :} \textbf{+} Industrialisation (barrage d'Akosombo/Volta, aluminium VALCO), Éducation gratuite, Santé, \cle{OUA} (siège Addis-Abeba), soutien mouvements libération (ANC, PAIGC). \textbf{-} \cle{Parti unique} (1964), \cle{Preventive Detention Act} (prison sans procès), économie ruinée (chute cours cacao, gaspillages), culte de la personnalité (\cle{Osagyefo}).
    \item \textbf{Fin :} Renversé par coup d'État militaire/police (1966) soutenu CIA (allégation) pendant voyage à Pékin/Hanoï. Mort exilé (Roumanie/Guinée).
    \item \textbf{Héritage :} Icône panafricaine mondiale. « Africa must unite ».
\end{itemize}

\textbf{2. Ahmed Sékou Touré (Guinée, 1922-1984) : Le Syndicaliste intransigeant}
\begin{itemize}
    \item \textbf{Origines :} Famille paysanne mandingue, éducation coranique/coloniale, syndicaliste CGT/CGAT. Homme de terrain, orateur hors pair.
    \item \textbf{Accession :} Leader RDA/PDG. \cle{Référendum 1958 : « Non » unique}. Indépendance immédiate 2 oct 1958. Rupture brutale avec France (retrait admin, argent, matériel).
    \item \textbf{Bilan :} \textbf{+} Souveraineté totale, dignité, alphabétisation (langues nationales), santé primaire, rôle médiateur (Accords d'Alger 1975 Iran/Irak), non-alignement actif. \textbf{-} \cle{Régime de terreur} (Camp Boiro, « complots » permanents, répression Peuls/opposants), \cle{Économie dirigiste} ruinée (nationalisations, contrôle changes, marché noir), exil massif (2-3 millions de Guinéens).
    \item \textbf{Fin :} Mort au pouvoir (opération cardiaque Cleveland/USA) 1984. Coup d'État immédiat (Conté).
    \item \textbf{Héritage :} Symbole de la \cle{rupture radicale} et de la \cle{fierté nationale}. Mais modèle autoritaire durable.
\end{itemize}

\textbf{3. Agostinho Neto (Angola, 1922-1979) : Le Poète-Médecin guérilléro-marxiste}
\begin{itemize}
    \item \textbf{Origines :} Méthodiste, études médecine (Coimbra, Lisbonne). Poète engagé. Arrêté par PIDE (police politique). Évasion, exil (Maroc, Zaïre, Congo-Brazza).
    \item \textbf{Accession :} Chef \cle{MPLA} (Mouvement Populaire de Libération de l'Angola - marxiste). Guerre de libération (1961-74) + Guerre civile (vs FNLA/UNITA). Indépendance 11 nov 1975 (jour fuite Portugais). Soutien massif \cle{Cuba} (troupes) / \cle{URSS}.
    \item \textbf{Bilan :} \textbf{+} Indépendance arrachée aux armes, unité nationale relative (MPLA hégémonique), alphabétisation, santé. \textbf{-} \cle{Guerre civile continue} (jusqu'à 2002), \cle{État-parti marxiste-léniniste} rigide, \cle{Purges} (27 mai 1977 - nitistas, milliers de morts), dépendance pétrole/soviétique, népotisme.
    \item \textbf{Fin :} Mort maladie (Moscou) 1979. Successeur Dos Santos (38 ans pouvoir).
    \item \textbf{Héritage :} Héros national (aéroport, université, poèmes). Mais architecte d'un système autoritaire pétrolier.
\end{itemize}

\textbf{4. Analyse comparative (Réponse à la question « Diversité des chemins »)}
\begin{center}
\begin{tabularx}{\linewidth}{|l|X|X|X|}\hline
\rowcolor{popblueL}
\textbf{Critère} & \textbf{Nkrumah (Ghana)} & \textbf{Sékou Touré (Guinée)} & \textbf{Neto (Angola)} \\ \hline
\textbf{Voie d'accès} & \cle{Négociée / Électorale} (Modèle britannique) & \cle{Rupture référendaire} (Refus Communauté) & \cle{Guérilla / Guerre} (Effondrement Portugal) \\ \hline
\textbf{Idéologie} & \cle{Panafricanisme} / Socialisme africain (non-aligné) & \cle{Nationalisme radical} / Socialisme africain (non-aligné) & \cle{Marxisme-Léninisme} (Bloc de l'Est) \\ \hline
\textbf{Appui extérieur} & UK/Commonwealth + Bloc de l'Est (plus tard) & URSS/Chine + Pays non-alignés (rejet France) & \cle{Cuba / URSS} (décisif militairement) \\ \hline
\textbf{Type de régime} & Parti unique (1964), autoritaire modernisateur & Parti unique (PDG), totalitaire, terroriste & Parti unique (MPLA), marxiste, militarisé \\ \hline
\textbf{Héritage principal} & Idée panafricaine, infrastructure (barrage) & Dignité/Souveraineté, traumatisme Boiro & Indépendance armée, guerre civile 27 ans \\ \hline
\end{tabularx}
\end{center}

\par \textbf{Conclusion commune :} Les trois incarnent la \cle{génération des « Pères fondateurs »} : légitimité historique (lutte), projet modernisateur (État-nation, éducation, santé), mais \cle{confiscation de la démocratie} au nom de l'unité et du développement, facilité par le contexte de \cle{Guerre Froide} (soutiens extérieurs conditionnés) et la \cle{fragilité étatique} post-coloniale.
\end{exoresolu}

\begin{attention}[Les 5 erreurs qui coûtent des points au Bac (Histoire - Terminale Tech)]
\begin{enumerate}
    \item \textbf{Confondre « Décolonisation » et « Indépendance » :} La décolonisation est un \cle{processus} (politique, économique, mental, culturel) qui \emph{dépasse} la date officielle de l'indépendance. Ne pas écrire : « La décolonisation du Ghana a eu lieu le 6 mars 1957. » Écrire : « L'indépendance du Ghana (6 mars 1957) est l'aboutissement politique d'un processus de décolonisation entamé dès les années 1940... ».
    \item \textbf{Ignorer la Guerre Froide comme facteur structurel :} Traiter la décolonisation (surtout Algérie, Angola, Congo, Éthiopie) comme une affaire purement locale « Colons vs Colonisés ». \textbf{Obligatoire :} Montrer comment l'URSS/USA/Chine/Cuba instrumentalisent les mouvements (armes, formation, diplomatie, ONU) et comment les métropoles (France, Portugal) justifient leur maintien par l'anticommunisme.
    \item \textbf{Traiter le « Néocolonialisme » comme une simple opinion politique :} C'est un \cle{concept d'analyse scientifique} (Nkrumah, Amin, Suret-Canale). Au Bac, il faut le \textbf{démontrer par des mécanismes concrets} (Franc CFA, Accords de défense, Multinationales, Dette, Termes de l'échange) et non juste l'affirmer (« La France continue de piller »). Nuancer : responsabilité des élites africaines (\cle{compradores}, corruption, mauvaise gouvernance).
    \item \textbf{Erreurs chronologiques / Géographiques « tueuses » :}
    \begin{itemize}
        \item Situer l'indépendance de l'Angola en 1962 (confusion avec Algérie) ou 1957 (Ghana).
        \item Confondre \cle{FNLA} (Holden Roberto, Bakongo, pro-Chine/US), \cle{MPLA} (Neto, Mbundu, marxiste, pro-URSS/Cuba), \cle{UNITA} (Savimbi, Ovimbundu, pro-RSA/US).
        \item Oublier que la \cle{Guinée} vote « Non » en \textbf{1958} (seule), les autres en 1960.
        \item Placer la \cle{Conférence de Bandung} en 1945 ou 1960 (c'est 1955).
    \end{itemize}
    \item \textbf{Rédaction : « Raconter » au lieu d'« Expliquer / Argumenter » :} Le correcteur attend une \textbf{démonstration} (Problématique $\to$ Arguments $\to$ Exemples $\to$ Synthèse), pas un récit chronologique linéaire (« Ensuite... Puis... Enfin... »). Utiliser les connecteurs logiques (\emph{cependant, par conséquent, en effet, en revanche, de surcroît}). \textbf{Sanction :} Hors sujet ou note maximale 6/20 si pas de problématique visible.
\end{enumerate}
\end{attention}

\newpage

\coursec{Exercices}

\courssub{Je vérifie mes connaissances}

\begin{exercice}[QCM : dates, acteurs et notions]
Pour chaque question, entoure la (ou les) bonne(s) réponse(s).

\textbf{1.} La date du \cle{6 mars 1957} correspond à :
\begin{itemize}
\item[a)] l'indépendance de l'Algérie ;
\item[b)] l'indépendance du Ghana ;
\item[c)] l'indépendance de l'Angola ;
\item[d)] le référendum guinéen.
\end{itemize}

\textbf{2.} Le \cle{1\ier{} novembre 1954} marque :
\begin{itemize}
\item[a)] la Toussaint rouge, début de la guerre d'Algérie ;
\item[b)] les accords d'Évian ;
\item[c)] l'indépendance de la Guinée Conakry ;
\item[d)] la création de la CPP.
\end{itemize}

\textbf{3.} Le \cle{MPLA} est :
\begin{itemize}
\item[a)] un parti ghanéen ;
\item[b)] un mouvement de libération de l'Angola dirigé par Agostinho Neto ;
\item[c)] une organisation française ;
\item[d)] un syndicat algérien.
\end{itemize}

\textbf{4.} Kwame Nkrumah est le fondateur :
\begin{itemize}
\item[a)] du FLN ;
\item[b)] de la Convention People's Party (CPP) ;
\item[c)] du MPLA ;
\item[d)] de la Françafrique.
\end{itemize}

\textbf{5.} Le \cle{28 septembre 1958}, la Guinée :
\begin{itemize}
\item[a)] vote << oui >> au référendum de De Gaulle ;
\item[b)] vote << non >> et accède à l'indépendance le 2 octobre 1958 ;
\item[c)] déclenche une guerre d'indépendance ;
\item[d)] rejoint la Communauté française.
\end{itemize}
\end{exercice}

\begin{exercice}[Vrai ou faux ? Justifie ta réponse]
Réponds par vrai ou faux, puis justifie en une ou deux phrases.
\begin{enumerate}
\item Le Ghana est la première colonie d'Afrique noire à accéder à l'indépendance.
\item La guerre d'Algérie s'est terminée par les accords d'Évian signés en mars 1962.
\item L'Angola a obtenu son indépendance le 11 novembre 1975 sans conflit armé.
\item Sékou Touré a accepté la Communauté proposée par le général de Gaulle.
\item Le néocolonialisme désigne la domination économique et politique exercée par les anciennes puissances coloniales après les indépendances.
\end{enumerate}
\end{exercice}

\begin{exercice}[Définitions]
Définis en deux ou trois phrases chacune des notions suivantes, en donnant pour chacune un exemple précis tiré du cours :
\begin{enumerate}
\item \cle{Décolonisation} ;
\item \cle{Néocolonialisme} ;
\item \cle{Françafrique} ;
\item \cle{Mouvement de libération nationale}.
\end{enumerate}
\end{exercice}

\begin{exercice}[Chronologie : remettre en ordre]
Replace dans l'ordre chronologique les événements suivants, en précisant la date exacte de chacun :
\begin{itemize}
\item indépendance de l'Angola ;
\item Toussaint rouge en Algérie ;
\item indépendance du Ghana ;
\item référendum du 28 septembre 1958 en Guinée ;
\item indépendance de l'Algérie.
\end{itemize}
Présente ensuite ces événements sous la forme d'une frise chronologique simple.
\end{exercice}

\courssub{Je m'entraîne}

\begin{exercice}[Question à développement construit : le Ghana]
En une vingtaine de lignes, montre que l'accession du Ghana à l'indépendance (1957) illustre une décolonisation négociée. Tu mobiliseras : le rôle de Kwame Nkrumah et de la CPP, les actions menées (boycott, désobéissance civile, << Positive Action >>) et les étapes institutionnelles de l'émancipation.
\end{exercice}

\begin{exercice}[Question à développement construit : la guerre d'Algérie]
En une vingtaine de lignes, explique pourquoi la décolonisation de l'Algérie a pris la forme d'une guerre (1954--1962). Tu évoqueras : les causes de la radicalisation (colonisation de peuplement, répression), le rôle du FLN, et l'issue du conflit (accords d'Évian, indépendance du 5 juillet 1962).
\end{exercice}

\begin{exercice}[Construction cartographique : l'Angola]
À l'aide des connaissances du cours, réalise un croquis de l'Angola faisant apparaître :
\begin{enumerate}
\item les zones d'influence des trois mouvements de libération : MPLA (région de Luanda et de l'est), FNLA (nord), UNITA (sud et centre) ;
\item les principales villes : Luanda, Huambo, Cabinda ;
\item les soutiens extérieurs de chaque mouvement, indiqués par des flèches légendées (URSS et Cuba pour le MPLA ; États-Unis et Zaïre pour le FNLA ; Afrique du Sud pour l'UNITA).
\end{enumerate}
Ton croquis comportera obligatoirement un titre, une orientation (flèche du nord) et une légende organisée.
\end{exercice}

\begin{exercice}[Analyse de tableau statistique]
Le tableau suivant donne les résultats du référendum du 28 septembre 1958 en Guinée.

\begin{center}
\begin{tabularx}{\linewidth}{|l|X|X|}
\hline
\rowcolor{popblueL} & Voix & Part en \% \\
\hline
<< Oui >> à la Communauté & \num{56981} & 4,8 \% \\
\hline
<< Non >> à la Communauté & \num{1136324} & 95,2 \% \\
\hline
\end{tabularx}
\end{center}

\begin{enumerate}
\item Calcule le nombre total de suffrages exprimés.
\item Vérifie par le calcul les pourcentages annoncés dans le tableau.
\item Que signifie politiquement ce résultat ? Rédige une conclusion de cinq à huit lignes montrant en quoi ce vote constitue une voie originale d'accession à l'indépendance.
\end{enumerate}
\end{exercice}

\courssub{Je me prépare au Bac}

\begin{exercice}[Dissertation type Baccalauréat]
\textbf{Sujet :} << Les voies de l'accession de l'Afrique à l'indépendance >>

En t'appuyant sur les cas étudiés dans l'unité (le Ghana, l'Algérie, la Guinée Conakry et l'Angola), tu rédigeras une dissertation organisée comprenant :
\begin{enumerate}
\item une introduction avec une accroche, une définition des termes du sujet, une problématique et l'annonce du plan ;
\item un développement en deux ou trois parties montrant la diversité des voies d'émancipation : voie négociée (Ghana), voie référendaire (Guinée), voie de la lutte armée (Algérie, Angola) ;
\item une conclusion qui répond à la problématique et ouvre sur les limites des indépendances (néocolonialisme).
\end{enumerate}
\emph{Durée conseillée : 2 heures. Barème indicatif : introduction 4 pts, développement 12 pts, conclusion 4 pts.}
\end{exercice}

\begin{exercice}[Étude critique de documents : le << non >> de la Guinée]
\textbf{Document 1 :} Extrait du discours de Sékou Touré au général de Gaulle, Conakry, 25 août 1958.

<< Nous préférons la liberté dans la pauvreté à la richesse dans l'esclavage [\dots] Nous exercerons notre droit à l'indépendance, mais nous resterons proches de la France. >>

\textbf{Document 2 :} Résultats du référendum du 28 septembre 1958 en Guinée : le << non >> l'emporte avec \num{1136324} voix (95,2 \%) contre \num{56981} voix (4,8 \%) pour le << oui >>. La Guinée proclame son indépendance le 2 octobre 1958. En représailles, la France retire immédiatement ses personnels administratifs et techniques et suspend toute aide.

\textbf{Travail à faire :}
\begin{enumerate}
\item Présente chaque document (nature, auteur, date, contexte). (4 pts)
\item Dégage le thème commun des deux documents. (2 pts)
\item Analyse l'attitude de Sékou Touré face à De Gaulle : que révèle-t-elle du rapport des forces en 1958 ? (6 pts)
\item Montre, à partir du document 2 et de tes connaissances, que le choix guinéen a eu un coût immédiat mais une portée historique majeure pour l'Afrique. (6 pts)
\item Conclus en situant le cas guinéen parmi les voies de la décolonisation africaine. (2 pts)
\end{enumerate}
\end{exercice}

\begin{exercice}[Débat argumenté type Probatoire : le franc CFA]
\textbf{Sujet :} << Le franc CFA : instrument de coopération ou chaîne du néocolonialisme ? >>

Après les indépendances, la plupart des anciennes colonies françaises d'Afrique subsaharienne sont restées dans la zone franc : leur monnaie, le franc CFA, est garantie par le Trésor français et arrimée à la monnaie française (puis à l'euro), avec une partie des réserves de change déposée en France.

\textbf{Travail à faire :}
\begin{enumerate}
\item Rappelle la définition du néocolonialisme et deux de ses manifestations en Afrique (accords de coopération, présence de firmes étrangères, interventions militaires). (4 pts)
\item Présente deux arguments en faveur de la thèse << instrument de coopération >> (stabilité monétaire, garantie de change). (4 pts)
\item Présente deux arguments en faveur de la thèse << chaîne du néocolonialisme >> (dépendance monétaire, contrôle français des réserves, frein à la souveraineté économique). (4 pts)
\item Rédige en une page maximum une prise de position personnelle et justifiée, en mobilisant au moins deux exemples africains précis. (8 pts)
\end{enumerate}
\end{exercice}

\newpage

\coursec{Corrigés des exercices}

\begin{corrige}[Exercice 1 — QCM : dates, acteurs et notions]
\textbf{1. Réponse b).} Le 6 mars 1957, la Gold Coast britannique devient le \cle{Ghana} indépendant, première colonie d'Afrique noire à accéder à l'indépendance. L'indépendance de l'Algérie date du 5 juillet 1962, celle de l'Angola du 11 novembre 1975, et le référendum guinéen du 28 septembre 1958.

\textbf{2. Réponse a).} Le 1\ier{} novembre 1954, une série d'attentats du \cle{FLN} en Algérie (la << Toussaint rouge >>) marque le début de la guerre d'Algérie. Les accords d'Évian datent du 18 mars 1962 ; l'indépendance de la Guinée du 2 octobre 1958 ; la CPP est fondée en 1949.

\textbf{3. Réponse b).} Le MPLA (Mouvement populaire de libération de l'Angola) est un mouvement nationaliste angolais dirigé par \cle{Agostinho Neto}, soutenu par l'URSS et Cuba.

\textbf{4. Réponse b).} Kwame Nkrumah fonde en 1949 la \cle{Convention People's Party} (CPP), parti de masse qui conduit la Gold Coast à l'indépendance. Le FLN est algérien, le MPLA angolais, et la Françafrique désigne les réseaux franco-africains de l'après-indépendance.

\textbf{5. Réponse b).} Le 28 septembre 1958, la Guinée vote << non >> à 95,2 \% au référendum de De Gaulle et proclame son indépendance le 2 octobre 1958, sous la conduite de Sékou Touré.

\textbf{Points de vigilance :} ne pas confondre le 6 mars 1957 (Ghana) avec le 5 juillet 1962 (Algérie) ; ne pas attribuer le MPLA à Nkrumah ; bien distinguer la date du référendum (28 septembre) de celle de la proclamation d'indépendance (2 octobre 1958).
\end{corrige}

\begin{corrige}[Exercice 2 — Vrai ou faux ?]
\textbf{1. Vrai.} Le Ghana (ex-Gold Coast britannique) accède à l'indépendance le 6 mars 1957, avant toutes les autres colonies d'Afrique noire.

\textbf{2. Vrai.} Les accords d'Évian, signés le 18 mars 1962 entre la France et le GPRA (FLN), instaurent un cessez-le-feu et organisent l'autodétermination, qui aboutit à l'indépendance du 5 juillet 1962.

\textbf{3. Faux.} L'indépendance de l'Angola est bien proclamée le 11 novembre 1975, mais elle est le fruit d'une longue lutte armée contre le Portugal (à partir de 1961) et s'accompagne d'une guerre civile entre MPLA, FNLA et UNITA.

\textbf{4. Faux.} Sékou Touré a refusé la Communauté : la Guinée a voté << non >> au référendum du 28 septembre 1958, seul territoire d'Afrique française à le faire.

\textbf{5. Vrai.} Le néocolonialisme désigne le maintien d'une domination économique, monétaire et politique des anciennes puissances coloniales (et des grandes puissances) sur les États africains devenus formellement indépendants (exemple : zone franc, accords de coopération).

\textbf{Points de vigilance :} une justification sans date ni exemple précis est insuffisante ; toujours citer un fait du cours.
\end{corrige}

\begin{corrige}[Exercice 3 — Définitions]
\textbf{1. Décolonisation.} Processus par lequel les territoires colonisés accèdent à la souveraineté nationale après la Seconde Guerre mondiale. Elle peut être négociée, référendaire ou armée. \emph{Exemple :} l'indépendance négociée du Ghana le 6 mars 1957.

\textbf{2. Néocolonialisme.} Ensemble des formes de domination économique, monétaire et politique exercées par les anciennes métropoles (ou de grandes puissances) sur des États officiellement indépendants. \emph{Exemple :} le maintien du franc CFA garanti par le Trésor français.

\textbf{3. Françafrique.} Réseau de relations privilégiées (accords de coopération, aides, interventions militaires, réseaux politiques et économiques) tissé entre la France et ses anciennes colonies africaines après 1960. \emph{Exemple :} les interventions militaires françaises au Gabon ou au Tchad.

\textbf{4. Mouvement de libération nationale.} Organisation politique et/ou armée qui lutte contre la domination coloniale pour obtenir l'indépendance. \emph{Exemples :} le FLN en Algérie (1954), le MPLA, le FNLA et l'UNITA en Angola (à partir de 1961).

\textbf{Points de vigilance :} une définition sans exemple tiré du cours ne vaut pas la totalité des points ; éviter les définitions circulaires (<< la décolonisation, c'est décoloniser >>).
\end{corrige}

\begin{corrige}[Exercice 4 — Chronologie]
Ordre chronologique correct :
\begin{enumerate}
\item \textbf{1\ier{} novembre 1954 :} Toussaint rouge en Algérie (début de la guerre d'Algérie) ;
\item \textbf{6 mars 1957 :} indépendance du Ghana ;
\item \textbf{28 septembre 1958 :} référendum en Guinée (victoire du << non >>, indépendance le 2 octobre 1958) ;
\item \textbf{5 juillet 1962 :} indépendance de l'Algérie (après les accords d'Évian du 18 mars 1962) ;
\item \textbf{11 novembre 1975 :} indépendance de l'Angola.
\end{enumerate}

\begin{popfigure}
\begin{center}
\begin{tikzpicture}[scale=1]
\draw[-{Stealth},thick,popdark] (0,0) -- (13.5,0) node[below left=2pt and 0pt]{};
\foreach \x/\d/\e in {1/1954/{Toussaint rouge\\ (Algérie)}, 4/1957/{Indépendance\\ du Ghana}, 7/1958/{Référendum\\ guinéen}, 10/1962/{Indépendance\\ de l'Algérie}, 13/1975/{Indépendance\\ de l'Angola}}{
  \draw[thick,popblue] (\x,0) -- (\x,0.25);
  \node[below,font=\small\bfseries,text=popdark] at (\x,0) {\d};
  \node[above,font=\footnotesize,align=center,text=popink] at (\x,0.25) {\e};
}
\end{tikzpicture}
\end{center}
\legende{Frise chronologique : les grandes étapes de la décolonisation africaine étudiées dans l'unité (1954--1975).}
\end{popfigure}

\textbf{Points de vigilance :} la frise doit comporter une flèche d'orientation du temps, des dates exactes (jour, mois, année) et des libellés courts ; l'échelle n'a pas besoin d'être proportionnelle mais l'ordre doit être strictement respecté.
\end{corrige}

\begin{corrige}[Exercice 5 — Question à développement construit : le Ghana]
\textbf{Corrigé rédigé (modèle).}

Le Ghana (ex-Gold Coast) illustre la voie négociée de la décolonisation : l'indépendance du 6 mars 1957 est obtenue sans guerre, par la combinaison d'une mobilisation populaire organisée et d'étapes institutionnelles acceptées par la Grande-Bretagne.

D'abord, la lutte est conduite par un leader charismatique, \cle{Kwame Nkrumah}, qui fonde en 1949 la \cle{Convention People's Party} (CPP), un parti de masse ouvert aux travailleurs, aux paysans et à la jeunesse scolarisée. Contrairement aux notables modérés, Nkrumah exige << l'autogouvernement maintenant >>.

Ensuite, la CPP mène des actions non violentes mais fermes : boycotts des produits britanniques, grèves et campagne de désobéissance civile regroupées sous le nom de << \cle{Positive Action} >> (1950). Nkrumah est emprisonné, ce qui accroît sa popularité.

Enfin, l'émancipation suit des étapes institutionnelles : la constitution de 1951 instaure un parlement ; la CPP remporte les élections de 1951 et Nkrumah, libéré, devient chef du gouvernement (Leader of Government Business), puis Premier ministre en 1952 ; de nouvelles victoires électorales (1954, 1956) confirment sa légitimité. La Grande-Bretagne accepte alors le transfert de souveraineté : le 6 mars 1957, la Gold Coast devient le Ghana indépendant, au sein du Commonwealth.

Ainsi, la décolonisation ghanéenne est négociée : la pression populaire non violente et les succès électoraux ont contraint la métropole à céder le pouvoir par voie constitutionnelle, sans conflit armé.

\textbf{Barème indicatif :} rôle de Nkrumah et de la CPP (5 pts) ; actions menées (boycott, désobéissance civile, Positive Action) (5 pts) ; étapes institutionnelles 1951--1957 (6 pts) ; cohérence et conclusion (4 pts).

\textbf{Points de vigilance :} ne pas écrire que le Ghana s'est libéré par les armes ; citer les dates des élections et de l'indépendance ; le terme << négociée >> doit être explicité (pas de guerre, transfert constitutionnel du pouvoir).
\end{corrige}

\begin{corrige}[Exercice 6 — Question à développement construit : la guerre d'Algérie]
\textbf{Corrigé rédigé (modèle).}

Contrairement au Ghana, l'Algérie accède à l'indépendance au prix d'une guerre longue et meurtrière (1954--1962). Trois facteurs expliquent cette radicalisation.

D'abord, l'Algérie est une colonie de peuplement : environ un million d'Européens y sont installés, et le pays est juridiquement rattaché à la France (départements). Les colons et une partie de l'armée refusent toute émancipation. De plus, les révoltes antérieures ont été durement réprimées : la répression de Sétif et Guelma (mai 1945) a fait des milliers de morts musulmans, détruisant l'espoir d'une évolution pacifique.

Ensuite, le \cle{FLN} (Front de libération nationale), créé en 1954, choisit la lutte armée : le 1\ier{} novembre 1954, la << Toussaint rouge >> lance l'insurrection. Le FLN organise une guérilla, mobilise les populations et porte la cause algérienne à l'ONU. La France répond par une guerre totale (plus de 400 000 soldats, torture, camps de regroupement), ce qui radicalise encore le conflit et provoque la chute de la IV\ieme{} République (1958).

Enfin, l'issue est politique : le général de Gaulle, revenu au pouvoir, finit par admettre l'autodétermination. Les \cle{accords d'Évian} (18 mars 1962) instaurent le cessez-le-feu ; le référendum d'autodétermination (1\ier{} juillet 1962) approuve massivement l'indépendance, proclamée le \cle{5 juillet 1962}. L'exode des pieds-noirs et les massacres de harkis accompagnent cette fin tragique.

Ainsi, la nature de la colonisation (peuplement, intégration juridique) et l'intransigeance des deux camps ont rendu la négociation impossible jusqu'en 1962 : l'Algérie illustre la voie armée de la décolonisation.

\textbf{Barème indicatif :} causes de la radicalisation (5 pts) ; rôle du FLN et déroulement (5 pts) ; issue (Évian, 5 juillet 1962) (6 pts) ; cohérence (4 pts).

\textbf{Points de vigilance :} distinguer le cessez-le-feu (18 mars 1962) de l'indépendance (5 juillet 1962) ; citer Sétif (1945) comme cause de radicalisation ; ne pas réduire la guerre à un simple conflit militaire (dimension diplomatique et sociale).
\end{corrige}

\begin{corrige}[Exercice 7 — Construction cartographique : l'Angola]
\textbf{Éléments attendus sur le croquis :} contour simplifié de l'Angola ; trois zones colorées (MPLA : Luanda et l'est ; FNLA : nord ; UNITA : centre-sud) ; villes Luanda (côte), Huambo (centre), Cabinda (enclave au nord) ; flèches de soutien extérieur légendées ; titre, flèche du nord, légende.

\begin{popfigure}
\begin{center}
\begin{tikzpicture}[scale=0.9,font=\small]
% Contour simplifié de l'Angola
\draw[thick,popdark,fill=popcream] (0,0) -- (6,0) -- (6.5,2) -- (6,5.5) -- (4.5,6.5) -- (1.5,6.2) -- (0.3,5.8) -- (0,4) -- cycle;
% Enclave de Cabinda
\draw[thick,popdark,fill=popgreenL] (0.4,7.2) -- (1.6,7.2) -- (1.6,8) -- (0.4,8) -- cycle;
\node[font=\footnotesize] at (1,7.6) {Cabinda};

% Zones d'influence (corrigées géographiquement)
% FNLA (Orange) : Nord (au-dessus de Luanda, ~ y > 5.0)
\fill[poporange!35] (0.16,5.0) -- (0.3,5.8) -- (1.5,6.2) -- (4.5,6.5) -- (6,5.5) -- (6,5.0) -- cycle;
% MPLA (Bleu) : Ouest (Luanda) + Est
\fill[popblue!35] (0,0) -- (0,4) -- (0.16,5.0) -- (1.5,4.5) -- (1.5,0) -- cycle; % Bande côtière ouest (Luanda)
\fill[popblue!35] (5,0) -- (6,0) -- (6.5,2) -- (6,5.5) -- (6,5.0) -- (5,4.5) -- cycle; % Bande est
% UNITA (Violet) : Centre-Sud (le reste)
\fill[poppurple!30] (1.5,0) -- (5,0) -- (5,4.5) -- (6,5.0) -- (0.16,5.0) -- (1.5,4.5) -- cycle;

% Villes
\fill[popdark] (0.6,4.6) circle (2pt) node[left,font=\footnotesize\bfseries]{Luanda};
\fill[popdark] (2.6,2.6) circle (2pt) node[below,font=\footnotesize]{Huambo};

% Nord
\draw[-{Stealth},thick,popdark] (7.3,0.5) -- (7.3,1.7);
\node[font=\bfseries] at (7.3,2) {N};

% Flèches de soutien extérieur
\draw[-{Stealth},very thick,popblue] (8.6,6.5) -- (5.2,5.2) node[midway,above,sloped,font=\footnotesize]{URSS + Cuba};
\draw[-{Stealth},very thick,poporange] (8.6,4.5) -- (1.6,5.5) node[midway,above,sloped,font=\footnotesize]{États-Unis + Zaïre};
\draw[-{Stealth},very thick,poppurple] (8.6,1.5) -- (4.5,2.0) node[midway,below,sloped,font=\footnotesize]{Afrique du Sud};

% Titre
\node[font=\bfseries,text=popink] at (3.5,8.6) {L'Angola : mouvements de libération et soutiens extérieurs (1975)};
\end{tikzpicture}
\end{center}
\legende{Croquis attendu (correction) : 1. zone MPLA (bleu, Luanda côte ouest et est) ; 2. zone FNLA (orange, nord) ; 3. zone UNITA (violet, centre-sud) ; 4. villes : Luanda, Huambo, Cabinda ; 5. flèches des soutiens extérieurs. Échelle indicative : 1 cm $\approx$ 200 km.}
\end{popfigure}

\textbf{Barème indicatif :} contour et orientation (2 pts) ; trois zones correctement placées (6 pts) ; villes (3 pts) ; flèches de soutien légendées (6 pts) ; titre et légende organisée (3 pts).

\textbf{Points de vigilance :} ne pas inverser les zones (FNLA au nord, pas au sud) ; Cabinda est une enclave séparée du reste du territoire ; chaque flèche doit être associée à une légende (un croquis sans légende est pénalisé).
\end{corrige}

\begin{corrige}[Exercice 8 — Analyse de tableau statistique]
\textbf{1. Total des suffrages exprimés :}
\[ \num{56981} + \num{1136324} = \num{1193305} \text{ suffrages exprimés.} \]

\textbf{2. Vérification des pourcentages :}
\begin{align*}
\text{Part du << oui >>} &= \frac{\num{56981}}{\num{1193305}} \times 100 \approx 4,78 \% \approx 4,8 \% \\
\text{Part du << non >>} &= \frac{\num{1136324}}{\num{1193305}} \times 100 \approx 95,22 \% \approx 95,2 \%
\end{align*}
Les pourcentages du tableau sont donc exacts (arrondis au dixième). On vérifie aussi : $4,8 + 95,2 = 100 \%$.

\textbf{3. Signification politique (conclusion rédigée).} Ce résultat signifie que le peuple guinéen a rejeté massivement la Communauté proposée par le général de Gaulle : plus de 19 électeurs sur 20 ont choisi le << non >>. C'est la seule colonie d'Afrique française à refuser l'offre française en 1958. Ce vote constitue une voie originale d'accession à l'indépendance : ni guerre (comme en Algérie), ni longue négociation (comme au Ghana), mais un référendum organisé par la métropole elle-même, dont le résultat est immédiatement appliqué — l'indépendance est proclamée dès le 2 octobre 1958. Le cas guinéen montre qu'un peuple colonisé pouvait conquérir sa souveraineté par les urnes, au prix d'une rupture brutale avec la France (retrait des personnels, suspension de l'aide), mais avec une portée symbolique immense pour toute l'Afrique.

\textbf{Points de vigilance :} présenter le calcul avec la formule (part = effectif / total $\times$ 100) ; ne pas confondre suffrages exprimés et inscrits ; la conclusion doit mobiliser la comparaison avec les autres voies de décolonisation.
\end{corrige}

\begin{corrige}[Exercice 9 — Dissertation : les voies de l'accession de l'Afrique à l'indépendance]
\textbf{Introduction (modèle).} Après 1945, l'Afrique colonisée entre dans l'ère des indépendances : entre 1957 (Ghana) et 1975 (Angola), la quasi-totalité du continent accède à la souveraineté. La décolonisation désigne le processus par lequel les territoires coloniaux deviennent des États indépendants. Mais toutes les colonies ne suivent pas le même chemin. Problématique : par quelles voies les peuples africains ont-ils accédé à l'indépendance, et qu'est-ce qui explique leur diversité ? Nous verrons que l'émancipation africaine emprunte des voies pacifiques — négociée ou référendaire (I) — mais aussi, lorsque la métropole refuse de céder, la voie de la lutte armée (II).

\textbf{I. Les voies pacifiques : négociation et référendum.}
\begin{itemize}
\item \emph{Le Ghana, voie négociée :} Nkrumah et la CPP (1949) mobilisent les masses (boycott, désobéissance civile, << Positive Action >> de 1950) ; les étapes institutionnelles (élections de 1951, 1954, 1956 ; Nkrumah Premier ministre en 1952) conduisent à l'indépendance du 6 mars 1957, sans guerre, au sein du Commonwealth.
\item \emph{La Guinée, voie référendaire :} le 28 septembre 1958, le << non >> à la Communauté de De Gaulle l'emporte avec 95,2 \% des voix ; l'indépendance est proclamée le 2 octobre 1958. C'est la seule colonie française à choisir la rupture immédiate, au prix de représailles françaises (retrait des personnels, suspension de l'aide).
\end{itemize}

\textbf{II. La voie de la lutte armée.}
\begin{itemize}
\item \emph{L'Algérie :} colonie de peuplement intégrée à la France, traumatisée par la répression de Sétif (1945) ; le FLN déclenche l'insurrection le 1\ier{} novembre 1954 (Toussaint rouge) ; après huit ans de guerre, les accords d'Évian (18 mars 1962) ouvrent la voie à l'indépendance du 5 juillet 1962.
\item \emph{L'Angola :} le Portugal de Salazar refuse toute émancipation ; MPLA (Neto), FNLA et UNITA prennent les armes dès 1961 ; la révolution des Œillets à Lisbonne (1974) permet l'indépendance le 11 novembre 1975, mais la rivalité des mouvements et les ingérences étrangères (URSS/Cuba, États-Unis/Zaïre, Afrique du Sud) plongent le pays dans la guerre civile.
\end{itemize}

\textbf{Conclusion (modèle).} L'Afrique a donc accédé à l'indépendance par des voies diverses : négociation constitutionnelle au Ghana, référendum en Guinée, guerres de libération en Algérie et en Angola. La voie suivie dépend de la nature de la colonisation (peuplement ou exploitation) et de l'attitude de la métropole. Mais ces indépendances restent inachevées : le néocolonialisme (zone franc, accords de coopération, firmes étrangères) perpétue la dépendance, posant la question de la véritable souveraineté africaine.

\textbf{Barème indicatif :} introduction (accroche, définitions, problématique, plan) 4 pts ; développement (deux parties équilibrées, faits et dates précis) 12 pts ; conclusion avec ouverture 4 pts.

\textbf{Points de vigilance :} annoncer le plan dans l'introduction ; chaque partie doit contenir au moins deux exemples datés ; éviter le catalogue de faits sans fil directeur (la problématique doit commander l'argumentation) ; l'ouverture sur le néocolonialisme est exigée par le sujet.
\end{corrige}

\begin{corrige}[Exercice 10 — Étude critique de documents : le << non >> de la Guinée]
\textbf{1. Présentation des documents (4 pts).}
\begin{itemize}
\item \emph{Document 1 :} extrait d'un discours politique prononcé par \cle{Sékou Touré}, leader du PDG et chef du gouvernement guinéen, à Conakry le 25 août 1958, devant le général de Gaulle en tournée africaine pour promouvoir la Communauté franco-africaine, un mois avant le référendum du 28 septembre 1958.
\item \emph{Document 2 :} tableau de résultats électoraux (données statistiques officielles) du référendum du 28 septembre 1958 en Guinée, suivi d'informations sur la proclamation d'indépendance (2 octobre 1958) et les représailles françaises.
\end{itemize}

\textbf{2. Thème commun (2 pts).} Les deux documents portent sur le choix de l'indépendance immédiate de la Guinée en 1958 : le refus de la Communauté française et ses conséquences.

\textbf{3. L'attitude de Sékou Touré (6 pts).} Sékou Touré affirme face à De Gaulle une position de dignité et de fermeté : la formule << nous préférons la liberté dans la pauvreté à la richesse dans l'esclavage >> hiérarchise les valeurs — la souveraineté prime sur l'aide matérielle. Mais il reste courtois et ouvert (<< nous resterons proches de la France >>) : il ne rompt pas avec la métropole, il revendique une relation d'égal à égal. Cette attitude révèle un rapport de forces nouveau en 1958 : les leaders africains, portés par l'opinion et le contexte international (ONU, mouvement de Bandung), peuvent désormais tenir tête à la métropole ; la France de De Gaulle doit accepter le principe de l'autodétermination, même si elle punit ceux qui l'exercent contre son gré.

\textbf{4. Coût immédiat et portée historique (6 pts).} Le document 2 montre le coût immédiat du choix guinéen : la France retire ses personnels administratifs et techniques et suspend toute aide, laissant le jeune État sans cadres ni ressources — c'est la sanction du << non >> plébiscité (95,2 \%). Mais la portée historique est majeure : la Guinée devient le premier État d'Afrique occidentale française indépendant, prouvant qu'une colonie pouvait dire << non >> à De Gaulle et survivre ; son exemple accélère les indépendances de 1960 (<< année de l'Afrique >>) et fait de Sékou Touré une figure du panafricanisme.

\textbf{5. Conclusion (2 pts).} Le cas guinéen illustre une voie originale de décolonisation : la voie référendaire, pacifique et légale, mais de rupture radicale — à mi-chemin entre la négociation ghanéenne et la lutte armée algérienne.

\textbf{Points de vigilance :} la présentation d'un document exige nature, auteur, date et contexte (4 éléments) ; la question 3 demande une analyse du rapport de forces, pas une paraphrase ; la question 4 doit articuler le document 2 ET les connaissances personnelles (année 1960).
\end{corrige}

\begin{corrige}[Exercice 11 — Débat argumenté : le franc CFA]
\textbf{1. Définition et manifestations du néocolonialisme (4 pts).} Le \cle{néocolonialisme} désigne la domination économique, monétaire et politique qu'exercent les anciennes puissances coloniales sur des États formellement indépendants. Manifestations : les accords de coopération inégaux signés au moment des indépendances (assistance technique, clauses commerciales préférentielles) ; la présence de firmes étrangères contrôlant les matières premières (pétrole, cacao, minerais) ; les interventions militaires de la France dans ses anciennes colonies (Gabon, Tchad) ; le maintien de la zone franc.

\textbf{2. Thèse << instrument de coopération >> (4 pts).}
\begin{itemize}
\item \emph{Stabilité monétaire :} la garantie du Trésor français et la parité fixe avec le franc puis l'euro protègent les pays de la zone franc des hyperinflations qui ont frappé d'autres États africains (ex. : le Zaïre de Mobutu avec le zaïre).
\item \emph{Garantie de change et crédibilité :} la convertibilité garantie rassure les investisseurs et facilite les échanges extérieurs des pays membres (ex. : la Côte d'Ivoire a longtemps attiré les capitaux grâce à cette stabilité).
\end{itemize}

\textbf{3. Thèse << chaîne du néocolonialisme >> (4 pts).}
\begin{itemize}
\item \emph{Dépendance monétaire :} les pays membres ne contrôlent pas leur politique monétaire (parité fixe décidée avec la France, dévaluation de 1994 négociée sous forte contrainte extérieure), ce qui limite leur souveraineté économique.
\item \emph{Contrôle des réserves :} une partie des réserves de change des États africains est déposée au Trésor français ; la France siège dans les instances monétaires, symbole d'une tutelle persistante dénoncée par de nombreux économistes africains.
\end{itemize}

\textbf{4. Prise de position personnelle (8 pts) — modèle.} Le franc CFA apparaît davantage comme une chaîne du néocolonialisme que comme un simple instrument de coopération, même s'il a rendu de réels services. Certes, la stabilité monétaire est un fait : les pays de la zone franc ont connu des inflations plus faibles que le Ghana ou la Guinée, qui ont dû créer puis réformer leurs monnaies nationales (le cedi ghanéen, le franc guinéen puis le syli de Sékou Touré, marques de souveraineté mais sources d'instabilité). Cependant, cette stabilité a un prix : la dévaluation du franc CFA de 50 \% en janvier 1994, négociée sous pression française et du FMI, a brutalement appauvri les populations du Sénégal, de la Côte d'Ivoire ou du Cameroun, sans que leurs gouvernements aient réellement le choix. De plus, le dépôt d'une partie des réserves de change en France prive ces États de ressources pour financer leur développement. La comparaison avec le Ghana, qui gère seul sa monnaie depuis 1957 malgré les difficultés, montre que la souveraineté monétaire est un élément constitutif de l'indépendance. Le franc CFA, hérité direct de l'époque coloniale (il signifiait à l'origine << franc des colonies françaises d'Afrique >>), perpétue donc une dépendance incompatible avec une souveraineté pleine et entière : c'est bien un instrument du néocolonialisme, même si sa suppression exige une préparation rigoureuse pour éviter le chaos monétaire.

\textbf{Barème indicatif :} définition + deux manifestations (4 pts) ; deux arguments << coopération >> (4 pts) ; deux arguments << néocolonialisme >> (4 pts) ; prise de position structurée avec au moins deux exemples précis (8 pts).

\textbf{Points de vigilance :} un débat argumenté exige de présenter honnêtement les DEUX thèses avant de trancher ; la prise de position doit mobiliser des exemples africains datés (dévaluation de 1994, monnaies du Ghana et de la Guinée) ; éviter les affirmations sans preuve et rester dans le cadre du cours (néocolonialisme, zone franc).
\end{corrige}

\newpage

\coursec{Fiche bilan}

\begin{popfigure}
\begin{tikzpicture}[mindmap, grow cyclic, every node/.style={concept, font=\small},
root concept/.append style={concept color=popblue, font=\small\bfseries, minimum size=3.2cm},
level 1/.append style={level distance=3.6cm, sibling angle=60, font=\scriptsize},
level 2/.append style={level distance=2.4cm, font=\tiny}]
\node[root concept]{Décolonisation et néocolonialisme en Afrique}
  child[concept color=poporange]{ node {Facteurs de la décolonisation}
    child { node {2\textsuperscript{e} Guerre mondiale} }
    child { node {ONU, USA--URSS} }
    child { node {Conscience nationaliste} }
  }
  child[concept color=popgreen]{ node {Ghana (britannique)}
    child { node {Nkrumah, CPP} }
    child { node {Indépendance 1957} }
  }
  child[concept color=poppink]{ node {Algérie (française)}
    child { node {FLN, 1954--1962} }
    child { node {Accords d'Évian 1962} }
  }
  child[concept color=popgold]{ node {Guinée Conakry}
    child { node {<< Non >> de 1958} }
    child { node {Sékou Touré} }
  }
  child[concept color=poppurple]{ node {Angola (portugaise)}
    child { node {MPLA, guerre armée} }
    child { node {Indépendance 1975} }
  }
  child[concept color=popmuted]{ node {Néocolonialisme}
    child { node {Dépendance économique} }
    child { node {Monnaies, accords, ingérences} }
  };
\end{tikzpicture}
\legende{Carte mentale de la leçon : au centre, la décolonisation africaine ; autour, les facteurs généraux, les quatre cas étudiés (Ghana, Algérie, Guinée, Angola) et la question du néocolonialisme.}
\end{popfigure}

\begin{bilan}
\textbf{Définitions clés}
\begin{itemize}
\item \cle{Décolonisation} : processus par lequel les territoires colonisés accèdent à la souveraineté et deviennent des États indépendants (en Afrique, surtout entre 1957 et 1975).
\item \cle{Nationalisme} : volonté d'un peuple de s'affirmer comme nation et de conquérir son indépendance politique.
\item \cle{Néocolonialisme} : domination économique, financière et politique exercée par les anciennes puissances coloniales (et les grandes puissances) sur des États pourtant officiellement indépendants.
\item \cle{Guerre de libération} : lutte armée menée par un mouvement nationaliste contre la puissance coloniale (Algérie, Angola).
\end{itemize}

\textbf{Repères chronologiques à connaître par c\oe ur}
\begin{center}
\begin{tabularx}{\linewidth}{|l|X|}\hline
\rowcolor{popblueL} \textbf{Date} & \textbf{Événement} \\ \hline
1945 & Conférence de Manchester (5\textsuperscript{e} congrès panafricain) ; création de l'ONU \\ \hline
1947 & Création du CPP de Kwame Nkrumah au Gold Coast \\ \hline
1\textsuperscript{er} nov. 1954 & Déclenchement de l'insurrection algérienne (Toussaint rouge) par le FLN \\ \hline
6 mars 1957 & Indépendance du Ghana, première colonie d'Afrique noire indépendante \\ \hline
28 sept. 1958 & Référendum : la Guinée vote << non >> à la Communauté de de Gaulle \\ \hline
2 oct. 1958 & Proclamation de l'indépendance de la Guinée (Sékou Touré) \\ \hline
18 mars 1962 & Accords d'Évian ; cessez-le-feu en Algérie \\ \hline
5 juillet 1962 & Indépendance de l'Algérie \\ \hline
1961--1974 & Guerre de libération en Angola (MPLA, FNLA, UNITA) \\ \hline
11 nov. 1975 & Indépendance de l'Angola (Agostinho Neto) \\ \hline
\end{tabularx}
\end{center}

\textbf{Les quatre cas comparés}
\begin{center}
\begin{tabularx}{\linewidth}{|l|X|X|X|}\hline
\rowcolor{popblueL} \textbf{Colonie} & \textbf{Puissance} & \textbf{Voie} & \textbf{Acteur clé} \\ \hline
Ghana & Grande-Bretagne & Négociée, lutte politique (CPP) & K. Nkrumah \\ \hline
Algérie & France & Guerre (1954--1962) & FLN \\ \hline
Guinée & France & Rupture référendaire (1958) & Sékou Touré \\ \hline
Angola & Portugal & Guerre armée (1961--1975) & MPLA / Neto \\ \hline
\end{tabularx}
\end{center}

\textbf{Méthodes express}
\begin{itemize}
\item \emph{Commentaire de document} : présenter (nature, auteur, date, contexte) $\rightarrow$ dégager les idées $\rightarrow$ expliquer avec les connaissances $\rightarrow$ conclure sur la portée.
\item \emph{Frise chronologique} : toujours dater et situer les acteurs avant d'expliquer causes et conséquences.
\item \emph{Sujet de synthèse} : distinguer décolonisation négociée (Ghana, Guinée) et décolonisation par les armes (Algérie, Angola).
\end{itemize}
\end{bilan}

\begin{aretenir}[Checklist avant le Bac]
\begin{itemize}
\item[$\square$] Je sais définir décolonisation, nationalisme et néocolonialisme.
\item[$\square$] Je cite les trois grands facteurs de la décolonisation africaine.
\item[$\square$] Je connais la date de l'indépendance du Ghana (6 mars 1957) et le rôle de Nkrumah.
\item[$\square$] Je sais que la guerre d'Algérie dure de 1954 à 1962 et s'achève par les accords d'Évian.
\item[$\square$] J'explique le << non >> guinéen du 28 septembre 1958 et ses conséquences.
\item[$\square$] Je situe la guerre de libération angolaise (1961--1975) et le rôle du MPLA.
\item[$\square$] Je distingue indépendance négociée et indépendance par la lutte armée.
\item[$\square$] Je cite au moins deux manifestations du néocolonialisme (monnaie, accords de coopération, sociétés étrangères, ingérences).
\item[$\square$] Je sais construire une frise chronologique de 1945 à 1975.
\item[$\square$] Je maîtrise les étapes du commentaire de document historique.
\end{itemize}
\end{aretenir}

\findecours

\end{document}
